% Leçon 40 — Grammaire : 一边……一边…….. ; la ponctuation en chinois — Chinois 1ère A — Cours signé OnBuch+
% Programme officiel MINESEC. Compiler avec : tectonic ZHA40-grammaire-la-ponctuation-en-chinois.tex
\newcommand{\DOCMATIERE}{Chinois}
\newcommand{\DOCNIVEAU}{1ère A}
\newcommand{\DOCMODULE}{Module 5 : Mass médias et télécommunication}
\newcommand{\DOCLECON}{ZHA40}
\newcommand{\DOCTITRE}{Leçon 40 — Grammaire : 一边……一边…….. ; la ponctuation en chinois}
% OnBuch+ — préambule « cours signés » ENRICHI (compilé avec tectonic / XeLaTeX)
% Système visuel « Pop » d'OnBuch+ : fond crème (jamais blanc), bordures encre
% épaisses, ombres « dures » décalées, titres Archivo Black en capitales.
% Version MATHÉMATIQUES : graphes (pgfplots/tikz), boîtes pédagogiques
% (définition, propriété, méthode, attention, le savais-tu, exercice résolu,
% exercice, TP, bilan), page de garde et signature OnBuch+ ancrée en bas.
%
% Macros attendues dans main.tex AVANT \input{preamble} :
%   \DOCMATIERE  \DOCNIVEAU  \DOCMODULE  \DOCLECON (numéro)  \DOCTITRE
\documentclass[11pt]{article}

\usepackage{fontspec}
\usepackage[french]{babel} % français = langue principale ; le chinois via \zh{..} / textebox (xeCJK)
\usepackage{xeCJK}
\usepackage{pifont} % \ding{51} \ding{55} utilisés par les modèles
\usepackage[a4paper,margin=2cm,top=3.4cm,bottom=2.8cm,headheight=1.4cm,footskip=1.3cm]{geometry}
\usepackage{amsmath,amssymb,amsfonts}
\usepackage{mathtools} % \xmapsto, \coloneqq... souvent utilisés par les modèles
\usepackage{cancel} % simplifications barrées (bilans de forces)
% \abs{z} (module, valeur absolue) et \norm{u} (norme), utilisés par les modèles
\providecommand{\abs}{}\let\abs\relax\DeclarePairedDelimiter{\abs}{\lvert}{\rvert}
\providecommand{\norm}{}\let\norm\relax\DeclarePairedDelimiter{\norm}{\lVert}{\rVert}
\usepackage[table]{xcolor} % [table] : fournit aussi \rowcolors (alternance de lignes)
\usepackage{pagecolor}
\usepackage{tikz}
\usetikzlibrary{arrows.meta,positioning,calc,shapes.geometric,shapes.misc,shapes.arrows,decorations.pathmorphing,decorations.pathreplacing,decorations.markings,patterns,shadows,mindmap,trees,fit,backgrounds,matrix,intersections,angles,quotes}
\usepackage{pgfplots}
\usepackage[version=4]{mhchem} % équations nucléaires \ce{^{235}_{92}U}
\usepackage[european,siunitx]{circuitikz} % schémas électriques
\pgfplotsset{compat=1.18}
\usepgfplotslibrary{fillbetween,groupplots} % aires entre courbes (calcul intégral)
\usepackage{enumitem}
\usepackage{fancyhdr}
% [most] charge toutes les bibliothèques tcolorbox (listings, minted, raster,
% documentation...) et gonfle inutilement la mémoire TeX sur un document long
% (~300 boîtes) : on ne charge que ce qui est réellement utilisé.
\usepackage[skins,breakable]{tcolorbox}
\usepackage{graphicx}
\usepackage{colortbl}
\usepackage{booktabs}
\usepackage{tabularx}
\usepackage{diagbox} % case d'en-tête diagonale des tableaux à double entrée
% Colonnes extensibles centrée / gauche / droite (C, L, R), souvent utilisées par les modèles
\newcolumntype{C}{>{\centering\arraybackslash}X}
\newcolumntype{L}{>{\raggedright\arraybackslash}X}
\newcolumntype{R}{>{\raggedleft\arraybackslash}X}
\usepackage{array}
\usepackage{multicol}
\usepackage{multirow}
\usepackage{siunitx}
% parse-numbers=false : accepte aussi \SI{2,5 \times 10^{-3}}{...}
\sisetup{locale=FR,output-decimal-marker={,},group-minimum-digits=4,parse-numbers=false}
\DeclareSIUnit\ohm{\ensuremath{\symup{\Omega}}} % U+2126 absent de la police
\DeclareSIUnit\division{div} % réglages d'oscilloscope (ms/div, V/div)
\DeclareSIUnit\tr{tr} % tours (vitesse de rotation, ex. \SI{1500}{\tr\per\minute})
\DeclareSIUnit\molar{M} % concentration molaire (ex. \SI{3}{\molar}, \SI{1}{\micro\molar})
\DeclareSIUnit\an{an} % année (le modèle confond parfois avec \year, un compteur TeX) — ex. \SI{5}{\kilo\meter\per\an}
\DeclareSIUnit\FCFA{FCFA} % franc CFA (ex. \SI{500}{\FCFA\per\kilo\gram})
\DeclareSIUnit\degreeBrix{\textdegree Brix} % teneur en sucre d'un sirop/jus (réfractométrie)
\DeclareSIUnit\month{mois} % le modèle utilise parfois \month, un compteur TeX, comme unité de durée
\usepackage{qrcode}
% \degree : commande fréquemment utilisée par les modèles pour un angle ou une
% température en dehors de \SI{}{} (non fournie par les paquets chargés).
\providecommand{\degree}{^{\circ}}
% \qed (fin de démonstration), utilisé par les modèles sans amsthm
\providecommand{\qed}{\hfill$\square$}
% \barre : double barre des valeurs interdites dans un tableau de variations (tkz-tab)
\providecommand{\barre}{\ensuremath{\|}}
% environnement proof (amsthm non chargé) utilisé par les modèles
\ifdefined\proof\else\newenvironment{proof}[1][Démonstration]{\par\noindent\textit{#1.}\ }{\qed\par}\fi
\usepackage{lastpage}
\usepackage{hyperref}
\hypersetup{hidelinks}

% ============================================================
% Palette « Pop » OnBuch+ — mêmes hex que la classe P (plus_ui.dart)
% ============================================================
\definecolor{poporange}{HTML}{F59321}
\definecolor{popdark}{HTML}{E07A0C}
\definecolor{popcream}{HTML}{FFF6E8}
\definecolor{popink}{HTML}{1C1714}
\definecolor{poppurple}{HTML}{7A5AE0}
\definecolor{popgreen}{HTML}{2E9E6A}
\definecolor{poppink}{HTML}{E0457B}
\definecolor{popblue}{HTML}{2F7DE1}
\definecolor{popgold}{HTML}{C98A12}
\definecolor{popmuted}{HTML}{8A7F76}
\definecolor{popline}{HTML}{EADFD2}
\definecolor{popgray}{HTML}{8A7F76}
\colorlet{popgrey}{popgray}
% Alias tolérants (le modèle nomme parfois les couleurs différemment)
\colorlet{popwhite}{white}
\colorlet{popblack}{popink}
\colorlet{popred}{poppink}
\colorlet{popyellow}{popgold}
% Teintes claires (fonds de tableaux/encadrés) — le modèle nomme parfois
% les couleurs avec un suffixe L (ex. popblueL) sans les avoir définies.
\colorlet{poporangeL}{poporange!20}
\colorlet{popdarkL}{popdark!20}
\colorlet{poppurpleL}{poppurple!20}
\colorlet{popgreenL}{popgreen!20}
\colorlet{poppinkL}{poppink!20}
\colorlet{popblueL}{popblue!20}
\colorlet{popgoldL}{popgold!20}
\colorlet{popredL}{popred!20}
\colorlet{popgrayL}{popgray!20}
% Teintes claires pour fonds de tableaux / zones
\colorlet{popblueL}{popblue!10!white}
\colorlet{poporangeL}{poporange!14!white}
\colorlet{popgreenL}{popgreen!12!white}
\colorlet{poppurpleL}{poppurple!12!white}
\colorlet{poppinkL}{poppink!12!white}
\colorlet{popgoldL}{popgold!14!white}

\pagecolor{popcream}

% ============================================================
% Typographie
% ============================================================
\defaultfontfeatures{Ligatures=TeX}
\setmainfont{PlusJakartaSans-Medium.ttf}[Path=../../../fonts/,BoldFont=PlusJakartaSans-Bold.ttf]
% Symboles, API et emojis absents de la police principale : repli sur DejaVu Sans (caractères actifs)
\newfontface\symfont{DejaVuSans.ttf}[Path=../../../fonts/]
\makeatletter
\newcommand\symactive[1]{\catcode#1=13 \begingroup\lccode`\~=#1 \lowercase{\endgroup\def~}{{\symfont\char#1}}}
\newcommand\symblank[1]{\catcode#1=13 \begingroup\lccode`\~=#1 \lowercase{\endgroup\def~}{}}
\newcommand\symhyph[1]{\catcode#1=13 \begingroup\lccode`\~=#1 \lowercase{\endgroup\def~}{\mbox{-}}}
\newcount\symc
\newcommand\symrange[2]{\symc=#1 \@whilenum\symc<#2 \do{\expandafter\symactive\expandafter{\the\symc}\advance\symc by 1 }}
\newcommand\symblankrange[2]{\symc=#1 \@whilenum\symc<#2 \do{\expandafter\symblank\expandafter{\the\symc}\advance\symc by 1 }}
\makeatother
\symrange{"0250}{"02FF}\symactive{"2713}\symactive{"2714}\symactive{"2717}\symactive{"2718}\symactive{"2610}\symactive{"2611}\symactive{"2612}
\symrange{"2600}{"27BF}
\catcode"2705=13 \begingroup\lccode`\~="2705 \lowercase{\endgroup\def~}{{\symfont\char"2713}}\symblank{"26BD}\symblank{"26A1}
\symrange{"2460}{"257F}\symactive{"223C}\symactive{"2248}\symactive{"2260}
\symactive{"01F8}\symactive{"01F9}\symactive{"1E3E}\symactive{"1E3F}
\symhyph{"2011}\symhyph{"2E1A}\symblankrange{"1F300}{"1FAFF}\symblank{"FE0F}
% Caractères chinois : WenQuanYi Zen Hei (sous-ensemble GB2312 + pinyin), gras/italique simulés
\setCJKmainfont{WenQuanYiZenHei-subset.ttf}[Path=../../../fonts/,AutoFakeBold=3,AutoFakeSlant=0.2]
\xeCJKsetup{CJKspace=false}
% Pinyin : voyelles à caron (ǎ ǐ ǒ ǔ ǖ ǘ ǚ ǜ) absentes de la police principale -> caractères actifs
\newfontface\pyfont{WenQuanYiZenHei-subset.ttf}[Path=../../../fonts/]
\newcommand\pyactive[1]{\catcode#1=13 \begingroup\lccode`\~=#1 \lowercase{\endgroup\def~}{{\pyfont\char#1}}}
\pyactive{"01CD}\pyactive{"01CE}\pyactive{"01CF}\pyactive{"01D0}\pyactive{"01D1}\pyactive{"01D2}\pyactive{"01D3}\pyactive{"01D4}\pyactive{"01D5}\pyactive{"01D6}\pyactive{"01D7}\pyactive{"01D8}\pyactive{"01D9}\pyactive{"01DA}\pyactive{"01DB}\pyactive{"01DC}
% Maths en sans-serif (Fira Math) avec chiffres et lettres droites de Plus
% Jakarta, pour que les formules se fondent dans le texte.
\usepackage[math-style=ISO,bold-style=ISO]{unicode-math}
\setmathfont{FiraMath-Regular.otf}[Path=../../../fonts/]
\setmathfont{PlusJakartaSans-Medium.ttf}[Path=../../../fonts/,range={up/num,up/{Latin,latin},\mathup}]
\usepackage{icomma} % virgule décimale sans espace en mode math
% Caractères mathématiques Unicode absents des polices de texte (Plus Jakarta,
% Archivo Black) : rendus par leur équivalent mathématique, en texte comme en
% formule — sinon ils s'affichent en carré vide (ex. « Arithmétique dans ℤ »).
\usepackage{newunicodechar}
\newunicodechar{ℝ}{\ensuremath{\mathbb{R}}}
\newunicodechar{ℤ}{\ensuremath{\mathbb{Z}}}
\newunicodechar{ℂ}{\ensuremath{\mathbb{C}}}
\newunicodechar{ℕ}{\ensuremath{\mathbb{N}}}
\newunicodechar{ℚ}{\ensuremath{\mathbb{Q}}}
\newunicodechar{𝔻}{\ensuremath{\mathbb{D}}}
\newunicodechar{α}{\ensuremath{\alpha}}
\newunicodechar{β}{\ensuremath{\beta}}
\newunicodechar{γ}{\ensuremath{\gamma}}
\newunicodechar{δ}{\ensuremath{\delta}}
\newunicodechar{ε}{\ensuremath{\varepsilon}}
\newunicodechar{θ}{\ensuremath{\theta}}
\newunicodechar{λ}{\ensuremath{\lambda}}
\newunicodechar{μ}{\ensuremath{\mu}}
\newunicodechar{σ}{\ensuremath{\sigma}}
\newunicodechar{τ}{\ensuremath{\tau}}
\newunicodechar{φ}{\ensuremath{\varphi}}
\newunicodechar{ψ}{\ensuremath{\psi}}
\newunicodechar{ω}{\ensuremath{\omega}}
\newunicodechar{Σ}{\ensuremath{\Sigma}}
% majuscules produites par \MakeUppercase dans les titres (α-aminés -> Α)
\newunicodechar{Α}{\ensuremath{\alpha}}
\newunicodechar{Β}{\ensuremath{\beta}}
\newunicodechar{Γ}{\ensuremath{\Gamma}}
\newunicodechar{Δ}{\ensuremath{\Delta}}
\newunicodechar{Ε}{\ensuremath{\varepsilon}}
\newunicodechar{Θ}{\ensuremath{\Theta}}
\newunicodechar{Λ}{\ensuremath{\Lambda}}
\newunicodechar{Μ}{\ensuremath{\mu}}
\newunicodechar{Τ}{\ensuremath{\tau}}
\newunicodechar{Φ}{\ensuremath{\Phi}}
\newunicodechar{Ψ}{\ensuremath{\Psi}}
\newunicodechar{Ω}{\ensuremath{\Omega}}
\newunicodechar{↦}{\ensuremath{\mapsto}}
\newunicodechar{∈}{\ensuremath{\in}}
\newunicodechar{∉}{\ensuremath{\notin}}
\newunicodechar{∧}{\ensuremath{\wedge}}
\newunicodechar{⇔}{\ensuremath{\Leftrightarrow}}
\newunicodechar{⟺}{\ensuremath{\Longleftrightarrow}}
\newunicodechar{⇒}{\ensuremath{\Rightarrow}}
\newunicodechar{⟹}{\ensuremath{\Longrightarrow}}
\newunicodechar{∘}{\ensuremath{\circ}}
\newunicodechar{∪}{\ensuremath{\cup}}
\newunicodechar{∩}{\ensuremath{\cap}}
\newunicodechar{≡}{\ensuremath{\equiv}}
\newunicodechar{⊂}{\ensuremath{\subset}}
\newunicodechar{⊥}{\ensuremath{\perp}}
\newunicodechar{∃}{\ensuremath{\exists}}
\newunicodechar{∀}{\ensuremath{\forall}}
\newunicodechar{∛}{\ensuremath{\sqrt[3]{\,}}}
\newunicodechar{∜}{\ensuremath{\sqrt[4]{\,}}}
\newunicodechar{⃗}{\ensuremath{{}^{\to}}}
\newunicodechar{ⁿ}{\ifmmode^{n}\else\textsuperscript{\ensuremath{n}}\fi}
\newunicodechar{ˣ}{\ifmmode^{x}\else\textsuperscript{\ensuremath{x}}\fi}
\newunicodechar{ᵇ}{\ifmmode^{b}\else\textsuperscript{\ensuremath{b}}\fi}
\newunicodechar{ʳ}{\ifmmode^{r}\else\textsuperscript{\ensuremath{r}}\fi}
\newunicodechar{ᵘ}{\ifmmode^{u}\else\textsuperscript{\ensuremath{u}}\fi}
\newunicodechar{ʸ}{\ifmmode^{y}\else\textsuperscript{\ensuremath{y}}\fi}
\newunicodechar{ᵃ}{\ifmmode^{a}\else\textsuperscript{\ensuremath{a}}\fi}
\newunicodechar{ᵉ}{\ifmmode^{e}\else\textsuperscript{\ensuremath{e}}\fi}
\newunicodechar{ᵏ}{\ifmmode^{k}\else\textsuperscript{\ensuremath{k}}\fi}
\newunicodechar{ᵖ}{\ifmmode^{p}\else\textsuperscript{\ensuremath{p}}\fi}
\newunicodechar{ᴺ}{\ifmmode^{N}\else\textsuperscript{\ensuremath{N}}\fi}
\newunicodechar{ᶜ}{\ifmmode^{c}\else\textsuperscript{\ensuremath{c}}\fi}
\newunicodechar{ʰ}{\ifmmode^{h}\else\textsuperscript{\ensuremath{h}}\fi}
\newunicodechar{ᵠ}{\ifmmode^{q}\else\textsuperscript{\ensuremath{q}}\fi}
\newunicodechar{ₙ}{\ifmmode_{n}\else\textsubscript{\ensuremath{n}}\fi}
\newunicodechar{ₐ}{\ifmmode_{a}\else\textsubscript{\ensuremath{a}}\fi}
\newunicodechar{ᵢ}{\ifmmode_{i}\else\textsubscript{\ensuremath{i}}\fi}
\newunicodechar{ᵣ}{\ifmmode_{r}\else\textsubscript{\ensuremath{r}}\fi}
\newunicodechar{ₖ}{\ifmmode_{k}\else\textsubscript{\ensuremath{k}}\fi}
\newunicodechar{ᵦ}{\ifmmode_{\beta}\else\textsubscript{\ensuremath{\beta}}\fi}
\newunicodechar{ₓ}{\ifmmode_{x}\else\textsubscript{\ensuremath{x}}\fi}
\newfontfamily\popdisplay{ArchivoBlack-Regular.ttf}[Path=../../../fonts/]
\newfontfamily\poplogo{SpaceGrotesk-Bold.ttf}[Path=../../../fonts/]
\newfontfamily\popsemi{PlusJakartaSans-SemiBold.ttf}[Path=../../../fonts/]
\newfontfamily\popxbold{PlusJakartaSans-ExtraBold.ttf}[Path=../../../fonts/]
\newfontfamily\popxboldls{PlusJakartaSans-ExtraBold.ttf}[Path=../../../fonts/,LetterSpace=3.0]

\setlist[enumerate]{leftmargin=1.7em,itemsep=5pt,topsep=4pt}
\setlist[itemize]{leftmargin=1.7em,itemsep=4pt,topsep=4pt,label={\color{poporange}\textbullet}}
\setlength{\parindent}{0pt}
\setlength{\parskip}{5pt}
\renewcommand{\arraystretch}{1.25}

% pgfplots : style Pop par défaut
\pgfplotsset{
  every axis/.append style={axis line style={popink,line width=1pt},tick style={popink},
    grid style={popline,line width=0.5pt},label style={font=\small},tick label style={font=\footnotesize}},
  popcurve/.style={poporange,line width=1.8pt,no markers,smooth},
}
% Même style disponible dans un \draw TikZ simple (hors pgfplots)
\tikzset{popcurve/.style={poporange,line width=1.8pt}}
\tikzset{popgrid/.style={popline,very thin}}
\colorlet{popcurve}{poporange} % le nom du style est parfois utilisé comme couleur

% ============================================================
% Pill / squiggle
% ============================================================
\newcommand{\poppill}[3]{%
  \tikz[baseline=-0.55ex]\node[draw=popink,line width=1.2pt,rounded corners=2.6pt,%
    fill=#2,inner xsep=6.5pt,inner ysep=3pt]{%
    \popxboldls\fontsize{7}{7}\selectfont\color{#3}\MakeUppercase{#1}};%
}
\newcommand{\popsquiggle}[1]{%
  \tikz[baseline=-0.4ex]{
    \draw[#1,line width=2.6pt,line cap=round]
      plot [smooth] coordinates {(0,0.09) (0.34,0.01) (0.68,0.21) (1.02,0.01) (1.36,0.10)};
  }%
}
% Mot-clé surligné (marqueur orange sous le texte)
\newcommand{\cle}[1]{{\popxbold\color{popdark}#1}}

% ============================================================
% En-tête / pied
% ============================================================
\pagestyle{fancy}
\fancyhf{}
\renewcommand{\headrulewidth}{0pt}
\renewcommand{\footrulewidth}{0pt}
\lhead{\raisebox{-0.15ex}{\poplogo\fontsize{13}{13}\selectfont\color{popink}OnBuch\color{poporange}+}}
\rhead{\poppill{\DOCMATIERE\ \textbullet\ \DOCNIVEAU\ \textbullet\ Leçon \DOCLECON}{white}{popink}}
\lfoot{\popxbold\fontsize{7.6}{7.6}\selectfont\color{popmuted}\MakeUppercase{Cours signé OnBuch+}\ \textbullet\ {\popsemi\DOCTITRE}}
\rfoot{\tikz[baseline=-0.6ex]\node[rounded corners=8.5pt,fill=popink,minimum height=17pt,minimum width=17pt,inner xsep=5pt,inner sep=0]{\popxbold\fontsize{8}{8}\selectfont\color{popcream}\thepage\,/\,\pageref*{LastPage}};}

% ============================================================
% Titres
% ============================================================
\newcommand{\chaptitre}[2]{%
  \begin{tcolorbox}[enhanced,colback=poporange,colframe=popink,boxrule=2.6pt,arc=11pt,
      drop shadow=popink,left=15pt,right=15pt,top=12pt,bottom=11pt,before skip=3pt,after skip=18pt]
    \poppill{#2}{popink}{popcream}\par\vspace{9pt}
    {\raggedright\hyphenpenalty=10000\exhyphenpenalty=10000\popdisplay\fontsize{22}{24}\selectfont\color{popcream}\MakeUppercase{#1}\par}
  \end{tcolorbox}
}
% Section numérotée : pastille orange avec le numéro + titre Archivo
\newcounter{popsec}
\newcommand{\coursec}[1]{%
  \stepcounter{popsec}\setcounter{popsub}{0}%
  \par\vspace{16pt}\noindent
  \begin{minipage}{\linewidth}
  \tikz[baseline=-0.7ex]\node[draw=popink,line width=1.6pt,fill=poporange,rounded corners=4pt,minimum size=22pt,inner sep=0]{\popdisplay\fontsize{12}{12}\selectfont\color{popcream}\thepopsec};%
  \hspace{8pt}{\popdisplay\fontsize{15}{17}\selectfont\color{popink}\MakeUppercase{#1}}\par
  \vspace{4pt}\noindent{\color{poporange}\rule{36pt}{3.6pt}}
  \end{minipage}\par\nopagebreak\vspace{8pt}
}
\newcounter{popsub}
\newcommand{\courssub}[1]{%
  \stepcounter{popsub}%
  \par\vspace{9pt}\noindent{\popxbold\fontsize{11.5}{13}\selectfont\color{popink}{\color{poporange}\thepopsec.\thepopsub}\hspace{6pt}#1}\par\nopagebreak\vspace{4pt}
}

% ============================================================
% Boîtes pédagogiques — même recette : fond blanc, bordure épaisse colorée,
% ombre dure encre, badge plein. Titre optionnel [#1].
% ============================================================
\tcbset{popbase/.style={breakable,enhanced,colback=white,boxrule=2.3pt,arc=9pt,
  shadow={3pt}{-3pt}{0pt}{popink},left=14pt,right=14pt,top=11pt,bottom=11pt,before skip=11pt,after skip=15pt,
  fontupper=\popsemi\fontsize{10.6}{14.5}\selectfont\color{popink},
  fontlower=\popsemi\fontsize{10.6}{14.5}\selectfont\color{popink}}}
% Coche dessinée (le glyphe ✓ n'existe pas dans les polices du document)
\newcommand{\popcheck}[1]{\tikz[baseline=-0.5ex]\draw[#1,line width=1.6pt,line cap=round,line join=round](-0.13,0.0)--(-0.04,-0.09)--(0.13,0.1);}
\providecommand{\coche}{\popcheck{popgreen}}
\providecommand{\resultat}[1]{\par\smallskip\noindent\fcolorbox{popgreen}{popgreenL}{\parbox{\dimexpr\linewidth-2\fboxsep-2\fboxrule\relax}{\textbf{Résultat.} #1}}\par\smallskip}
\newcommand{\popboxhead}[3]{\poppill{#1}{#2}{white}\ifx\relax#3\relax\else\hspace{6pt}{\popxbold\fontsize{10.6}{12}\selectfont\color{popink}#3}\fi\par\vspace{7pt}}

\newtcolorbox{aretenir}[1][]{popbase,colframe=popgreen,before upper={\popboxhead{À retenir}{popgreen}{#1}}}
% Texte en chinois (document de lecture, dialogue, extrait) : cadre
% d'étude. Un titre optionnel (source, auteur) s'affiche dans l'en-tête.
\newtcolorbox{textebox}[1][]{popbase,colframe=popblue,before upper={\popboxhead{文本 · Texte}{popblue}{#1}},after upper={}}
% Marques ✓ / ✗ (forme correcte / fautive) souvent utilisées par les modèles
\providecommand{\cross}{\textcolor{poppink}{\ensuremath{\times}}}
\providecommand{\xmark}{\cross}\providecommand{\crossmark}{\cross}\providecommand{\cmark}{\ensuremath{\checkmark}}
% Ligne de réponse / justification à compléter (inventée par les modèles)
\providecommand{\justification}[1]{\par\noindent\textit{Justification~:} \dotfill\par}
% \textipa (package tipa non chargé) : la police ne couvre pas l'API -> simple passage
\providecommand{\textipa}[1]{#1}
% Soulignement (ulem non chargé) : \uline, \ul
\providecommand{\uline}[1]{\underline{#1}}\providecommand{\ul}[1]{\underline{#1}}
% \glossaire{\item ...} inventée par les modèles : liste de glossaire
\providecommand{\glossaire}[1]{\begin{itemize}#1\end{itemize}}
% \alert (beamer) inventée par les modèles : texte mis en évidence
\providecommand{\alert}[1]{\textbf{\textcolor{poppink}{#1}}}
% variantes ASCII d'ordinaux écrites en macro
\providecommand{\iere}{\textsuperscript{re}}\providecommand{\ieme}{\textsuperscript{e}}\providecommand{\ere}{\textsuperscript{re}}\providecommand{\eme}{\textsuperscript{e}}\providecommand{\er}{\textsuperscript{er}}\providecommand{\ie}{\textsuperscript{e}}
% Ordinaux français écrits en macro par les modèles : 1\ière, 3\ième
\newcommand{\ière}{\textsuperscript{re}}\newcommand{\ième}{\textsuperscript{e}}
% \exemple (inventée par les modèles) : étiquette « Exemple : »
\providecommand{\exemple}{\textbf{Exemple~:}\ }
% \square utilisable aussi hors mode math (cases à cocher)
\AtBeginDocument{\let\squareMath\square\renewcommand{\square}{\ensuremath{\squareMath}}}
% Mot ou phrase en chinois à l'intérieur d'un paragraphe français
\newcommand{\zh}[1]{\ifmmode\text{#1}\else{#1}\fi}
\let\eng\zh \let\esp\zh \let\all\zh % alias tolérés
% flèches d'intonation inventées par les modèles
\providecommand{\textsearrow}{}\renewcommand{\textsearrow}{\ensuremath{\searrow}}\providecommand{\textnearrow}{}\renewcommand{\textnearrow}{\ensuremath{\nearrow}}
% \fr{..} : mot français (inventé par les modèles)
\providecommand{\fr}[1]{\textit{#1}}
% zhbox : encadré de texte chinois inventé par les modèles
\newenvironment{zhbox}{\begin{quote}}{\end{quote}}
% \courssubsub : titre de niveau 3 inventé par les modèles
\providecommand{\courssubsub}[1]{\par\medskip\noindent\textbf{#1}\par\smallskip}
% \zhbf : texte chinois en gras inventé par les modèles
\providecommand{\zhbf}[1]{\textbf{#1}}
% \vs : « versus » inventé par les modèles
\providecommand{\vs}{\textit{vs}\ }
% \blank : case à compléter inventée par les modèles
\providecommand{\blank}[1][]{\underline{\hspace{1.6em}}}
\newtcolorbox{exemplebox}[1][]{popbase,colframe=poporange,before upper={\popboxhead{Exemple}{poporange}{#1}}}
\newtcolorbox{definition}[1][]{popbase,colframe=popblue,before upper={\popboxhead{Définition}{popblue}{#1}}}
\newtcolorbox{propriete}[1][]{popbase,colframe=poppurple,before upper={\popboxhead{Propriété}{poppurple}{#1}}}
\newtcolorbox{methode}[1][]{popbase,colframe=popgold,before upper={\popboxhead{Méthode}{popgold}{#1}}}
\newtcolorbox{attention}[1][]{popbase,colframe=poppink,before upper={\popboxhead{Attention}{poppink}{#1}}}
\newtcolorbox{experience}[1][]{popbase,colframe=popblue,colback=popblueL,before upper={\popboxhead{Expérience}{popblue}{#1}}}
% « Le savais-tu ? » : carte violette pleine (contexte, culture, Cameroun)
\newtcolorbox{savaistu}[1][]{popbase,colback=poppurple,colframe=popink,
  fontupper=\popsemi\fontsize{10.4}{14.2}\selectfont\color{white},
  before upper={\poppill{Le savais-tu\,?}{white}{poppurple}\ifx\relax#1\relax\else\hspace{6pt}{\popxbold\color{white}#1}\fi\par\vspace{7pt}}}
% Exercice résolu : énoncé en haut, \tcblower puis la solution sur fond crème
\newcounter{popexo}\newcounter{popexr}
\newtcolorbox{exoresolu}[1][]{popbase,colframe=poporange,colbacklower=popcream,
  segmentation style={popink,line width=1.2pt,dashed},
  before upper={\stepcounter{popexr}\popboxhead{Exercice résolu \thepopexr}{poporange}{#1}},
  before lower={\poppill{Solution}{popink}{popcream}\par\vspace{6pt}}}
% Exercice (non résolu en place) — la correction est dans la partie Corrigés
\newtcolorbox{exercice}[1][]{popbase,colframe=popink,
  before upper={\stepcounter{popexo}\popboxhead{Exercice \thepopexo}{popink}{#1}}}
% Corrigé
\newtcolorbox{corrige}[1][]{popbase,colframe=popgreen,colback=popgreenL,
  before upper={\popboxhead{Corrigé}{popgreen}{#1}}}
% Objectifs / prérequis / bilan
\newtcolorbox{objectifs}{popbase,colframe=popink,colback=poporangeL,
  before upper={\popboxhead{Objectifs de la leçon}{popink}{}}}
\newtcolorbox{prerequis}{popbase,colframe=popink,colback=popblueL,
  before upper={\popboxhead{Prérequis}{popblue}{}}}
\newtcolorbox{bilan}{popbase,colframe=popink,colback=white,boxrule=2.8pt,
  before upper={\popboxhead{Fiche bilan}{poporange}{L'essentiel à connaître pour l'examen}}}
% Figure encadrée avec légende
\newtcolorbox{popfigure}[1][]{enhanced,breakable=false,colback=white,colframe=popline,boxrule=1.4pt,arc=8pt,
  left=10pt,right=10pt,top=10pt,bottom=8pt,before skip=10pt,after skip=12pt,halign=center,
  lower separated=false,halign lower=center,
  fontlower=\popsemi\fontsize{9}{11.5}\selectfont\color{popmuted},
  #1}
\newcommand{\legende}[1]{\par\vspace{6pt}{\popsemi\fontsize{9}{11.5}\selectfont\color{popmuted}#1}}

% ============================================================
% Signature OnBuch+
% ============================================================
\newcommand{\poplogoinline}[1]{{\poplogo\fontsize{#1}{#1}\selectfont\color{popink}OnBuch\color{poporange}+}}
\newcommand{\signatureonbuch}{%
  \begin{tcolorbox}[enhanced,colback=popink,colframe=popink,boxrule=2.2pt,arc=10pt,
      drop shadow=poporange,left=14pt,right=14pt,top=10pt,bottom=10pt,before skip=0pt,after skip=0pt]
    \tikz[baseline=-0.7ex]\node[circle,fill=popgreen,draw=popcream,line width=1.3pt,minimum size=22pt,inner sep=0]{\popcheck{white}};%
    \hspace{9pt}\begin{minipage}[c]{0.62\linewidth}
      {\popxbold\fontsize{12}{13}\selectfont\color{popcream}Signé\ \textbullet\ {\poplogo OnBuch\color{poporange}+}}\par\vspace{2pt}
      {\raggedright\popsemi\fontsize{8}{10.5}\selectfont\color{popline}\MakeUppercase{Équipe pédagogique OnBuch+}\\\MakeUppercase{Conforme au programme officiel MINESEC}\par}
    \end{minipage}\hfill
    {\poplogo\fontsize{22}{22}\selectfont\color{popcream}OnBuch\color{poporange}+}
  \end{tcolorbox}%
}

% ============================================================
% Page de garde
% #1 titre · #2 matière · #3 niveau · #4 module · #5 n° leçon · #6 durée ·
% #7 description / objectifs (texte ou itemize)
% ============================================================
\newcommand{\pagegarde}[7]{%
  \thispagestyle{empty}
  \newgeometry{a4paper,margin=1.7cm,top=1.6cm,bottom=1.4cm}%
  \begin{tikzpicture}[remember picture,overlay]
    \fill[poporange!18] ([xshift=-3.2cm,yshift=-3.6cm]current page.north east) circle (4.6cm);
    \fill[poppurple!14] ([xshift=2.4cm,yshift=7.6cm]current page.south west) circle (3.8cm);
  \end{tikzpicture}%
  {\poplogo\fontsize{30}{30}\selectfont\color{popink}OnBuch\color{poporange}+}\hfill
  \raisebox{0.6ex}{\poppill{Cours signé \textbullet\ Premium}{popink}{popcream}}\par
  \vspace{1pt}\popsquiggle{poppurple}\par
  \vspace{8pt}
  \begin{tcolorbox}[enhanced,colback=poporange,colframe=popink,boxrule=2.8pt,arc=13pt,
      drop shadow=popink,left=18pt,right=18pt,top=12pt,bottom=13pt,after skip=10pt]
    \poppill{#2 \textbullet\ #3}{popink}{popcream}\hspace{5pt}\poppill{#4}{white}{popink}\par\vspace{8pt}
    {\popdisplay\fontsize{14}{14}\selectfont\color{popink}LEÇON #5}\par\vspace{4pt}
    {\raggedright\hyphenpenalty=10000\exhyphenpenalty=10000\popdisplay\fontsize{25}{27}\selectfont\color{popcream}\MakeUppercase{#1}\par}
    \vspace{8pt}
    {\popxbold\fontsize{9.5}{9.5}\selectfont\color{popink}\MakeUppercase{Durée conseillée : #6}}
  \end{tcolorbox}
  \begin{tcolorbox}[enhanced,colback=white,colframe=popink,boxrule=2.2pt,arc=10pt,
      drop shadow=popink,left=16pt,right=16pt,top=10pt,bottom=10pt,after skip=8pt]
    \poppill{Dans cette leçon}{poppurple}{white}\par\vspace{5pt}
    \popsemi\fontsize{9.3}{12.6}\selectfont\color{popink}#7
  \end{tcolorbox}
  \noindent\begin{minipage}[t]{0.487\linewidth}
    \begin{tcolorbox}[enhanced,colback=popgreen,colframe=popink,boxrule=2.2pt,arc=10pt,
        drop shadow=popink,left=11pt,right=11pt,top=10pt,bottom=10pt,height=3.1cm,valign=center]
      \tcbox[colback=white,colframe=popink,boxrule=1.3pt,arc=5pt,boxsep=0pt,nobeforeafter,
        left=3pt,right=3pt,top=3pt,bottom=3pt,valign=center]{\qrcode[height=1.9cm]{https://play.google.com/store/apps/details?id=com.luvvix.onbuch&hl=fr}}%
      \hspace{8pt}\begin{minipage}[c]{0.5\linewidth}
        {\popdisplay\fontsize{11}{12.5}\selectfont\color{white}\MakeUppercase{Télécharge OnBuch}}\par\vspace{4pt}
        {\raggedright\popsemi\fontsize{8.4}{11}\selectfont\color{white}Gratuit \textbullet\ Google Play\\Tuteur IA, annales, concours, résultats\par}
      \end{minipage}
    \end{tcolorbox}
  \end{minipage}\hfill
  \begin{minipage}[t]{0.487\linewidth}
    \begin{tcolorbox}[enhanced,colback=poppurple,colframe=popink,boxrule=2.2pt,arc=10pt,
        drop shadow=popink,left=11pt,right=11pt,top=10pt,bottom=10pt,height=3.1cm,valign=center]
      \tikz[baseline=-0.7ex]\node[circle,fill=white,draw=popink,line width=1.3pt,minimum size=1.6cm,inner sep=0]{\popdisplay\fontsize{24}{24}\selectfont\color{poppurple}+};%
      \hspace{8pt}\begin{minipage}[c]{0.6\linewidth}
        {\popdisplay\fontsize{11}{12.5}\selectfont\color{white}\MakeUppercase{Passe à OnBuch+}}\par\vspace{4pt}
        {\raggedright\popsemi\fontsize{8.4}{11}\selectfont\color{white}Tous les cours signés, Tuteur IA illimité, fiches PDF — onglet OnBuch+ de l'app\par}
      \end{minipage}
    \end{tcolorbox}
  \end{minipage}
  \vspace{8pt}
  \popcontenu
  \vfill
  \signatureonbuch
  \restoregeometry
  \newpage
}

% Rangée « Ce cours contient » (page de garde)
\newcommand{\poptile}[3]{% #1 couleur #2 gros texte #3 légende
  \begin{tcolorbox}[enhanced,colback=#1,colframe=popink,boxrule=2pt,arc=9pt,shadow={2.5pt}{-2.5pt}{0pt}{popink},
      left=6pt,right=6pt,top=9pt,bottom=9pt,halign=center,valign=center,height=1.75cm,width=\linewidth,nobeforeafter]
    {\popdisplay\fontsize{12.5}{13}\selectfont\color{white}\MakeUppercase{#2}}\par\vspace{3pt}
    {\centering\popsemi\fontsize{7.8}{9.5}\selectfont\color{white}#3\par}
  \end{tcolorbox}}
\newcommand{\popcontenu}{%
  \par\noindent{\popxboldls\fontsize{7.5}{7.5}\selectfont\color{popmuted}CE COURS SIGNÉ CONTIENT}\par\vspace{7pt}
  \noindent\begin{minipage}[t]{0.235\linewidth}\poptile{popblue}{Cours}{complet, illustré, pas à pas}\end{minipage}\hfill
  \begin{minipage}[t]{0.235\linewidth}\poptile{poporange}{Méthodes}{et exercices résolus}\end{minipage}\hfill
  \begin{minipage}[t]{0.235\linewidth}\poptile{poppink}{Exercices}{type examen + corrigés}\end{minipage}\hfill
  \begin{minipage}[t]{0.235\linewidth}\poptile{popgreen}{Fiche bilan}{l'essentiel à retenir}\end{minipage}\par}

% Sommaire
\newtcolorbox{sommaire}{enhanced,colback=white,colframe=popink,boxrule=2.2pt,arc=10pt,
  drop shadow=popink,left=18pt,right=18pt,top=13pt,bottom=13pt,before skip=6pt,after skip=20pt,
  before upper={\poppill{Sommaire}{poppurple}{white}\par\vspace{9pt}\popsemi\fontsize{10.6}{14.5}\selectfont\color{popink}}}

% Signature de fin de cours — poussée tout en bas de la dernière page
\newcommand{\findecours}{%
  \par\vfill\nopagebreak
  \begin{center}\popsquiggle{poporange}\end{center}\vspace{4pt}
  \signatureonbuch
}
\AtBeginDocument{\def\llbracket{\mathopen{[\mkern-3.5mu[}}\def\rrbracket{\mathclose{]\mkern-3.5mu]}}} % intervalles d'entiers (glyphe ⟦ absent de la police)
% Glyphes absents de Fira Math : redéfinis à partir de glyphes disponibles
\AtBeginDocument{%
  \def\setminus{\mathbin{\backslash}}\let\smallsetminus\setminus
  \def\vdots{\mathord{\vbox{\baselineskip3.2pt\lineskiplimit0pt\kern4pt\hbox{.}\hbox{.}\hbox{.}}}}%
  \def\ddots{\mathinner{\mkern1mu\raise6.5pt\hbox{.}\mkern2mu\raise3.6pt\hbox{.}\mkern2mu\raise0.7pt\hbox{.}\mkern1mu}}%
  \def\triangle{\mathord{\scalebox{0.9}{$\symup{\Delta}$}}}%
  \def\nabla{\mathord{\raisebox{\depth}{\scalebox{1}[-1]{$\symup{\Delta}$}}}}%
  \def\checkmark{\popcheck{popdark}}%
  \def\star{\mathbin{\ast}}\def\bigstar{\mathord{\ast}}%
  \def\diamond{\mathbin{\scalebox{0.6}{$\Diamond$}}}%
  \def\bigcup{\mathop{\mathchoice{\raisebox{-0.3em}{\scalebox{1.7}{$\cup$}}}{\scalebox{1.25}{$\cup$}}{\cup}{\cup}}\displaylimits}%
  \def\bigcap{\mathop{\mathchoice{\raisebox{-0.3em}{\scalebox{1.7}{$\cap$}}}{\scalebox{1.25}{$\cap$}}{\cap}{\cap}}\displaylimits}%
  \def\lesseqqgtr{\mathrel{\lessgtr}}\def\bigoplus{\mathop{\oplus}\displaylimits}%
}
\providecommand{\ssi}{si et seulement si} % abréviation fréquente
\AtBeginDocument{\providecommand{\wideparen}{\overparen}} % arc de cercle

%% FIGFIX-BEGIN v3 — correctifs globaux des figures (docs/onbuch-generation/tools/figfix)
\usepackage{adjustbox}\usepackage{etoolbox}
% toute figure TikZ/pgfplots plus large que la ligne est réduite à la largeur disponible
\BeforeBeginEnvironment{tikzpicture}{\begin{adjustbox}{max width=\linewidth}}
\AfterEndEnvironment{tikzpicture}{\end{adjustbox}}
% légendes pgfplots lisibles par-dessus les courbes
\pgfplotsset{every axis/.append style={legend style={fill=white,fill opacity=0.92,text opacity=1,draw=black!15,inner sep=3pt}}}
% étiquettes d'arêtes (syntaxe edge["..."]) avec fond blanc
\tikzset{every edge quotes/.append style={fill=white,inner sep=1.2pt}}
% bulles de cartes mentales : texte à la ligne dans la bulle, bulle un peu plus grande
\tikzset{every concept/.append style={text width=2.0cm,align=center,minimum size=2.3cm,inner sep=2pt}}
%% FIGFIX-END
\begin{document}

\pagegarde{Leçon 40 — Grammaire : 一边……一边…….. ; la ponctuation en chinois}{Chinois}{1ère A}{Module 5 : Mass médias et télécommunication}{ZHA40}{2 h}{Cette leçon de grammaire présente la structure 一边……一边…… qui exprime deux actions simultanées, ainsi que les signes de ponctuation propres à l'écriture chinoise. L'élève apprend à construire, employer et ponctuer correctement des phrases décrivant des actions parallèles, en lien avec le thème des médias et de la communication.
\vspace{4pt}
\begin{multicols}{2}\small
\begin{itemize}
\item À la fin de cette leçon, je sais identifier la structure 一边……一边…… dans un texte chinois
\item À la fin de cette leçon, je sais construire correctement une phrase avec 一边……一边…… (sujet + 一边 + verbe1, 一边 + verbe2)
\item À la fin de cette leçon, je sais distinguer les emplois corrects et incorrects de la structure (actions simultanées réalisables, même sujet)
\item À la fin de cette leçon, je sais reconnaître et utiliser les principaux signes de ponctuation chinois : 。，、；：？！“”《》……
\item À la fin de cette leçon, je sais comparer la ponctuation chinoise et la ponctuation française
\item À la fin de cette leçon, je sais éviter les erreurs fréquentes des francophones (place du sujet, oubli du second 一边, virgule 、 mal employée)
\end{itemize}\end{multicols}}

\begin{sommaire}
\begin{multicols}{2}
\begin{enumerate}
\item Observation et découverte de la structure
\item Schéma et construction de la structure
\item Emplois, limites et exceptions
\item Comparaison avec le français et pièges des francophones
\item La ponctuation en chinois
\item Entraînement intégré : structure + ponctuation
\item Activité d'intégration
\item Méthodes et exercices résolus
\item Exercices
\item Corrigés des exercices
\item Fiche bilan
\end{enumerate}
\end{multicols}
\end{sommaire}

\begin{objectifs}
\begin{itemize}
\item À la fin de cette leçon, je sais identifier la structure 一边……一边…… dans un texte chinois
\item À la fin de cette leçon, je sais construire correctement une phrase avec 一边……一边…… (sujet + 一边 + verbe1, 一边 + verbe2)
\item À la fin de cette leçon, je sais distinguer les emplois corrects et incorrects de la structure (actions simultanées réalisables, même sujet)
\item À la fin de cette leçon, je sais reconnaître et utiliser les principaux signes de ponctuation chinois : 。，、；：？！“”《》……
\item À la fin de cette leçon, je sais comparer la ponctuation chinoise et la ponctuation française
\item À la fin de cette leçon, je sais éviter les erreurs fréquentes des francophones (place du sujet, oubli du second 一边, virgule 、 mal employée)
\item À la fin de cette leçon, je sais transformer des phrases simples en phrases à actions simultanées
\item À la fin de cette leçon, je sais rédiger un court paragraphe ponctué correctement sur le thème des médias
\end{itemize}
\end{objectifs}

\begin{prerequis}
\begin{itemize}
\item Les pronoms personnels et la structure de base Sujet + Verbe + Complément (Seconde)
\item Les verbes d'action courants : 看, 听, 说, 吃, 写, 唱, 学 (Seconde)
\item Le vocabulaire des médias du Module 5 : 电视, 手机, 报纸, 网络, 新闻
\item L'adverbe 在 pour exprimer une action en cours (leçons précédentes de 1ère)
\item La négation 不 et 没 (Seconde)
\end{itemize}
\end{prerequis}

\newpage

\coursec{Observation et découverte de la structure}

\courssub{Situation d'accroche : la pause au lycée de Yaoundé}

Pendant la pause au lycée de Yaoundé, Amina écoute de la musique sur son téléphone tout en révisant ses leçons. Son ami Kevin, lui, mange des beignets tout en regardant un match des Lions Indomptables à la télévision. En français, on dit « en même temps » ou « tout en ». Mais comment dire cela en chinois ? Comment un journaliste chinois de CCTV décrirait-il ces deux actions simultanées ? Et comment ponctuer correctement ses phrases dans un article ? C'est ce que nous allons découvrir aujourd'hui.

\textbf{Problématique :} comment exprimer en chinois deux actions réalisées \emph{en même temps} par \emph{la même personne} ? Et comment ponctuer correctement une phrase chinoise à l'écrit ?

\courssub{Lecture guidée : quatre exemples à observer}

Lisons d'abord les quatre phrases suivantes, tirées de la vie quotidienne et du monde des médias. Lisez-les à voix haute, en faisant attention aux tons, puis observez-les attentivement.

\begin{textebox}[Exemple 1 — À la maison]
他一边看电视，一边吃饭。\\
\emph{Tā yìbiān kàn diànshì, yìbiān chīfàn.}\\
« Il regarde la télévision tout en mangeant. »
\end{textebox}

\begin{textebox}[Exemple 2 — Pendant les devoirs]
我们一边听音乐，一边做作业。\\
\emph{Wǒmen yìbiān tīng yīnyuè, yìbiān zuò zuòyè.}\\
« Nous écoutons de la musique tout en faisant nos devoirs. »
\end{textebox}

\begin{textebox}[Exemple 3 — Dans la cuisine]
妈妈一边打电话，一边做饭。\\
\emph{Māma yìbiān dǎ diànhuà, yìbiān zuòfàn.}\\
« Maman téléphone tout en préparant le repas. »
\end{textebox}

\begin{textebox}[Exemple 4 — En classe de chinois]
学生们一边读报纸，一边记生词。\\
\emph{Xuéshengmen yìbiān dú bàozhǐ, yìbiān jì shēngcí.}\\
« Les élèves lisent le journal tout en notant les mots nouveaux. »
\end{textebox}

\courssub{Vocabulaire des exemples}

Avant d'observer la structure, assurons-nous de bien comprendre chaque mot.

\begin{center}
\begin{tabularx}{\linewidth}{|X|X|X|X|}
\hline
\rowcolor{popblueL} Caractères & Pinyin & Nature & Traduction \\
\hline
一边……一边…… & yìbiān... yìbiān... & corrélatif & tout en..., en même temps... \\
\hline
看 & kàn & verbe & regarder \\
\hline
电视 & diànshì & nom & télévision \\
\hline
吃饭 & chīfàn & verbe & manger (prendre un repas) \\
\hline
听 & tīng & verbe & écouter \\
\hline
音乐 & yīnyuè & nom & musique \\
\hline
做作业 & zuò zuòyè & verbe + nom & faire ses devoirs \\
\hline
打电话 & dǎ diànhuà & verbe + nom & téléphoner \\
\hline
做饭 & zuòfàn & verbe & préparer le repas, cuisiner \\
\hline
读 & dú & verbe & lire (à voix haute), étudier \\
\hline
报纸 & bàozhǐ & nom & journal, presse écrite \\
\hline
记 & jì & verbe & noter, mémoriser \\
\hline
生词 & shēngcí & nom & mot nouveau, vocabulaire nouveau \\
\hline
\end{tabularx}
\end{center}

\begin{attention}[Le mot 音乐 : attention au caractère 乐]
Le caractère \zh{乐} a deux lectures : \zh{lè} dans \zh{快乐} (kuàilè, « joyeux ») et \zh{yuè} dans \zh{音乐} (yīnyuè, « musique »). On dit qu'il est \cle{polyphonique}. Retenez : la musique, c'est \zh{yuè}, pas \zh{lè} !
\end{attention}

\courssub{Questions d'observation}

\begin{methode}[Comment observer une structure grammaticale nouvelle]
Face à des exemples nouveaux, le bon sinisant ne mémorise pas : il \emph{observe}. Voici la démarche en quatre étapes :
\begin{enumerate}
\item \textbf{Repérer le mot qui se répète} dans toutes les phrases : c'est souvent lui, le mot-clé de la structure.
\item \textbf{Souligner les verbes} et noter leur place dans la phrase.
\item \textbf{Vérifier le sujet} : y a-t-il un seul sujet ou plusieurs ? Où est-il placé ?
\item \textbf{Formuler la règle avec ses propres mots}, en français d'abord, puis sous forme de schéma.
\end{enumerate}
\end{methode}

Appliquons cette méthode à nos quatre exemples.

\textbf{Question 1 — Quel mot se répète ?}

Dans les quatre phrases, le mot \zh{一边} (yìbiān) apparaît \emph{deux fois}. C'est un mot \cle{corrélé} : il fonctionne par paire, comme « d'une part... d'autre part... » en français. Le premier \zh{一边} annonce la première action, le second \zh{一边} annonce la seconde action.

\textbf{Question 2 — Où se placent les verbes ?}

Chaque \zh{一边} est immédiatement suivi d'un \cle{verbe} (ou d'un groupe verbal) : \zh{一边看电视}, \zh{一边吃饭}. Jamais de nom, jamais d'adjectif seul après \zh{一边} : cette structure n'encadre que des \emph{actions}.

\textbf{Question 3 — Combien d'actions ? Qui les fait ?}

Chaque phrase contient \textbf{deux actions}, et c'est \textbf{le même sujet} qui les accomplit toutes les deux : dans l'exemple 3, c'est \zh{妈妈} qui téléphone \emph{et} qui cuisine. Le sujet n'est dit qu'\emph{une seule fois}, au tout début de la phrase.

\textbf{Question 4 — Quel rapport de temps entre les deux actions ?}

Les deux actions se déroulent \emph{en même temps} : elles sont \cle{simultanées}. Maman ne téléphone pas d'abord puis ne cuisine pas ensuite : elle fait les deux ensemble.

\begin{popfigure}
\begin{tikzpicture}[font=\small]
\node[draw=popdark, fill=popblueL, rounded corners, align=center, minimum width=2.6cm, minimum height=1cm] (sujet) at (0,0) {\textbf{Sujet}\\他 (tā)};
\node[draw=popblue, fill=popblueL!60, rounded corners, align=center, minimum width=4.2cm, minimum height=1cm] (act1) at (7,1.2) {一边 + \textbf{action 1}\\看电视 (kàn diànshì)};
\node[draw=poporange, fill=poporangeL, rounded corners, align=center, minimum width=4.2cm, minimum height=1cm] (act2) at (7,-1.2) {一边 + \textbf{action 2}\\吃饭 (chīfàn)};
\draw[-{Stealth[length=3mm]}, very thick, popblue] (sujet.east) -- (act1.west) node[midway, above, sloped, popblue]{en même temps};
\draw[-{Stealth[length=3mm]}, very thick, poporange] (sujet.east) -- (act2.west) node[midway, below, sloped, poporange]{en même temps};
\end{tikzpicture}
\legende{Un seul sujet, deux actions parallèles : les deux flèches partent du même point et avancent ensemble dans le temps.}
\end{popfigure}

\courssub{Dégager collectivement la règle}

Grâce à nos observations, nous pouvons maintenant formuler la règle.

\begin{definition}[La structure 一边……一边……]
\zh{一边……一边……} (yìbiān... yìbiān...) est une \cle{structure corrélative} qui indique que \textbf{deux actions (ou plus) sont accomplies simultanément par le même sujet}. Elle correspond au français « \textbf{tout en + gérondif} » ou « \textbf{en même temps} ».
\begin{itemize}
\item Schéma : \textbf{Sujet + 一边 + verbe 1, 一边 + verbe 2.}
\item Le sujet est placé \emph{une seule fois}, avant le premier \zh{一边}.
\item Les deux actions ont la même importance : aucune n'est secondaire.
\end{itemize}
\end{definition}

\begin{exemplebox}[Exemple supplémentaire]
她一边唱歌，一边跳舞。\\
\emph{Tā yìbiān chànggē, yìbiān tiàowǔ.}\\
« Elle chante tout en dansant. »\\[4pt]
Autres exemples :\\
老师一边讲课，一边写字。\emph{(Lǎoshī yìbiān jiǎngkè, yìbiān xiězì.)} « Le professeur explique la leçon tout en écrivant. »\\
他一边开车，一边听广播。\emph{(Tā yìbiān kāichē, yìbiān tīng guǎngbō.)} « Il conduit tout en écoutant la radio. »
\end{exemplebox}

Comparons maintenant avec le français :

\begin{center}
\begin{tabularx}{\linewidth}{|X|X|X|}
\hline
\rowcolor{popblueL} Français & Chinois (caractères + pinyin) & Observation \\
\hline
Il mange \emph{en regardant} la télé. & 他一边吃饭，一边看电视。Tā yìbiān chīfàn, yìbiān kàn diànshì. & Le gérondif français devient un verbe plein après 一边. \\
\hline
Elle chante \emph{tout en} dansant. & 她一边唱歌，一边跳舞。Tā yìbiān chànggē, yìbiān tiàowǔ. & « Tout en » = 一边……一边…… (obligatoirement répété). \\
\hline
Nous travaillons \emph{en même temps}. & 我们一边工作，一边学习。Wǒmen yìbiān gōngzuò, yìbiān xuéxí. & Le chinois nomme les deux actions explicitement. \\
\hline
\end{tabularx}
\end{center}

\begin{attention}[Piège n° 1 des francophones : ne pas oublier le second 一边]
En français, on peut dire « il mange en regardant la télé » avec un seul mot-outil (« en »). En chinois, la structure est \textbf{par paire} : on ne peut pas dire \zh{他一边吃饭，看电视}. Les deux \zh{一边} sont nécessaires, comme les deux moitiés d'une paire de chaussures. Piège n° 2 : le sujet ne se répète pas — on ne dit pas \zh{他一边吃饭，他一边看电视}.
\end{attention}

\begin{popfigure}
\begin{tikzpicture}[mindmap, grow cyclic, every node/.style={concept, align=center, font=\small},
root concept/.append style={concept color=popblue, text=white, font=\bfseries},
level 1/.append style={level distance=3.4cm, sibling angle=90}]
\node[root concept, align=center]{一边\\yìbiān}
  child[concept color=poporange] { node[align=center]{place :\\avant chaque\\verbe} }
  child[concept color=popgreen] { node[align=center]{sujet unique\\placé avant\\le 1er 一边} }
  child[concept color=poppurple] { node[align=center]{actions\\simultanées\\(2 ou +)} }
  child[concept color=popgold] { node[align=center]{synonyme :\\一面……一面……\\yìmiàn... yìmiàn...} };
\end{tikzpicture}
\legende{Carte mentale de la structure 一边……一边…… : ses quatre caractéristiques essentielles.}
\end{popfigure}

\begin{savaistu}[Le synonyme 一面……一面……]
Il existe une variante plus littéraire : \zh{一面……一面……} (yìmiàn... yìmiàn...), que l'on rencontre surtout dans les articles de presse et les textes écrits. Le sens est identique : \zh{他一面看报，一面喝茶} (Tā yìmiàn kànbào, yìmiàn hēchá, « il lit le journal tout en buvant son thé »). À l'oral et à votre niveau, utilisez \zh{一边……一边……}, la forme la plus courante.
\end{savaistu}

\begin{exoresolu}[Observer et dégager la règle]
\textbf{Énoncé.} Voici une phrase nouvelle : \zh{记者一边采访，一边录音。} (Jìzhě yìbiān cǎifǎng, yìbiān lùyīn.) 1) Quel mot se répète ? 2) Quels sont les verbes ? 3) Qui fait les actions ? 4) Traduisez la phrase.
\tcblower
\textbf{Solution détaillée.}
\begin{enumerate}
\item Le mot qui se répète est \zh{一边} (yìbiān), deux fois : c'est le signal de la structure corrélative.
\item Les verbes sont \zh{采访} (cǎifǎng, « interviewer ») et \zh{录音} (lùyīn, « enregistrer »), chacun placé juste après un \zh{一边}.
\item Le sujet unique est \zh{记者} (jìzhě, « le journaliste »), placé au début, avant le premier \zh{一边} : c'est lui qui interviewe \emph{et} qui enregistre.
\item Traduction : « Le journaliste interviewe tout en enregistrant. » Les deux actions sont simultanées : la structure 一边……一边…… est donc bien respectée.
\end{enumerate}
\end{exoresolu}

\begin{exercice}[À vous d'observer]
Observez la phrase suivante et répondez aux quatre questions de la méthode : \zh{爸爸一边开车，一边打电话。} (Bàba yìbiān kāichē, yìbiān dǎ diànhuà.) Puis traduisez-la en français avec « tout en ».
\end{exercice}

\begin{corrige}[Exercice « À vous d'observer »]
1) Le mot répété est \zh{一边}. 2) Les verbes sont \zh{开车} (kāichē, « conduire ») et \zh{打电话} (dǎ diànhuà, « téléphoner »). 3) Le sujet unique est \zh{爸爸} (bàba, « papa »), placé avant le premier \zh{一边}. 4) Traduction : « Papa conduit tout en téléphonant. » (Au passage : au Cameroun comme en Chine, c'est dangereux et interdit !)
\end{corrige}

\begin{aretenir}[L'essentiel de cette première étape]
\begin{itemize}
\item \zh{一边……一边……} exprime \textbf{deux actions simultanées} accomplies par \textbf{un seul sujet}.
\item Schéma : \textbf{Sujet + 一边 + verbe 1, 一边 + verbe 2.}
\item Équivalents français : « tout en + gérondif », « en même temps ».
\item Les deux \zh{一边} sont obligatoires ; le sujet ne se dit qu'une fois, au début.
\item Variante écrite : \zh{一面……一面……}.
\end{itemize}
\end{aretenir}

\coursec{Schéma et construction de la structure}

\courssub{Le schéma général}

\begin{propriete}[Construction de 一边……一边……]
\textbf{Schéma de base :} Sujet + 一边 + verbe 1 (+ complément) ，一边 + verbe 2 (+ complément) 。
\begin{itemize}
\item Le \cle{sujet} (nom ou pronom) ouvre la phrase, une seule fois.
\item Chaque \zh{一边} est \emph{collé} au verbe qui le suit : pas de mot entre \zh{一边} et son verbe.
\item Les deux segments sont séparés par une \cle{virgule chinoise} \zh{，}.
\item On peut enchaîner trois actions : \zh{他一边吃饭，一边看报，一边喝茶。} (Tā yìbiān chīfàn, yìbiān kànbào, yìbiān hēchá. « Il mange, lit le journal et boit son thé en même temps. »)
\end{itemize}
\end{propriete}

\courssub{Construire pas à pas}

\begin{methode}[Fabriquer une phrase avec 一边……一边…… en quatre étapes]
\begin{enumerate}
\item \textbf{Choisir le sujet} unique : \zh{我} (wǒ, « je »).
\item \textbf{Choisir deux actions compatibles} (qu'on peut faire en même temps) : \zh{走路} (zǒulù, « marcher ») et \zh{唱歌} (chànggē, « chanter »).
\item \textbf{Placer 一边 devant chaque verbe} : \zh{一边走路} + \zh{一边唱歌}.
\item \textbf{Assembler} avec la virgule chinoise et le point final : \zh{我一边走路，一边唱歌。} (Wǒ yìbiān zǒulù, yìbiān chànggē. « Je chante tout en marchant. »)
\end{enumerate}
\end{methode}

\begin{exemplebox}[Phrases construites avec ce schéma]
妹妹一边哭，一边说。\emph{(Mèimei yìbiān kū, yìbiān shuō.)} « La petite sœur parle tout en pleurant. »\\
我们一边吃饭，一边谈工作。\emph{(Wǒmen yìbiān chīfàn, yìbiān tán gōngzuò.)} « Nous parlons du travail tout en mangeant. »\\
他一边上网，一边和朋友聊天。\emph{(Tā yìbiān shàngwǎng, yìbiān hé péngyou liáotiān.)} « Il discute avec ses amis tout en naviguant sur Internet. »
\end{exemplebox}

\courssub{Emplois, limites et exceptions}

\begin{propriete}[Conditions d'emploi de 一边……一边……]
\begin{enumerate}
\item \textbf{Même sujet obligatoire.} Si deux personnes différentes font deux actions en même temps, on n'utilise pas cette structure ; on dit par exemple : \zh{我看电视的时候，妈妈在做饭。} (Wǒ kàn diànshì de shíhou, māma zài zuòfàn. « Pendant que je regarde la télé, maman cuisine. »)
\item \textbf{Deux actions volontaires et compatibles.} Les deux verbes doivent désigner des actions que l'on peut réellement faire ensemble. On ne dira pas \zh{他一边睡觉，一边跑步} (« il dort tout en courant ») : c'est impossible !
\item \textbf{Pas d'adjectif après 一边.} Pour deux états simultanés, on utilise \zh{又……又……} (yòu... yòu...) : \zh{这个菜又辣又好吃。} (Zhège cài yòu là yòu hǎochī. « Ce plat est à la fois piquant et délicieux. »)
\item \textbf{Aspect.} On n'ajoute en général pas \zh{了} ni \zh{过} directement après les verbes de la structure, car elle décrit des actions en cours, non accomplies.
\end{enumerate}
\end{propriete}

\begin{attention}[一边……一边…… ou 又……又…… ?]
Les francophones confondent souvent ces deux structures corrélatives. Le test est simple :
\begin{itemize}
\item Deux \textbf{actions} (verbes d'action) simultanées $\rightarrow$ \zh{一边……一边……} : \zh{他一边走，一边看。} « Il regarde tout en marchant. »
\item Deux \textbf{qualités ou états} (adjectifs) simultanés $\rightarrow$ \zh{又……又……} : \zh{这个手机又便宜又好。} (Zhège shǒujī yòu piányi yòu hǎo. « Ce téléphone est à la fois bon marché et bon. »)
\end{itemize}
\end{attention}

\coursec{Comparaison avec le français et pièges des francophones}

\courssub{Tableau récapitulatif des correspondances}

\begin{center}
\begin{tabularx}{\linewidth}{|X|X|X|}
\hline
\rowcolor{popblueL} Structure & Exemple en caractères + pinyin & Traduction \\
\hline
Sujet + 一边 + V1, 一边 + V2 & 她一边听音乐，一边做作业。Tā yìbiān tīng yīnyuè, yìbiān zuò zuòyè. & Elle écoute de la musique tout en faisant ses devoirs. \\
\hline
Avec compléments d'objet & 他一边喝咖啡，一边看报纸。Tā yìbiān hē kāfēi, yìbiān kàn bàozhǐ. & Il boit son café tout en lisant le journal. \\
\hline
Trois actions & 孩子们一边唱，一边跳，一边笑。Háizimen yìbiān chàng, yìbiān tiào, yìbiān xiào. & Les enfants chantent, dansent et rient en même temps. \\
\hline
Variante écrite 一面……一面…… & 他一面工作，一面学习。Tā yìmiàn gōngzuò, yìmiàn xuéxí. & Il travaille tout en étudiant. \\
\hline
\end{tabularx}
\end{center}

\courssub{Les quatre erreurs typiques des élèves francophones}

\begin{attention}[Erreurs à éviter absolument]
\begin{enumerate}
\item \textbf{Oublier le second 一边} : \zh{他一边吃饭，看电视} est fautif. En français un seul « en » suffit ; en chinois, il faut les deux moitiés de la paire.
\item \textbf{Répéter le sujet} : \zh{他一边吃饭，他一边看电视} est fautif. Le sujet ne se dit qu'une fois, au début.
\item \textbf{Mettre un adjectif après 一边} : \zh{她一边漂亮，一边聪明} est fautif ; pour deux qualités, utilisez \zh{又……又……} : \zh{她又漂亮又聪明。}
\item \textbf{Traduire mot à mot « en même temps »} : le français « il mange et il regarde la télé en même temps » ne se traduit pas par \zh{同时} à votre niveau ; la structure naturelle est bien \zh{一边……一边……}.
\end{enumerate}
\end{attention}

\begin{exoresolu}[Corriger une phrase fautive]
\textbf{Énoncé.} Un élève a écrit : \zh{我们一边学习，我们一边说话。} Relevez l'erreur et corrigez la phrase. Traduisez ensuite la phrase corrigée.
\tcblower
\textbf{Solution détaillée.} L'erreur est la répétition du sujet \zh{我们} devant le second \zh{一边}. Dans la structure 一边……一边……, le sujet unique n'apparaît qu'une seule fois, en tête de phrase. Phrase corrigée : \zh{我们一边学习，一边说话。} (Wǒmen yìbiān xuéxí, yìbiān shuōhuà.) Traduction : « Nous parlons tout en étudiant. »
\end{exoresolu}

\coursec{La ponctuation en chinois}

\courssub{Pourquoi apprendre la ponctuation chinoise ?}

Un journaliste de CCTV ou un élève chinois n'écrit jamais avec la ponctuation française. Le chinois possède ses propres signes, \emph{pleine chasse} (ils occupent la largeur d'un caractère). Les utiliser correctement est indispensable à l'écrit, notamment dans les copies d'examen.

\begin{definition}[Les principaux signes de ponctuation chinois]
\begin{itemize}
\item \zh{。} \cle{le point} (jùhào 句号) : termine une phrase déclarative. C'est un petit rond, pas un point.
\item \zh{，} \cle{la virgule} (dòuhào 逗号) : sépare les parties d'une phrase, comme entre les deux \zh{一边}.
\item \zh{、} \cle{la virgule d'énumération} (dùnhào 顿号) : sépare des mots ou groupes de mots coordonnés : \zh{苹果、香蕉、橙子} (píngguǒ, xiāngjiāo, chéngzi « pommes, bananes, oranges »).
\item \zh{？} \cle{le point d'interrogation} : termine une question : \zh{你好吗？}
\item \zh{！} \cle{le point d'exclamation} : exprime l'étonnement, l'ordre, l'émotion : \zh{太好了！} (Tài hǎo le ! « C'est formidable ! »)
\item \zh{：} \cle{les deux-points} : introduisent une explication ou une citation.
\item \zh{；} \cle{le point-virgule} : sépare deux propositions longues et parallèles.
\item \zh{“ ”} \cle{les guillemets} : citent une parole ou un titre d'œuvre courte.
\item \zh{《 》} \cle{les guillemets de titre} : encadrent les titres de livres, journaux, films : \zh{《人民日报》} (Rénmín Rìbào, « le Quotidien du Peuple »).
\item \zh{……} \cle{les points de suspension} : indiquent une énumération inachevée ou une hésitation.
\end{itemize}
\end{definition}

\begin{attention}[Différences majeures avec le français]
\begin{enumerate}
\item Le point chinois est un \textbf{rond} \zh{。}, jamais un point « . » dans un texte en caractères.
\item Pour énumérer des \textbf{mots}, on n'utilise pas la virgule \zh{，} mais \zh{、} : on écrit \zh{我喜欢唱歌、跳舞和画画。} (Wǒ xǐhuan chànggē, tiàowǔ hé huàhuà. « J'aime chanter, danser et dessiner. »)
\item Les titres de journaux et de livres prennent \zh{《》}, pas des guillemets français ni l'italique.
\item Pas d'espace avant la ponctuation : en chinois, le signe est \emph{collé} au caractère qui précède.
\end{enumerate}
\end{attention}

\begin{exemplebox}[Un mini-texte bien ponctué]
晚上，我们一家人都很忙：爸爸一边看《人民日报》，一边喝茶；妈妈一边做饭，一边听广播；我一边做作业，一边听英语。弟弟呢？他一边吃苹果，一边看电视！\\
\emph{Wǎnshang, wǒmen yìjiārén dōu hěn máng: bàba yìbiān kàn 《Rénmín Rìbào》, yìbiān hēchá; māma yìbiān zuòfàn, yìbiān tīng guǎngbō; wǒ yìbiān zuò zuòyè, yìbiān tīng Yīngyǔ. Dìdi ne? Tā yìbiān chī píngguǒ, yìbiān kàn diànshì!}\\
« Le soir, toute la famille est occupée : papa lit le Quotidien du Peuple tout en buvant son thé ; maman cuisine tout en écoutant la radio ; moi, je fais mes devoirs tout en écoutant l'anglais. Et mon petit frère ? Il mange une pomme tout en regardant la télé ! »
\end{exemplebox}

Observez dans ce texte : les deux-points \zh{：} qui annoncent la liste, les points-virgules \zh{；} qui séparent les membres de la famille, le point d'interrogation \zh{？}, le point d'exclamation \zh{！}, les guillemets de titre \zh{《》} et les points \zh{。} finaux.

\coursec{Entraînement intégré : structure + ponctuation}

\courssub{Exercices d'application}

\begin{exoresolu}[Traduire en chinois avec la bonne ponctuation]
\textbf{Énoncé.} Traduisez en chinois : « À la maison, mon grand frère écoute la radio tout en réparant son téléphone. »
\tcblower
\textbf{Solution détaillée.}
\begin{enumerate}
\item Sujet unique : \zh{我哥哥} (wǒ gēge, « mon grand frère »).
\item Complément de lieu en tête : \zh{在家里} (zài jiālǐ, « à la maison »), suivi de la virgule \zh{，}.
\item Deux actions simultanées : \zh{听广播} (tīng guǎngbō, « écouter la radio ») et \zh{修手机} (xiū shǒujī, « réparer le téléphone »).
\item Assemblage : \zh{在家里，我哥哥一边听广播，一边修手机。} (Zài jiālǐ, wǒ gēge yìbiān tīng guǎngbō, yìbiān xiū shǒujī.) N'oubliez ni le second \zh{一边}, ni le point rond final \zh{。}.
\end{enumerate}
\end{exoresolu}

\begin{exercice}[Exercice 1 — Complétez avec 一边……一边……]
Complétez les phrases avec les verbes entre parenthèses :
\begin{enumerate}
\item 她……，……。（唱歌 / 跳舞）
\item 我们……，……。（吃饭 / 看电视）
\item 老师……，……。（讲课 / 写字）
\end{enumerate}
\end{exercice}

\begin{corrige}[Exercice 1]
1) \zh{她一边唱歌，一边跳舞。} (Tā yìbiān chànggē, yìbiān tiàowǔ.) « Elle chante tout en dansant. » 2) \zh{我们一边吃饭，一边看电视。} (Wǒmen yìbiān chīfàn, yìbiān kàn diànshì.) « Nous regardons la télé tout en mangeant. » 3) \zh{老师一边讲课，一边写字。} (Lǎoshī yìbiān jiǎngkè, yìbiān xiězì.) « Le professeur explique la leçon tout en écrivant. »
\end{corrige}

\begin{exercice}[Exercice 2 — Corrigez les phrases fautives]
Chaque phrase contient une erreur. Corrigez-la et traduisez la phrase corrigée.
\begin{enumerate}
\item 他一边喝茶，看报纸。
\item 妈妈一边做饭，妈妈一边打电话。
\item 这个菜一边辣，一边好吃。
\end{enumerate}
\end{exercice}

\begin{corrige}[Exercice 2]
1) Il manque le second \zh{一边} : \zh{他一边喝茶，一边看报纸。} (Tā yìbiān hēchá, yìbiān kàn bàozhǐ.) « Il lit le journal tout en buvant son thé. » 2) Le sujet est répété : \zh{妈妈一边做饭，一边打电话。} (Māma yìbiān zuòfàn, yìbiān dǎ diànhuà.) « Maman téléphone tout en cuisinant. » 3) Deux qualités (adjectifs) : il faut \zh{又……又……} : \zh{这个菜又辣又好吃。} (Zhège cài yòu là yòu hǎochī.) « Ce plat est à la fois piquant et délicieux. »
\end{corrige}

\begin{exercice}[Exercice 3 — Ponctuez le texte]
Recopiez ce texte en ajoutant la ponctuation chinoise correcte (point, virgules, deux-points, point d'exclamation) :\\
\zh{晚上 弟弟一边吃水果 一边看电视 妈妈一边做饭 一边听音乐 我们家真热闹}
\end{exercice}

\begin{corrige}[Exercice 3]
\zh{晚上，弟弟一边吃水果，一边看电视；妈妈一边做饭，一边听音乐。我们家真热闹！}\\
\emph{Wǎnshang, dìdi yìbiān chī shuǐguǒ, yìbiān kàn diànshì; māma yìbiān zuòfàn, yìbiān tīng yīnyuè. Wǒmen jiā zhēn rènao!}\\
« Le soir, mon petit frère mange des fruits tout en regardant la télé ; maman cuisine tout en écoutant de la musique. Comme c'est animé chez nous ! »
\end{corrige}

\begin{exercice}[Exercice 4 — Production écrite guidée]
Rédigez un court texte de quatre phrases intitulé \zh{在我家} (Zài wǒ jiā, « Chez moi ») : décrivez ce que font les membres de votre famille le soir, en utilisant au moins trois fois la structure \zh{一边……一边……} et une ponctuation chinoise correcte (virgules, deux-points, point final).
\end{exercice}

\begin{corrige}[Exercice 4 — Modèle de rédaction]
\zh{晚上，我们一家人都很忙。爸爸一边看报纸，一边喝茶。妈妈一边做饭，一边打电话。我一边做作业，一边听音乐。}\\
\emph{Wǎnshang, wǒmen yìjiārén dōu hěn máng. Bàba yìbiān kàn bàozhǐ, yìbiān hēchá. Māma yìbiān zuòfàn, yìbiān dǎ diànhuà. Wǒ yìbiān zuò zuòyè, yìbiān tīng yīnyuè.}\\
« Le soir, toute la famille est occupée. Papa lit le journal tout en buvant son thé. Maman cuisine tout en téléphonant. Moi, je fais mes devoirs tout en écoutant de la musique. »
\end{corrige}

\begin{experience}[Jeu de rôle : le journaliste multitâche]
Par groupes de trois : un élève joue un journaliste de CRTV ou de CCTV qui décrit en direct ce qu'il voit dans la cour du lycée ; les deux autres miment deux actions simultanées (manger + parler, marcher + téléphoner, lire + écrire...). Le journaliste doit produire des phrases correctes avec \zh{一边……一边……}, par exemple : \zh{那个学生一边走，一边打电话。} (Nàge xuésheng yìbiān zǒu, yìbiān dǎ diànhuà. « Cet élève téléphone tout en marchant. ») La classe vérifie : sujet unique, deux \zh{一边}, verbes d'action.
\end{experience}

\begin{savaistu}[Cameroun / Chine : la ponctuation, une histoire moderne]
La ponctuation chinoise est récente : avant le XX\textsuperscript{e} siècle, les textes chinois s'écrivaient presque sans signes, et le lecteur devait découper lui-même les phrases ! Les signes modernes (\zh{。，？！《》}) ont été officialisés au XX\textsuperscript{e} siècle, sous l'influence des langues européennes, puis adaptés à l'écriture chinoise. Au Cameroun, où le français et l'anglais imposent leurs propres règles de ponctuation, apprendre celles du chinois, c'est découvrir qu'une langue peut inventer sa manière de « respirer » à l'écrit.
\end{savaistu}

\begin{aretenir}[Bilan de la leçon]
\begin{itemize}
\item \textbf{Structure} : Sujet + 一边 + verbe 1, 一边 + verbe 2 = deux actions simultanées, un seul sujet.
\item \textbf{Limites} : même sujet obligatoire ; actions compatibles ; pas d'adjectif (utiliser \zh{又……又……}) ; variante écrite \zh{一面……一面……}.
\item \textbf{Pièges} : ne jamais oublier le second \zh{一边} ; ne jamais répéter le sujet.
\item \textbf{Ponctuation chinoise} : point rond \zh{。}, virgule \zh{，}, virgule d'énumération \zh{、}, deux-points \zh{：}, point-virgule \zh{；}, \zh{？}, \zh{！}, guillemets \zh{“”}, titres \zh{《》}, suspension \zh{……}. Aucune espace avant le signe.
\end{itemize}
\end{aretenir}

\coursec{Leçon 40 — Grammaire : 一边……一边…… ; La ponctuation en chinois}

\begin{textebox}[Situation de communication : Au cybercafé de Yaoundé]
\zh{—— 你看，那个学生一边上网，一边查资料。} \\
\emph{Nǐ kàn, nàgè xuésheng yībiān shàngwǎng, yībiān chá zīliào.} \\
« Regarde, cet étudiant navigue sur Internet tout en cherchant des informations. » \\[4pt]
\zh{—— 是的，我们一边听课，一边记笔记。} \\
\emph{Shì de, wǒmen yībiān tīngkè, yībiān jì bǐjì.} \\
« Oui, nous, on écoute le cours tout en prenant des notes. »
\end{textebox}

\courssub{Observation et découverte de la structure}

\begin{experience}[Repérage guidé]
Lisez le dialogue ci-dessus. Identifiez : 
\begin{enumerate}
\item Le sujet commun aux deux actions.
\item Les deux verbes d'action.
\item Le mot-outil qui se répète devant chaque verbe.
\item Le signe de ponctuation qui sépare les deux parties.
\end{enumerate}
\tcblower
\textbf{Réponses.} 
\begin{enumerate}
\item Sujet : \zh{那个学生} / \zh{我们}.
\item Verbes : \zh{上网} / \zh{查资料} ; \zh{听课} / \zh{记笔记}.
\item Mot-outil : \zh{一边} (répété deux fois).
\item Ponctuation : la virgule chinoise \zh{，}.
\end{enumerate}
\end{experience}

\begin{definition}[Notion : La corrélation \zh{一边……一边……}]
La structure \zh{一边……一边……} (souvent prononcé \emph{yíbiān……yíbiān} par assimilation de ton) est une \textbf{corrélation} qui exprime la \textbf{simultanéité de deux actions volontaires} accomplies par le même sujet. Elle correspond au français « tout en + infinitif », « en + gérondif » ou « pendant que + indicatif ».
\end{definition}

\courssub{Schéma et construction de la structure}

\begin{definition}[Schéma canonique]
La structure suit le schéma immuable suivant :\\[4pt]
\textbf{Sujet + 一边 + Verbe 1 (+ Complément 1)，一边 + Verbe 2 (+ Complément 2)。}\\[4pt]
\begin{itemize}
\item Le \textbf{sujet} n'apparaît qu'\textbf{une seule fois}, placé \textbf{avant} le premier \zh{一边}.
\item Les deux membres sont séparés par la \textbf{virgule chinoise} \zh{，}.
\item La phrase se termine par le \textbf{point chinois} \zh{。}.
\end{itemize}
\end{definition}

\begin{popfigure}
\begin{tikzpicture}[
  box/.style={draw=popblue, fill=popblueL, rounded corners=2pt, align=center, minimum height=0.9cm, font=\small, minimum width=2.2cm},
  arrow/.style={-{Stealth[length=2mm]}, thick, popdark}
]
\node[box, align=center] (s) at (0,0){\textbf{Sujet}\\ \zh{我们}};
\node[box, align=center] (a) at (3.2,0){\zh{一边} + V1 (+ C1)\\ \zh{一边上网}};
\node[box, fill=poporangeL, draw=poporange, align=center] (v) at (6.0,0){virgule\\ \zh{，}};
\node[box, align=center] (b) at (8.8,0){\zh{一边} + V2 (+ C2)\\ \zh{一边查资料}};
\node[box, fill=poporangeL, draw=poporange, align=center] (p) at (11.6,0){point\\ \zh{。}};
\draw[arrow] (s) -- (a);
\draw[arrow] (a) -- (v);
\draw[arrow] (v) -- (b);
\draw[arrow] (b) -- (p);
\end{tikzpicture}
\legende{Schéma chaîné : Sujet → 一边 + V1 → ，→ 一边 + V2 → 。}
\end{popfigure}

\begin{propriete}[Un seul sujet, actions volontaires]
La structure exige \textbf{un seul sujet} (le même agent pour les deux actions). Les verbes doivent exprimer des \textbf{actions actives, volontaires, contrôlables} (marcher, manger, lire, écouter, écrire…). On \textbf{ne peut pas} utiliser cette structure avec des verbes d'état (\zh{是}, \zh{在}, \zh{有}, \zh{喜欢}, \zh{知道}…) ni des verbes involontaires (\zh{生病}, \zh{睡着}, \zh{忘记}…).
\end{propriete}

\begin{exemplebox}[Exemples analysés (contexte camerounais et chinois)]
\begin{itemize}
\item \zh{我们一边上网，一边查资料。}\\ \emph{Wǒmen yībiān shàngwǎng, yībiān chá zīliào.}\\ « Nous naviguons sur Internet tout en cherchant des infos. »
\item \zh{他一边听音乐，一边写作业。}\\ \emph{Tā yībiān tīng yīnyuè, yībiān xiě zuòyè.}\\ « Il écoute de la musique tout en faisant ses devoirs. »
\item \zh{妈妈一边做饭，一边打电话。}\\ \emph{Māma yībiān zuòfàn, yībiān dǎ diànhuà.}\\ « Maman cuisine tout en téléphonant. »
\item \zh{学生们一边看足球赛，一边讨论。}\\ \emph{Xuéshengmen yībiān kàn zúqiúsài, yībiān tǎolùn.}\\ « Les élèves regardent le match de foot tout en discutant. »
\item \zh{我在喀麦隆一边学中文，一边教法语。}\\ \emph{Wǒ zài Kāmàilóng yībiān xué Zhōngwén, yībiān jiāo Fǎyǔ.}\\ « Au Cameroun, j'apprends le chinois tout en enseignant le français. »
\end{itemize}
\end{exemplebox}

\courssub{La place des compléments (temps, lieu, objet)}

\begin{propriete}[Règle de placement]
Les compléments (lieu, temps, objet) se placent \textbf{à l'intérieur de chaque membre}, selon l'ordre canonique chinois : \textbf{Sujet + Temps + Lieu + Verbe + Objet}.
\end{propriete}

\begin{exemplebox}[Compléments de lieu et de temps]
\zh{他一边在房间里听音乐，一边写作业。}\\ \emph{Tā yībiān zài fángjiān lǐ tīng yīnyuè, yībiān xiě zuòyè.}\\ « Tout en écoutant de la musique \textbf{dans sa chambre}, il fait ses devoirs. » (Lieu avant V1)\\[4pt]
\zh{我们一边在教室里上课，一边记笔记。}\\ \emph{Wǒmen yībiān zài jiàoshì lǐ shàngkè, yībiān jì bǐjì.}\\ « Tout en suivant le cours \textbf{dans la salle}, nous prenons des notes. »\\[4pt]
\zh{弟弟一边吃早饭，一边看手机。}\\ \emph{Dìdi yībiān chī zǎofàn, yībiān kàn shǒujī.}\\ « Le petit frère mange son petit-déjeuner tout en regardant son téléphone. » (Objet après V1 et V2)
\end{exemplebox}

\courssub{Deux propriétés essentielles : virgule obligatoire et ordre interchangeable}

\begin{propriete}[La virgule \zh{，} est obligatoire]
Les deux membres de la corrélation sont \textbf{toujours} séparés par la virgule chinoise \zh{，}. C'est une règle syntaxique et prosodique : elle marque la pause entre les deux actions simultanées.
\end{propriete}

\begin{propriete}[L'ordre des actions est interchangeable]
Les deux actions ayant le même poids (simultanéité stricte), on peut inverser les membres \textbf{sans changer le sens global} :\\[4pt]
\zh{她一边唱歌，一边跳舞。} $\Leftrightarrow$ \zh{她一边跳舞，一边唱歌。}\\ \emph{Tā yībiān chànggē, yībiān tiàowǔ.}\\ « Elle chante tout en dansant. » $\Leftrightarrow$ « Elle danse tout en chantant. »
\end{propriete}

\begin{methode}[Construire une phrase avec 一边……一边…… (5 étapes)]
\begin{enumerate}
\item \textbf{Identifier le sujet unique} : \zh{我}.
\item \textbf{Écrire \zh{一边} + Action 1 (+ compléments)} : \zh{我一边喝茶}.
\item \textbf{Insérer la virgule chinoise} : \zh{，}.
\item \textbf{Écrire \zh{一边} + Action 2 (+ compléments)} : \zh{一边看报纸}.
\item \textbf{Terminer par le point chinois} : \zh{。} $\rightarrow$ \zh{我一边喝茶，一边看报纸。} \emph{Wǒ yībiān hē chá, yībiān kàn bàozhǐ.}
\end{enumerate}
\end{methode}

\courssub{Variantes de la structure (niveau lecture / culture)}

\begin{definition}[Variante à deux sujets (rare, compréhension écrite)]
Schéma : \textbf{Sujet 1 + 一边 + V1，Sujet 2 + 一边 + V2。}\\[4pt]
\zh{老师一边讲，学生一边记。}\\ \emph{Lǎoshī yībiān jiǎng, xuésheng yībiān jì.}\\ « Le prof explique, les élèves prennent des notes (en même temps). »\\[4pt]
\emph{À l'écrit (examen), utilisez toujours la forme à sujet unique.}
\end{definition}

\begin{definition}[Variante abrégée \zh{边……边……} (registre soutenu, presse, littérature)]
\zh{边} est la forme brève de \zh{一边}. Schéma : \textbf{Sujet + 边 + V1，边 + V2。} (La virgule est souvent omise si V1 et V2 sont monosyllabiques).\\[4pt]
\zh{他边走边说。}\\ \emph{Tā biān zǒu biān shuō.}\\ « Il parle en marchant. »\\[4pt]
\zh{代表团边参观边交流。}\\ \emph{Dàibiǎotuán biān cānguān biān jiāoliú.}\\ « La délégation visite tout en échangeant. »
\end{definition}

\begin{center}
\begin{tabularx}{\linewidth}{|X|X|X|}
\hline
\rowcolor{popblueL} \textbf{Forme} & \textbf{Exemple (caractères + pinyin)} & \textbf{Traduction} \\
\hline
Pleine \zh{一边……一边……} (courant, oral/écrit) & \zh{他一边走路，一边说话。} \emph{Tā yībiān zǒulù, yībiān shuōhuà.} & Il marche tout en parlant. \\
\hline
Abrégée \zh{边……边……} (soutenu, presse) & \zh{他边走边说。} \emph{Tā biān zǒu biān shuō.} & Il parle en marchant. \\
\hline
Deux sujets (lecture) & \zh{你一边问，我一边答。} \emph{Nǐ yībiān wèn, wǒ yībiān dá.} & Tu questions, je réponds (simultanément). \\
\hline
\end{tabularx}
\end{center}

\begin{savaistu}[L'expression figée \zh{边走边看}]
\zh{边走边看} \emph{biān zǒu biān kàn} (litt. « marcher et regarder ») signifie au figuré : « avancer à tâtons », « décider au fur et à mesure », « voir comment les choses évoluent ». Exemple : \zh{这个项目比较复杂，我们只能边走边看。} \emph{Zhège xiàngmù bǐjiào fùzá, wǒmen zhǐ néng biān zǒu biān kàn.} « Ce projet est complexe, on ne peut qu'avancer à tâtons. »
\end{savaistu}

\courssub{Comparaison avec le français et pièges des francophones}

Le français marque la simultanéité une seule fois (\emph{tout en, en, pendant que}). Le chinois la marque \textbf{deux fois} (corrélation). C'est la source des trois erreurs majeures ci-dessous.

\begin{center}
\begin{tabularx}{\linewidth}{|X|X|X|}
\hline
\rowcolor{popblueL} \textbf{Français} & \textbf{Chinois correct} & \textbf{Erreur fréquente (à bannir)} \\
\hline
Il écoute de la musique en faisant ses devoirs. & \zh{他一边听音乐，一边写作业。} & \zh{他一边听音乐，写作业。} (2\textsuperscript{e} \zh{一边} oublié) \\
\hline
Nous mangeons en regardant la télé. & \zh{我们一边吃饭，一边看电视。} & \zh{我们一边吃饭，一边我们看电视。} (Sujet répété) \\
\hline
Elle chante en dansant. & \zh{她一边唱歌，一边跳舞。} & \zh{她一边唱歌一边跳舞。} (Virgule oubliée) \\
\hline
Il est étudiant et il travaille. (États) & \zh{他又是学生，又工作。} (\zh{又……又……}) & \zh{他一边是学生，一边工作。} (Verbes d'état interdits) \\
\hline
\end{tabularx}
\end{center}

\begin{attention}[Les trois erreurs fatales des francophones]
\begin{enumerate}
\item \textbf{Oublier le second \zh{一边}} : Calque du français « en + V ». \emph{Règle :} Corrélation = double marquage obligatoire.
\item \textbf{Répéter le sujet} devant le second \zh{一边} : Calque de « il fait X et il fait Y ». \emph{Règle :} Sujet unique en tête de phrase.
\item \textbf{Utiliser avec des verbes d'état / involontaires} : \zh{一边生病，一边吃药} (faux), \zh{一边喜欢，一边讨厌} (faux). \emph{Règle :} Actions volontaires uniquement. Pour états/qualités simultanés $\rightarrow$ \zh{又……又……} (leçon ultérieure).
\end{enumerate}
\end{attention}

\begin{exoresolu}[Construction guidée]
\textbf{Énoncé.} Construisez une phrase avec \zh{一边……一边……} : Sujet \zh{妹妹} (petite sœur) ; Action 1 : \zh{唱歌} ; Action 2 : \zh{跳舞}.
\tcblower
\textbf{Solution détaillée.}\\
1. Sujet : \zh{妹妹}.\\
2. \zh{一边} + V1 : \zh{妹妹一边唱歌}.\\
3. Virgule : \zh{，}.\\
4. \zh{一边} + V2 : \zh{一边跳舞}.\\
5. Point : \zh{。}\\[4pt]
\zh{妹妹一边唱歌，一边跳舞。} \emph{Mèimei yībiān chànggē, yībiān tiàowǔ.}\\
« Ma petite sœur chante tout en dansant. » (Ordre inversé possible : \zh{妹妹一边跳舞，一边唱歌。})
\end{exoresolu}

\courssub{La ponctuation en chinois (标点符号)}

\begin{definition}[Principes de base]
La ponctuation chinoise utilise des \textbf{signes pleins (pleine chasse)} qui occupent la place d'un caractère. \textbf{Pas d'espace} avant ni après les signes (sauf cas rares de fin de ligne). Les noms propres ne prennent pas de majuscule visuelle (pas de capitale), mais on peut souligner les noms propres en manuscrit (trait de soulignement).
\end{definition}

\begin{propriete}[Tableau des signes principaux (niveau HSK 1-4)]
\begin{center}
\begin{tabularx}{\linewidth}{|c|X|X|}
\hline
\rowcolor{popblueL} \textbf{Signe} & \textbf{Nom / Fonction} & \textbf{Exemple (Caractères + Pinyin)} \\
\hline
\zh{。} & Point final (句号 \emph{jùhào}) & \zh{我爱喀麦隆。} \emph{Wǒ ài Kāmàilóng.} \\
\hline
\zh{，} & Virgule (逗号 \emph{dòuhào}) — séparation clauses / items & \zh{他一边吃，一边看。} \emph{Tā yībiān chī, yībiān kàn.} \\
\hline
\zh{、} & Virgule d'énumération (顿号 \emph{dùnhào}) — liste de noms/verbes & \zh{我喜欢吃米饭、面条、水饺。} \emph{Wǒ xǐhuan chī mǐfàn、miàntiáo、shuǐjiǎo.} \\
\hline
\zh{？} & Point d'interrogation (问号 \emph{wènhào}) & \zh{你叫什么名字？} \emph{Nǐ jiào shénme míngzi?} \\
\hline
\zh{！} & Point d'exclamation (感叹号 \emph{gǎntànhào}) & \zh{太美了！} \emph{Tài měi le!} \\
\hline
\zh{：} & Deux-points (冒号 \emph{màohào}) — annonce, citation & \zh{他说：} \zh{「我来了。」} \emph{Tā shuō: « Wǒ lái le. »} \\
\hline
\zh{；} & Point-virgule (分号 \emph{fēnhào}) — clauses indépendantes longues & \zh{天很热；我们去游泳吧。} \emph{Tiān hěn rè; wǒmen qù yóuyǒng ba.} \\
\hline
\zh{“ ”} & Guillemets doubles (引号 \emph{yǐnhào}) — citation, dialogue & \zh{他问：「你好吗？」} \emph{Tā wèn: « Nǐ hǎo ma? »} \\
\hline
\zh{‘ ’} & Guillemets simples (单引号) — citation dans citation & \zh{他说：『我听见有人喊“救命”』} \\
\hline
\zh{《 》} & Guillemets de titres d'ouvrages (书名号 \emph{shūmínghào}) & \zh{我读过《红楼梦》。} \emph{Wǒ dúguò « Hónglóumèng ».} \\
\hline
\zh{……} & Points de suspension (省略号 \emph{shěnglüèhào}) — 6 points, occupent 2 cases & \zh{这个故事很长……} \emph{Zhège gùshì hěn cháng……} \\
\hline
\zh{·} & Point médian (间隔号 \emph{jiàngéhào}) — noms étrangers, dates & \zh{雅温得} \emph{Yǎwēndé} (Yaoundé), \zh{2024·10·01} \\
\hline
\end{tabularx}
\end{center}
\end{propriete}

\begin{attention}[Différences majeures avec le français]
\begin{enumerate}
\item \textbf{Pas d'espace} autour des signes : \zh{你好！} (pas \zh{你好 ！}).
\item \textbf{Virgules distinctes} : \zh{，} pour clauses, \zh{、} pour énumérations simples (noms/verbes). Ne les confondez pas.
\item \textbf{Guillemets} : \zh{“ ”} et \zh{‘ ’} (pleine chasse), pas \zh{« »}.
\item \textbf{Titres} : \zh{《 》} obligatoires pour livres, films, journaux, navires. Pas d'italique.
\item \textbf{Point médian \zh{·}} : Sépare prénom/nom étrangers (\zh{马克·吐温} Mǎkè·Tǔwēn) ou éléments de date.
\end{enumerate}
\end{attention}

\courssub{Entraînement intégré : Structure \zh{一边……一边……} + Ponctuation}

\begin{exercice}[Transformation et ponctuation]
Transformez les paires de phrases en une seule phrase avec \zh{一边……一边……}. Ajoutez la ponctuation correcte (virgule \zh{，}, point \zh{。}). Traduisez.
\begin{enumerate}
\item \zh{哥哥听广播。} + \zh{哥哥做运动。}
\item \zh{同学们在教室里上课。} + \zh{同学们记笔记。}
\item \zh{爸爸开车。} + \zh{爸爸打电话。} (Ajoutez un commentaire de sécurité avec \zh{！})
\item \zh{我喝茶。} + \zh{我看报纸。} + \zh{我听音乐。} (Trois actions : utilisez \zh{一边……一边……一边……} possible ? Non, structure binaire. Choisissez deux actions.)
\end{enumerate}
\end{exercice}

\begin{exercice}[Correction de ponctuation]
Corrigez la ponctuation (signes manquants, mauvais signes, espaces) dans les phrases suivantes :
\begin{enumerate}
\item \zh{他一边跑步 一边听歌}
\item \zh{我们喜欢吃饭 、 睡觉 、 打球 。}
\item \zh{他说 「 我 是 中国 人 」}
\item \zh{红楼梦 是 名著}
\item \zh{雅温得 是 喀麦隆 的 首都}
\end{enumerate}
\end{exercice}

\begin{corrige}[Exercices intégrés]
\textbf{Exercice 1 (Transformation) :}
\begin{enumerate}
\item \zh{哥哥一边听广播，一边做运动。} \emph{Gēge yībiān tīng guǎngbō, yībiān zuò yùndòng.} « Grand frère écoute la radio tout en faisant du sport. »
\item \zh{同学们一边在教室里上课，一边记笔记。} \emph{Tóngxuémen yībiān zài jiàoshì lǐ shàngkè, yībiān jì bǐjì.} « Les camarades suivent le cours en salle tout en prenant des notes. »
\item \zh{爸爸一边开车，一边打电话！} \emph{Bàba yībiān kāichē, yībiān dǎ diànhuà!} « Papa conduit tout en téléphonant ! » \zh{这非常危险！} \emph{Zhè fēicháng wēixiǎn!} « C'est extrêmement dangereux ! »
\item \zh{我一边喝茶，一边看报纸。} (Choix de deux actions parmi trois). \emph{Wǒ yībiān hē chá, yībiān kàn bàozhǐ.} « Je bois du thé tout en lisant le journal. »
\end{enumerate}

\textbf{Exercice 2 (Correction ponctuation) :}
\begin{enumerate}
\item \zh{他一边跑步，一边听歌。} (Ajout virgule \zh{，}, point \zh{。})
\item \zh{我们喜欢吃饭、睡觉、打球。} (Remplacement virgules françaises par \zh{、}, suppression espaces, point \zh{。})
\item \zh{他说：「我是中国人。」} (Guillemets \zh{「 」}, deux-points \zh{：}, point \zh{。} à l'intérieur).
\item \zh{《红楼梦》是名著。} (Guillemets titres \zh{《 》}, point \zh{。})
\item \zh{雅温得·是喀麦隆的首都。} (Point médian \zh{·} pour nom étranger \zh{雅温得} ? Non, \zh{雅温得} est un mot unique transcrit. Correction : \zh{雅温得是喀麦隆的首都。} Pas de point médian nécessaire ici sauf si on sépare \zh{雅·温·得} ce qui ne se fait pas. Phrase simple : Sujet + Verbe + Complément. Point final \zh{。})
\end{enumerate}
\end{corrige}

\courssub{Vocabulaire clé de la leçon (HSK 1-3)}

\begin{center}
\begin{tabularx}{\linewidth}{|X|X|X|X|}
\hline
\rowcolor{popblueL} \textbf{Caractères} & \textbf{Pinyin} & \textbf{Nature} & \textbf{Traduction} \\
\hline
一边 & yībiān & nom / mot-outil & côté ; (structure) tout en… \\
\hline
上网 & shàngwǎng & verbe & naviguer sur Internet, aller en ligne \\
\hline
查资料 & chá zīliào & verbe + objet & chercher des informations / des documents \\
\hline
听音乐 & tīng yīnyuè & verbe + objet & écouter de la musique \\
\hline
写作业 & xiě zuòyè & verbe + objet & faire ses devoirs \\
\hline
做饭 & zuòfàn & verbe & cuisiner, faire la cuisine \\
\hline
打电话 & dǎ diànhuà & verbe + objet & téléphoner, passer un coup de fil \\
\hline
看比赛 & kàn bǐsài & verbe + objet & regarder un match \\
\hline
讨论 & tǎolùn & verbe & discuter, débattre \\
\hline
喝茶 & hē chá & verbe + objet & boire du thé \\
\hline
看报纸 & kàn bàozhǐ & verbe + objet & lire le journal \\
\hline
吃早饭 & chī zǎofàn & verbe + objet | verbe & prendre le petit-déjeuner \\
\hline
看手机 & kàn shǒujī & verbe + objet & regarder son téléphone \\
\hline
听广播 & tīng guǎngbō & verbe + objet & écouter la radio \\
\hline
做运动 & zuò yùndòng & verbe + objet & faire du sport, s'exercer \\
\hline
记笔记 & jì bǐjì & verbe + objet & prendre des notes \\
\hline
开车 & kāichē & verbe | verbe + objet & conduire (une voiture) \\
\hline
边走边看 & biān zǒu biān kàn & expression & avancer à tâtons ; décider au fur et à mesure \\
\hline
标点符号 & biāodiǎn fúhào & nom & signes de ponctuation \\
\hline
逗号 & dòuhào & nom & virgule (clauses) \zh{，} \\
\hline
顿号 & dùnhào & nom & virgule d'énumération \zh{、} \\
\hline
书名号 & shūmínghào & nom & guillemets de titre \zh{《 》} \\
\hline
\end{tabularx}
\end{center}

\courssub{Écriture des caractères : ordre des traits et radicaux}

\begin{definition}[Caractères à maîtriser (reconnaissance + écriture)]
\begin{itemize}
\item \zh{边} \emph{biān} (côté, bord) — \textbf{Radical} : \zh{辶} (chuò, « marche »). \textbf{Traits} (5 + 3 = 8) : \zh{自} (haut→bas, gauche→droite : 丨, (trait brisé), 一, 一, 丨) + \zh{辶} (点, 横折折撇, 捺). \emph{Astuce :} \zh{自} (soi-même) à côté de \zh{辶} (chemin) = le bord du chemin.
\item \zh{听} \emph{tīng} (écouter) — \textbf{Radical} : \zh{立} (lì, « debout »). \textbf{Traits} : \zh{立} (5 traits) + \zh{十} (2) + \zh{一} (1) + \zh{心} (4, cœur en bas) = 12 traits. \emph{Logique :} Oreille \zh{耳} (dans \zh{圣} version trad.) + Œil \zh{目} + Un \zh{一} + Cœur \zh{心} = Écouter de tout son cœur. (Simplifié : \zh{立} + \zh{十} + \zh{一} + \zh{心}).
\item \zh{看} \emph{kàn} (regarder) — \textbf{Radical} : \zh{目} (mù, « œil »). \textbf{Traits} : \zh{目} (5) + \zh{手} (4, main simplifiée en haut) = 9 traits. \emph{Logique :} Main sur les yeux pour mieux voir.
\item \zh{写} \emph{xiě} (écrire) — \textbf{Radical} : \zh{冖} (mì, « couvrir »). \textbf{Traits} : \zh{冖} (2) + \zh{与} (3, donner) = 5 traits. \emph{Logique :} Couvrir + Donner = Transmettre par écrit.
\item \zh{吃} \emph{chī} (manger) — \textbf{Radical} : \zh{口} (kǒu, « bouche »). \textbf{Traits} : \zh{口} (3) + \zh{乞} (qǐ, mendiant, 3 traits) = 6 traits. \emph{Logique :} Bouche + Mendiant = Manger (quémander de la nourriture).
\item \zh{做} \emph{zuò} (faire) — \textbf{Radical} : \zh{亻} (rén, « homme »). \textbf{Traits} : \zh{亻} (2) + \zh{故} (gù, ancien, 9 traits) = 11 traits. \emph{Note :} \zh{故} = \zh{古} (ancien) + \zh{攵} (pū, action/taper).
\end{itemize}
\end{definition}

\begin{aretenir}[Synthèse de la leçon 40]
\begin{itemize}
\item \textbf{Structure \zh{一边……一边……}} : Sujet unique + \zh{一边} + V1 (+ C1) \zh{，} \zh{一边} + V2 (+ C2) \zh{。}
\item \textbf{Conditions} : Actions volontaires, simultanées, même sujet. Ordre interchangeable.
\item \textbf{Pièges} : Oubli 2\textsuperscript{e} \zh{一边}, répétition sujet, verbes d'état, oubli virgule.
\item \textbf{Variantes} : \zh{边……边……} (soutenu), deux sujets (lecture).
\item \textbf{Ponctuation} : Signes pleins, pas d'espaces. Distinction \zh{，} (clauses) / \zh{、} (listes). Guillemets \zh{“ ”} / \zh{《 》}. Point médian \zh{·} (noms étrangers).
\end{itemize}
\end{aretenir}

\coursec{Emplois, limites et exceptions}

Dans la section précédente, nous avons établi le schéma de la structure : \zh{一边 + Verbe$_1$，一边 + Verbe$_2$}, avec un sujet unique qui accomplit deux actions en même temps. Nous savons donc \emph{construire} la phrase. Mais savoir construire ne suffit pas : il faut savoir \emph{quand} l'employer, et surtout quand \emph{ne pas} l'employer. C'est l'objet de cette section : délimiter précisément le territoire d'emploi de \zh{一边……一边……}, ses limites et ses exceptions, afin d'éviter les contresens que les francophones commettent très souvent sous l'influence du gérondif français « en + participe ».

\courssub{Observation : quand la structure est-elle possible ?}

\begin{textebox}[Dialogue : au cybercafé de Yaoundé]
\zh{A：你昨天一边吃饭，一边看足球比赛，对吗？}\\
\emph{A : Nǐ zuótiān yìbiān chīfàn, yìbiān kàn zúqiú bǐsài, duì ma ?}\\
A : Hier, tu regardais le match de football en mangeant, c'est ça ?\\
\zh{B：对！我一边吃，一边给哥哥打电话。}\\
\emph{B : Duì ! Wǒ yìbiān chī, yìbiān gěi gēge dǎ diànhuà.}\\
B : Oui ! Je mangeais tout en téléphonant à mon grand frère.\\
\zh{A：你不能一边睡觉，一边看比赛吧？}\\
\emph{A : Nǐ bù néng yìbiān shuìjiào, yìbiān kàn bǐsài ba ?}\\
A : Tu ne peux quand même pas dormir et regarder le match en même temps ?\\
\zh{B：当然不能！那是笑话。}\\
\emph{B : Dāngrán bù néng ! Nà shì xiàohua.}\\
B : Bien sûr que non ! Ce serait une blague.
\end{textebox}

Observons. Les deux premières phrases avec \zh{一边……一边……} sont naturelles : manger + regarder, manger + téléphoner sont des couples d'actions \cle{compatibles}. En revanche, dormir + regarder un match est physiquement impossible : le locuteur B le signale lui-même avec humour. Remarquons aussi la ponctuation : les deux membres de la structure sont séparés par la \cle{virgule chinoise} \zh{，}, jamais par un point. Cette observation nous donne la première grande condition d'emploi.

\courssub{Emploi principal : deux actions simultanées, durables et compatibles}

\begin{propriete}[Condition fondamentale d'emploi]
\zh{一边……一边……} exprime que \textbf{le même sujet} accomplit \textbf{deux actions simultanées}, qui remplissent trois conditions :
\begin{enumerate}
\item \textbf{Simultanéité réelle} : les deux actions se déroulent vraiment pendant le même laps de temps ;
\item \textbf{Durée} : ce sont des actions qui peuvent durer (duratives), pas des événements instantanés ;
\item \textbf{Compatibilité physique et logique} : il doit être possible de faire les deux en même temps.
\end{enumerate}
\end{propriete}

\begin{exemplebox}[Couples d'actions compatibles]
\zh{他一边走路，一边打电话。}\\
\emph{Tā yìbiān zǒulù, yìbiān dǎ diànhuà.}\\
« Il téléphone en marchant. » (marcher et téléphoner : compatibles)\\[4pt]
\zh{妈妈一边做饭，一边听音乐。}\\
\emph{Māma yìbiān zuòfàn, yìbiān tīng yīnyuè.}\\
« Maman écoute de la musique en cuisinant. »\\[4pt]
\zh{学生们一边听老师讲课，一边写汉字。}\\
\emph{Xuéshengmen yìbiān tīng lǎoshī jiǎngkè, yìbiān xiě Hànzì.}\\
« Les élèves écrivent les caractères tout en écoutant le cours du professeur. »\\[4pt]
\zh{我们一边吃恩多莱，一边聊天。}\\
\emph{Wǒmen yìbiān chī ndolé, yìbiān liáotiān.}\\
« Nous bavardons en mangeant le ndolé. »
\end{exemplebox}

\begin{exemplebox}[Correct et absurde : le test de compatibilité]
\textbf{正确 (correct) :} \zh{他一边走路，一边打电话。} \emph{Tā yìbiān zǒulù, yìbiān dǎ diànhuà.} « Il téléphone en marchant. » — on peut réellement faire les deux.\\[4pt]
\textbf{错误 (incorrect) :} \zh{*他一边睡觉，一边跑步。} \emph{*Tā yìbiān shuìjiào, yìbiān pǎobù.} — absurde : personne ne peut dormir et courir en même temps. La grammaire est correcte, mais le sens est impossible : la phrase est rejetée.
\end{exemplebox}

\begin{methode}[Le test de compatibilité en trois questions]
Avant d'écrire une phrase avec \zh{一边……一边……}, pose-toi trois questions :
\begin{enumerate}
\item Les deux actions ont-elles lieu \textbf{au même moment} ? (Sinon : utiliser \zh{先……然后……}.)
\item Une personne réelle peut-elle \textbf{physiquement} faire les deux ensemble ? (Sinon : phrase absurde.)
\item Les deux verbes sont-ils des \textbf{actions actives} ? (Sinon : voir la limite suivante.)
\end{enumerate}
Si la réponse est « oui » aux trois questions, la structure est correcte.
\end{methode}

\courssub{Limite 1 : les verbes d'état sont exclus}

\begin{propriete}[Actions volontaires et durables uniquement]
Les deux membres de \zh{一边……一边……} doivent exprimer des \textbf{actions volontaires et durables}. Les \textbf{verbes d'état} sont exclus : \zh{是} (être), \zh{有} (avoir), \zh{像} (ressembler), \zh{在} au sens d'« exister / se trouver », ainsi que les verbes purement psychologiques duratifs comme \zh{知道} (savoir), \zh{爱} (aimer) employés seuls.
\end{propriete}

Pourquoi ? Parce que \zh{一边……一边……} décrit un \emph{partage de l'attention et de l'activité} entre deux tâches. Un état ne « s'accomplit » pas : on ne fait pas activement « être étudiant » ou « avoir un frère ».

\begin{exemplebox}[Verbes d'état : phrases à éviter]
\zh{*他一边是学生，一边是老师。} — incorrect avec le sens « il est à la fois étudiant et professeur ». On dira plutôt : \zh{他既是学生，又是老师。} \emph{Tā jì shì xuésheng, yòu shì lǎoshī.} « Il est à la fois étudiant et professeur. »\\[4pt]
\zh{*她一边有电脑，一边有手机。} — incorrect. On dira : \zh{她既有电脑，又有手机。} \emph{Tā jì yǒu diànnǎo, yòu yǒu shǒujī.} « Elle a à la fois un ordinateur et un téléphone. »\\[4pt]
En revanche, avec des verbes d'action, tout redevient possible : \zh{她一边用电脑，一边打电话。} \emph{Tā yìbiān yòng diànnǎo, yìbiān dǎ diànhuà.} « Elle téléphone tout en utilisant l'ordinateur. »
\end{exemplebox}

\courssub{Limite 2 : jamais pour deux actions successives}

C'est la limite la plus importante pour vous, francophones. En français, le gérondif « en + participe » sert parfois à enchaîner deux événements qui se suivent : « Il est rentré \emph{en} mangeant un beignet » peut vouloir dire « il est rentré, \emph{puis} il a mangé ». Le chinois \zh{一边……一边……} \textbf{refuse absolument} cet usage : il impose la simultanéité stricte.

\begin{attention}[Le piège du gérondif français « en »]
Ne traduis jamais mécaniquement « en + participe » par \zh{一边……一边……} quand les actions sont \textbf{successives}.\\[4pt]
\textbf{错误 :} \zh{*他一边回家，一边吃饭。} pour dire « il rentre chez lui, puis il mange » — FAUX, car rentrer et manger se succèdent, ils ne se déroulent pas en même temps.\\[4pt]
\textbf{正确 :} \zh{他先回家，然后吃饭。} \emph{Tā xiān huí jiā, ránhòu chīfàn.} « Il rentre d'abord chez lui, puis il mange. »\\[4pt]
Règle pratique : si tu peux remplacer « en » par « puis » en français, alors \zh{一边……一边……} est interdit ; utilise \zh{先……然后……}.
\end{attention}

\begin{center}
\begin{tabularx}{\linewidth}{|X|X|X|}
\hline
\rowcolor{popblueL} Situation & Structure chinoise & Exemple \\
\hline
Deux actions simultanées & \zh{一边……一边……} & \zh{他一边吃饭，一边看电视。} Il regarde la télé en mangeant. \\
\hline
Deux actions successives & \zh{先……然后……} & \zh{他先吃饭，然后看电视。} Il mange d'abord, puis il regarde la télé. \\
\hline
Un moment + un événement & \zh{……的时候} & \zh{吃饭的时候，电话响了。} Pendant le repas, le téléphone a sonné. \\
\hline
\end{tabularx}
\end{center}

\courssub{Limite 3 : la négation}

La négation de \zh{一边……一边……} est \textbf{rare} et maladroite. On ne dit presque jamais \zh{*不一边……一边……}. En pratique, les Chinois procèdent autrement :

\begin{itemize}
\item soit ils nient \textbf{chaque membre séparément} : \zh{他既不看电视，也不听音乐。} \emph{Tā jì bù kàn diànshì, yě bù tīng yīnyuè.} « Il ne regarde pas la télé et n'écoute pas non plus la musique. » ;
\item soit ils \textbf{reformulent} : \zh{他吃饭的时候不看电视。} \emph{Tā chīfàn de shíhou bù kàn diànshì.} « Quand il mange, il ne regarde pas la télé. » ;
\item soit la négation porte sur la possibilité, avec \zh{不能} : \zh{你不能一边开车，一边打电话。} \emph{Nǐ bù néng yìbiān kāichē, yìbiān dǎ diànhuà.} « Tu ne peux pas téléphoner en conduisant. » — ici \zh{不能} nie la possibilité de combiner les deux actions, ce qui est naturel et très fréquent (pensez aux campagnes de sécurité routière, au Cameroun comme en Chine).
\end{itemize}

\courssub{Limite 4 : la place de 了}

\begin{propriete}[\zh{了} et \zh{一边……一边……}]
L'usage courant préfère la structure \textbf{sans} \zh{了}, car elle décrit des actions \textbf{duratives} en cours, et \zh{了} marque l'accomplissement d'actions ponctuelles. Cependant, \zh{了} est \textbf{possible dans chaque membre} pour présenter deux actions ponctuelles réalisées simultanément dans le passé :
\zh{他一边笑了笑，一边点了点头。} \emph{Tā yìbiān xiào le xiào, yìbiān diǎn le diǎn tóu.} « Il a souri tout en hochant la tête. »
\end{propriete}

À votre niveau, retenez la règle simple : \textbf{par défaut, pas de \zh{了}} dans \zh{一边……一边……}. Si vous décrivez une habitude ou une scène en cours, la forme nue est toujours correcte : \zh{她一边唱歌，一边洗衣服。} \emph{Tā yìbiān chànggē, yìbiān xǐ yīfu.} « Elle chante en faisant la lessive. »

\courssub{Pour l'examen : la variante écrite 一面……一面……}

\begin{definition}[\zh{一面……一面……}]
\zh{一面……一面……} \emph{yìmiàn……yìmiàn……} est la \textbf{variante écrite et soutenue} de \zh{一边……一边……}. Le sens est identique (deux actions simultanées du même sujet), mais le registre est plus littéraire : on la rencontre dans les articles de presse, les dissertations et les textes officiels. À l'oral courant, on préfère \zh{一边……一边……}. À l'oral familier, on entend aussi la forme raccourcie \zh{边……边……} : \zh{我们边走边说。} \emph{Wǒmen biān zǒu biān shuō.} « Nous parlons en marchant. »
\end{definition}

\begin{exemplebox}[Registre écrit]
\zh{他一面工作，一面学习。}\\
\emph{Tā yìmiàn gōngzuò, yìmiàn xuéxí.}\\
« Il travaille tout en poursuivant ses études. » (phrase typique d'un article de presse)\\[4pt]
\zh{政府一面发展经济，一面保护环境。}\\
\emph{Zhèngfǔ yìmiàn fāzhǎn jīngjì, yìmiàn bǎohù huánjìng.}\\
« Le gouvernement développe l'économie tout en protégeant l'environnement. »
\end{exemplebox}

\begin{center}
\begin{tabularx}{\linewidth}{|X|X|X|X|}
\hline
\rowcolor{popblueL} Structure & Registre & Emploi & Exemple \\
\hline
\zh{一边……一边……} & oral et écrit courant & 2 actions simultanées & \zh{我一边走，一边看。} \\
\hline
\zh{一面……一面……} & écrit soutenu, presse & 2 actions simultanées & \zh{他一面工作，一面学习。} \\
\hline
\zh{先……然后……} & tous registres & 2 actions successives & \zh{先吃饭，然后学习。} \\
\hline
\zh{既……又……} & tous registres & 2 états ou qualités & \zh{他既是学生，又是老师。} \\
\hline
\end{tabularx}
\end{center}

\courssub{Organigramme de décision}

\begin{popfigure}
\begin{tikzpicture}[
node distance=0.9cm,
decision/.style={diamond, draw=popblue, fill=popblueL, text=popdark, align=center, aspect=2.2, inner sep=1pt, font=\small},
action/.style={rectangle, rounded corners, draw=popgreen, fill=popgreenL, text=popdark, align=center, font=\small},
other/.style={rectangle, rounded corners, draw=poporange, fill=poporangeL, text=popdark, align=center, font=\small},
fleche/.style={-{Stealth}, thick, popdark}
]
\node[decision, align=center] (q1){Deux actions\\en même temps ?};
\node[decision, below=of q1, align=center] (q2){Même\\sujet ?};
\node[decision, below=of q2, align=center] (q3){Actions actives\\et compatibles ?};
\node[action, below=of q3, align=center] (ok){一边……一边……\\(écrit soutenu : 一面……一面……)};
\node[other, right=2.6cm of q1] (n1) {先……然后……};
\node[other, right=2.6cm of q2, align=center] (n2){……的时候\\ou deux phrases};
\node[other, right=2.6cm of q3, align=center] (n3){既……又……\\ou reformuler};
\draw[fleche] (q1) -- node[left, font=\small]{oui} (q2);
\draw[fleche] (q2) -- node[left, font=\small]{oui} (q3);
\draw[fleche] (q3) -- node[left, font=\small]{oui} (ok);
\draw[fleche] (q1) -- node[above, font=\small]{non} (n1);
\draw[fleche] (q2) -- node[above, font=\small]{non} (n2);
\draw[fleche] (q3) -- node[above, font=\small]{non} (n3);
\end{tikzpicture}
\legende{Choisir la bonne structure : l'arbre de décision de \zh{一边……一边……}}
\end{popfigure}

\begin{exoresolu}[Choix de structure : simultané ou successif ?]
Énoncé. Pour chaque phrase française, choisis la structure chinoise correcte (\zh{一边……一边……}, \zh{先……然后……} ou \zh{既……又……}), puis traduis en chinois.\\
a) « Mon frère écoute la radio en préparant le petit-déjeuner. »\\
b) « Mon frère prépare le petit-déjeuner, puis il écoute la radio. »\\
c) « Elle est à la fois médecin et professeur. »
\tcblower
Solution détaillée.\\
a) Écouter la radio et préparer le petit-déjeuner peuvent se dérouler \textbf{en même temps} (actions durables, compatibles, même sujet) : \zh{一边……一边……}.\\
\zh{我哥哥一边做早饭，一边听广播。} \emph{Wǒ gēge yìbiān zuò zǎofàn, yìbiān tīng guǎngbō.}\\
b) « Puis » indique clairement une \textbf{succession} : \zh{一边……一边……} est interdit ; on utilise \zh{先……然后……}.\\
\zh{我哥哥先做早饭，然后听广播。} \emph{Wǒ gēge xiān zuò zǎofàn, ránhòu tīng guǎngbō.}\\
c) « Être médecin » et « être professeur » sont des \textbf{états}, pas des actions : on utilise \zh{既……又……}.\\
\zh{她既是医生，又是老师。} \emph{Tā jì shì yīshēng, yòu shì lǎoshī.}
\end{exoresolu}

\begin{exercice}[À vous : correct ou incorrect ?]
Pour chaque phrase, dis si elle est correcte (正确) ou incorrecte (错误), puis corrige les phrases fautives.\\
1. \zh{他一边洗澡，一边看书。}\\
2. \zh{*她一边是医生，一边在医院工作。}\\
3. \zh{*我们先一边吃饭，一边然后看电视。}\\
4. \zh{爸爸一边开车，一边听新闻。}
\end{exercice}

\begin{corrige}[Exercice : correct ou incorrect ?]
1. \textbf{Incorrecte} : se laver et lire un livre sont physiquement incompatibles (le livre serait mouillé !). Reformulation possible : \zh{他先洗澡，然后看书。} « Il se lave d'abord, puis il lit. »\\
2. \textbf{Incorrecte} : \zh{是} est un verbe d'état. Correction : \zh{她是医生，在医院工作。} « Elle est médecin, elle travaille à l'hôpital. »\\
3. \textbf{Incorrecte} : on ne peut pas mélanger \zh{先……然后……} (succession) et \zh{一边……一边……} (simultanéité). Correction : \zh{我们先吃饭，然后看电视。} « Nous mangeons d'abord, puis nous regardons la télé. »\\
4. \textbf{Correcte} : conduire et écouter les informations sont deux actions simultanées, durables et compatibles, accomplies par le même sujet.
\end{corrige}

\begin{attention}[Récapitulatif des pièges à éviter]
\begin{itemize}
\item Ne combine jamais deux actions physiquement incompatibles : \zh{*一边睡觉，一边跑步} est absurde.
\item N'emploie jamais \zh{一边……一边……} avec des verbes d'état (\zh{是}, \zh{有}, \zh{像}) : utilise \zh{既……又……}.
\item Ne traduis pas « il rentre en mangeant » (succession) par \zh{*他一边回家，一边吃饭} : dis \zh{他先回家，然后吃饭。}.
\item N'ajoute pas \zh{了} par réflexe : la forme simple sans \zh{了} est la norme.
\item Évite la négation directe \zh{*不一边……} : reformule avec \zh{不能}, \zh{既不……也不……} ou \zh{……的时候}.
\item Sépare toujours les deux membres par la virgule chinoise \zh{，} : \zh{他一边吃饭，一边看电视。}
\end{itemize}
\end{attention}

\begin{aretenir}[L'essentiel de la section]
\zh{一边……一边……} n'est correcte que si les trois conditions sont réunies : \cle{simultanéité réelle}, \cle{même sujet}, \cle{actions actives, durables et compatibles}. Pour la succession, utilise \zh{先……然后……} ; pour les états, \zh{既……又……} ; pour le registre écrit soutenu, la variante \zh{一面……一面……}. En cas de doute, applique l'arbre de décision : trois questions, trois « oui », et la structure est sûre.
\end{aretenir}

\coursec{Comparaison avec le français et pièges des francophones}

Dans la section précédente, nous avons vu les emplois, les limites et les exceptions de la structure \zh{一边……一边……} : elle exprime deux actions simultanées, accomplies par un même sujet, et elle refuse les actions successives ou incompatibles. Nous allons maintenant confronter cette structure au français. C'est ici que se gagnent — ou se perdent — la plupart des points à l'examen : les erreurs des élèves francophones sont prévisibles, donc évitables. Observons d'abord, puis énonçons les règles, puis traquons les pièges un par un.

\courssub{Observation : ce que le français fait une fois, le chinois le fait deux fois}

\begin{exemplebox}[Observer avant de comprendre]
Comparons quatre phrases françaises et leur traduction chinoise :\\

\zh{他一边唱歌，一边跳舞。} \emph{Tā yìbiān chàng gē, yìbiān tiào wǔ.} « Il chante tout en dansant. »\\

\zh{妈妈一边做饭，一边听音乐。} \emph{Māma yìbiān zuò fàn, yìbiān tīng yīnyuè.} « Maman écoute de la musique en faisant la cuisine. »\\

\zh{我们一边喝茶，一边聊天。} \emph{Wǒmen yìbiān hē chá, yìbiān liáo tiān.} « Nous buvons du thé et nous discutons en même temps. »\\

\zh{学生们一边走路，一边说法语。} \emph{Xuéshengmen yìbiān zǒu lù, yìbiān shuō Fǎyǔ.} « Les élèves parlent français tout en marchant. »\\

Que remarque-t-on ? Le français utilise un seul marqueur (\emph{en + gérondif}, \emph{tout en}, \emph{en même temps}) placé devant une seule des deux actions. Le chinois, lui, place \zh{一边} devant \cle{chacun} des deux verbes, sans exception.
\end{exemplebox}

\begin{attention}[La source d'erreur n°1 des francophones]
En français, le marqueur de simultanéité apparaît \textbf{une seule fois} : « Il chante \emph{en} dansant » — on ne dit jamais « \emph{en} chantant, \emph{en} dansant ». En chinois, c'est exactement l'inverse : le marqueur \zh{一边} se \cle{répète obligatoirement} devant chaque verbe : \zh{一边……一边……}. Oublier le second \zh{一边} est la faute la plus fréquente des élèves camerounais francophones. Retenez : \textbf{deux actions = deux 一边}.
\end{attention}

\courssub{Les correspondances françaises de 一边……一边……}

Le français possède plusieurs outils pour exprimer la simultanéité. Le chinois, dans les phrases simples du niveau HSK 2-3, n'en utilise qu'un : la structure \zh{一边……一边……}. Voici le tableau des correspondances.

\begin{center}
\begin{tabularx}{\linewidth}{|X|X|X|}
\hline
\rowcolor{popblueL} \textbf{Français} & \textbf{Structure chinoise} & \textbf{Exemple traduit} \\
\hline
tout en + gérondif & 一边……一边…… & \zh{他一边吃饭，一边看手机。} \emph{Tā yìbiān chī fàn, yìbiān kàn shǒujī.} « Il regarde son téléphone tout en mangeant. » \\
\hline
en + gérondif & 一边……一边…… & \zh{她一边跑步，一边听广播。} \emph{Tā yìbiān pǎo bù, yìbiān tīng guǎngbō.} « Elle écoute la radio en courant. » \\
\hline
en même temps (que) & 一边……一边…… & \zh{我们一边工作，一边学习汉语。} \emph{Wǒmen yìbiān gōngzuò, yìbiān xuéxí Hànyǔ.} « Nous travaillons et apprenons le chinois en même temps. » \\
\hline
à la fois (deux actions) & 一边……一边…… & \zh{他一边开车，一边打电话。} \emph{Tā yìbiān kāi chē, yìbiān dǎ diànhuà.} « Il conduit et téléphone à la fois. » \\
\hline
\end{tabularx}
\end{center}

\begin{propriete}[Règle de correspondance]
Quelle que soit la formule française (\emph{tout en}, \emph{en + gérondif}, \emph{en même temps}, \emph{à la fois}), la traduction chinoise standard au lycée est toujours :

\textbf{Sujet + 一边 + Verbe 1 (+ complément)，一边 + Verbe 2 (+ complément)。}

Le français choisit souvent de subordonner une action à l'autre (gérondif) ; le chinois met les deux actions sur le même plan, chacune précédée de \zh{一边}. Il n'y a pas de hiérarchie entre les deux verbes chinois.
\end{propriete}

\begin{popfigure}
\begin{tikzpicture}[node distance=0.4cm and 1.2cm]
\node[draw=popblue, fill=popblueL, rounded corners, align=center, font=\small] (fr) {Français :\\ Il chante \textbf{en} dansant.\\ (un seul marqueur)};
\node[draw=poporange, fill=poporangeL, rounded corners, align=center, font=\small, right=of fr] (zh) {Chinois :\\ 他\textbf{一边}唱歌，\textbf{一边}跳舞。\\ \emph{Tā yìbiān chàng gē, yìbiān tiào wǔ.}};
\node[draw=popgreen, fill=popgreenL, rounded corners, align=center, font=\small, below=1.1cm of zh, xshift=-1.6cm] (v1) {一边 + 唱歌\\ (action 1)};
\node[draw=popgreen, fill=popgreenL, rounded corners, align=center, font=\small, right=of v1] (v2) {一边 + 跳舞\\ (action 2)};
\draw[-{Stealth}, thick, popdark] (fr) -- (zh) node[midway, above, font=\small\itshape] {traduction};
\draw[-{Stealth}, thick, poporange] (zh.south) -- (v1.north);
\draw[-{Stealth}, thick, poporange] (zh.south) -- (v2.north);
\draw[{Stealth}-{Stealth}, very thick, poppurple] (v1.east) -- (v2.west) node[midway, above, font=\small\bfseries] {répétition obligatoire};
\end{tikzpicture}
\legende{Le français marque la simultanéité une fois ; le chinois répète 一边 devant chaque verbe.}
\end{popfigure}

\courssub{Piège 1 : oublier le second 一边}

\begin{exemplebox}[Faux et juste]
FAUX : \zh{*他一边看电视，吃饭。} \emph{*Tā yìbiān kàn diànshì, chī fàn.}\\

JUSTE : \zh{他一边看电视，一边吃饭。} \emph{Tā yìbiān kàn diànshì, yìbiān chī fàn.} « Il regarde la télévision tout en mangeant. »\\

Sans le second \zh{一边}, la phrase est boiteuse : le premier membre annonce une simultanéité que le second membre ne confirme pas. C'est comme si, en français, on disait « tout en regardant la télé, il mange » avec une construction cassée.
\end{exemplebox}

\courssub{Piège 2 : placer le sujet entre les deux 一边}

Le sujet, unique, se place \cle{en tête de phrase}, avant le premier \zh{一边}. Un francophone, influencé par « en chantant, \emph{il} danse », a tendance à répéter le sujet au milieu.

FAUX : \zh{*一边他看电视，一边吃饭。} ou \zh{*一边看电视，他一边吃饭。}\\

JUSTE : \zh{他一边看电视，一边吃饭。}

\begin{propriete}[Place du sujet]
Schéma obligatoire : \textbf{Sujet + 一边 + V1，一边 + V2。} Le sujet n'apparaît qu'\cle{une seule fois}, tout au début. Si les deux actions ont des sujets différents, la structure \zh{一边……一边……} est impossible : il faut deux phrases séparées.
\end{propriete}

\courssub{Piège 3 : le pléonasme avec 同时}

Le mot \zh{同时} \emph{tóngshí} signifie « en même temps ». Le francophone qui traduit mot à mot « il chante et danse en même temps » produit :

LOURD : \zh{?他一边唱歌，一边同时跳舞。}\\

JUSTE : \zh{他一边唱歌，一边跳舞。}

La structure \zh{一边……一边……} contient \cle{déjà} l'idée de simultanéité : ajouter \zh{同时} dans une phrase simple est un pléonasme. \zh{同时} s'emploie plutôt seul, entre deux phrases : \zh{他在学汉语，同时也在工作。} \emph{Tā zài xué Hànyǔ, tóngshí yě zài gōngzuò.} « Il apprend le chinois et, en même temps, il travaille. »

\courssub{Piège 4 : oublier la virgule centrale}

Entre les deux membres, la virgule chinoise \zh{，} est obligatoire à l'écrit :

FAUX : \zh{*他一边唱歌一边跳舞。} (à l'écrit)\\

JUSTE : \zh{他一边唱歌，一边跳舞。}

À l'oral, cette virgule correspond à une petite pause. À l'écrit — et notamment au BEPC et au Probatoire — son oubli est sanctionné comme une faute de ponctuation (nous y reviendrons dans la section 5).

\courssub{Piège 5 : actions incompatibles ou successives}

La structure exige deux actions \cle{réellement simultanées} et compatibles. Le français « en + gérondif » tolère parfois des enchaînements (« il partit en claquant la porte ») ; le chinois est plus strict.

FAUX (actions successives) : \zh{*他一边起床，一边睡觉。} — on ne dort pas en se levant !\\

FAUX (succession logique) : \zh{*我一边吃晚饭，一边做作业。} si l'on dîne \emph{puis} on travaille : dire plutôt \zh{我吃完晚饭以后做作业。} \emph{Wǒ chī wán wǎnfàn yǐhòu zuò zuòyè.} « Après avoir dîné, je fais mes devoirs. »\\

JUSTE : \zh{我一边做作业，一边听音乐。} \emph{Wǒ yìbiān zuò zuòyè, yìbiān tīng yīnyuè.} « Je fais mes devoirs en écoutant de la musique. »

\begin{savaistu}[L'ordre fixe du chinois]
En français, on peut inverser librement : « En mangeant, il regarde la télé » ou « Il regarde la télé en mangeant ». Le chinois, lui, garde un ordre \cle{fixe} : sujet → action 1 → action 2. On ne peut pas mettre un membre \zh{一边} devant le sujet. Cette rigidité est une chance pour vous : une seule structure à mémoriser, aucune inversion possible, donc aucune hésitation à l'examen !
\end{savaistu}

\begin{methode}[Vérifier sa phrase en trois contrôles]
Avant de rendre votre copie, appliquez ce contrôle systématique à toute phrase avec \zh{一边……一边……} :
\begin{enumerate}
\item \textbf{Compter les 一边} : il en faut exactement \textbf{deux}, un devant chaque verbe.
\item \textbf{Vérifier le sujet unique en tête} : un seul sujet, placé avant le premier \zh{一边}, jamais entre les deux.
\item \textbf{Vérifier la virgule centrale} : \zh{，} entre les deux membres, et aucun mot superflu (\zh{同时}, \zh{和}) entre eux.
\end{enumerate}
Bonus : demandez-vous si les deux actions peuvent vraiment se faire en même temps. Si non, changez de structure (\zh{以后}, \zh{先……然后……}).
\end{methode}

\begin{exoresolu}[Exercice flash de repérage d'erreurs]
Énoncé — Les quatre phrases suivantes, écrites par des élèves de Première A, contiennent chacune une faute. Trouvez l'erreur et corrigez la phrase.

1. \zh{*他一边听广播，看报纸。}

2. \zh{*一边她做饭，一边打电话。}

3. \zh{*我们一边踢足球一边同时唱歌。}

4. \zh{*我一边睡觉，一边吃早饭。}

\tcblower
Solution détaillée —

1. Il manque le second \zh{一边} (piège 1). Correction : \zh{他一边听广播，一边看报纸。} \emph{Tā yìbiān tīng guǎngbō, yìbiān kàn bàozhǐ.} « Il écoute la radio en lisant le journal. »

2. Le sujet \zh{她} est mal placé : il doit précéder le premier \zh{一边} (piège 2). Correction : \zh{她一边做饭，一边打电话。} \emph{Tā yìbiān zuò fàn, yìbiān dǎ diànhuà.} « Elle téléphone tout en faisant la cuisine. »

3. Double faute : virgule centrale absente (piège 4) et \zh{同时} pléonastique (piège 3). Correction : \zh{我们一边踢足球，一边唱歌。} \emph{Wǒmen yìbiān tī zúqiú, yìbiān chàng gē.} « Nous jouons au football en chantant. »

4. Actions incompatibles : on ne peut pas dormir et prendre son petit-déjeuner en même temps (piège 5). Correction possible avec succession : \zh{我先吃早饭，然后睡觉。} \emph{Wǒ xiān chī zǎofàn, ránhòu shuì jiào.} « Je prends d'abord mon petit-déjeuner, puis je dors. » Ou choisir deux actions compatibles : \zh{我一边吃早饭，一边听新闻。} « Je prends mon petit-déjeuner en écoutant les informations. »
\end{exoresolu}

\begin{exercice}[À vous de jouer]
Traduisez en chinois, puis vérifiez chaque phrase avec la méthode des trois contrôles :

1. Mon père lit le journal en buvant son café.

2. Les élèves du lycée de Yaoundé chantent tout en dansant.

3. Elle apprend le chinois et prépare l'examen en même temps.

4. Corrigez : \zh{*一边我说汉语，一边写汉字。}
\end{exercice}

\begin{corrige}[Exercice « À vous de jouer »]
1. \zh{我爸爸一边看报纸，一边喝咖啡。} \emph{Wǒ bàba yìbiān kàn bàozhǐ, yìbiān hē kāfēi.}

2. \zh{雅温得中学的学生们一边唱歌，一边跳舞。} \emph{Yǎwēndé zhōngxué de xuéshengmen yìbiān chàng gē, yìbiān tiào wǔ.}

3. \zh{她一边学汉语，一边准备考试。} \emph{Tā yìbiān xué Hànyǔ, yìbiān zhǔnbèi kǎoshì.} (Ne pas ajouter \zh{同时} !)

4. Sujet mal placé : \zh{我一边写汉字，一边说汉语。} \emph{Wǒ yìbiān xiě Hànzì, yìbiān shuō Hànyǔ.} « J'écris des caractères tout en parlant chinois. »
\end{corrige}

\begin{aretenir}[L'essentiel de la section]
\begin{itemize}
\item Le français marque la simultanéité \textbf{une fois} (\emph{en}, \emph{tout en}) ; le chinois répète \zh{一边} devant \textbf{chaque} verbe.
\item Schéma unique et fixe : \textbf{Sujet + 一边 + V1，一边 + V2。}
\item Jamais de sujet entre les deux \zh{一边}, jamais de \zh{同时} en plus, toujours la virgule \zh{，} au centre.
\item Les deux actions doivent être simultanées et compatibles ; sinon, utilisez \zh{先……然后……} ou \zh{以后}.
\end{itemize}
\end{aretenir}

\coursec{Leçon 40 — Grammaire : \zh{一边……一边……} ; La ponctuation en chinois}

\courssub{Observation et découverte de la structure}

\begin{textebox}[Dialogue : Une soirée à Yaoundé]
\zh{— 小王，你昨晚做什么了？}\\
\zh{— 我一边做作业，一边听收音机。收音机里放《新闻联播》，我一边听，一边记笔记。}\\
\zh{— 真认真！我呢，一边看手机，一边吃薯片，结果作业没写完。}\\
\emph{— Xiǎo Wáng, nǐ zuówǎn zuò shénme le ?}\\
\emph{— Wǒ yìbiān zuò zuòyè, yìbiān tīng shōuyīnjī. Shōuyīnjī li fàng 《Xīnwén Liánbō》, wǒ yìbiān tīng, yìbiān jì bǐjì.}\\
\emph{— Zhēn rènzhēn ! Wǒ ne, yìbiān kàn shǒujī, yìbiān chī shǔpiàn, jiéguǒ zuòyè méi xiě wán.}\\
« — Xiao Wang, qu'as-tu fait hier soir ?\\
— J'ai fait mes devoirs tout en écoutant la radio. La radio diffusait le Journal télévisé national, je l'écoutais tout en prenant des notes.\\
— Vraiment sérieux ! Moi, je regardais mon téléphone tout en mangeant des chips, résultat : mes devoirs n'étaient pas finis. »
\end{textebox}

\begin{definition}[La structure \zh{一边……一边……} (\emph{yìbiān……yìbiān……})]
La structure \zh{一边……一边……} exprime la \cle{simultanéité de deux actions} effectuées par \cle{le même sujet}. Elle correspond au français « tout en faisant A, faire B » ou « faire A pendant qu'on fait B ». Les deux actions sont généralement \cle{duratives} (elles durent dans le temps).
\end{definition}

\begin{exemplebox}[Analyse du dialogue]
\begin{itemize}
\item \zh{我一边做作业，一边听收音机。} \emph{Wǒ yìbiān zuò zuòyè, yìbiān tīng shōuyīnjī.} « Je fais mes devoirs tout en écoutant la radio. » (Deux actions duratives, même sujet \zh{我}).
\item \zh{我一边听，一边记笔记。} \emph{Wǒ yìbiān tīng, yìbiān jì bǐjì.} « J'écoute tout en prenant des notes. » (Verbes monosyllabiques, objet après le second verbe).
\item \zh{我一边看手机，一边吃薯片。} \emph{Wǒ yìbiān kàn shǒujī, yìbiān chī shǔpiàn.} « Je regarde mon téléphone tout en mangeant des chips. »
\end{itemize}
\end{exemplebox}

\courssub{Schéma et construction de la structure}

\begin{propriete}[Schéma syntaxique]
\zh{Sujet + 一边 + Verbe₁ (+ Complément) + ， + 一边 + Verbe₂ (+ Complément) 。}
\end{propriete}

\begin{center}
\begin{tabularx}{\linewidth}{|X|X|X|}
\hline
\rowcolor{popblueL} Composant & Exemple & Rôle \\
\hline
Sujet & \zh{我} (\emph{wǒ}) / \zh{他} (\emph{tā}) / \zh{小李} (\emph{Xiǎo Lǐ}) & Doit être \textbf{identique} pour les deux actions. \\
\hline
\zh{一边} (\emph{yìbiān}) & \zh{一边} & Marqueur de simultanéité (premier volet). \\
\hline
Verbe₁ (+ Compl.) & \zh{做作业} (\emph{zuò zuòyè}) / \zh{听音乐} (\emph{tīng yīnyuè}) & Action 1 (durative). \\
\hline
\zh{，} (\emph{dòuhào}) & \zh{，} & \textbf{Obligatoire} : sépare les deux propositions. \\
\hline
\zh{一边} (\emph{yìbiān}) & \zh{一边} & Marqueur de simultanéité (second volet). \\
\hline
Verbe₂ (+ Compl.) & \zh{看电视} (\emph{kàn diànshì}) / \zh{吃饭} (\emph{chīfàn}) & Action 2 (durative), souvent l'action « principale ». \\
\hline
\end{tabularx}
\end{center}

\begin{attention}[Règles de construction strictes]
\begin{enumerate}
\item \textbf{Même sujet :} Le sujet des deux actions est unique et placé \emph{avant} le premier \zh{一边}. On ne le répète pas. \textbf{Faux :} \zh{我一边做作业，我一边听音乐。}
\item \textbf{Virgule obligatoire :} La virgule \zh{，} (pleine chasse) sépare \cle{impérativement} les deux volets. \textbf{Faux :} \zh{我一边做作业一边听音乐。}
\item \textbf{Verbes duratifs :} Les verbes doivent exprimer une action qui dure (étudier, manger, marcher, écouter, regarder). Les verbes d'état (\zh{喜欢}, \zh{知道}) ou d'action ponctuelle achevée (\zh{买}, \zh{死}) ne s'utilisent généralement pas.
\item \textbf{Ordre des actions :} L'action après le second \zh{一边} est souvent l'action « principale » ou celle sur laquelle on porte l'attention, mais les deux sont simultanées.
\end{enumerate}
\end{attention}

\courssub{Emplois, limites et exceptions}

\begin{propriete}[Emplois principaux]
\begin{enumerate}
\item \textbf{Deux actions volontaires simultanées :} \zh{她一边唱歌，一边跳舞。} \emph{Tā yìbiān chànggē, yìbiān tiàowǔ.} « Elle chante tout en dansant. »
\item \textbf{Action principale + action accompagnatrice :} \zh{请一边走，一边看路。} \emph{Qǐng yìbiān zǒu, yìbiān kàn lù.} « Marchez tout en regardant la route. » (Impératif / Conseil).
\item \textbf{Description de scène :} \zh{大街上，人们一边走路，一边看手机。} \emph{Dàjiē shàng, rénmen yìbiān zǒulù, yìbiān kàn shǒujī.} « Dans la rue, les gens marchent tout en regardant leur téléphone. »
\end{enumerate}
\end{propriete}

\begin{propriete}[Limites et cas particuliers]
\begin{itemize}
\item \textbf{Négation :} On nie l'ensemble avec \zh{没(有)} ou \zh{不} placé \emph{avant} le premier \zh{一边} (rare) ou on nie le résultat. Ex: \zh{他没一边吃饭一边说话。} (Il n'a pas mangé en parlant). Plus naturel : \zh{他吃饭的时候没说话。}
\item \textbf{Aspect \zh{着} (zhe) :} Pour insister sur l'état d'accompagnement, on peut ajouter \zh{着} au V₁ : \zh{他一边听着音乐，一边睡着了。} \emph{Tā yìbiān tīngzhe yīnyuè, yìbiān shuìzháo le.} « Il écoutait de la musique (en fond) et s'est endormi. » (Structure avancée HSK 4).
\item \textbf{Verbes monosyllabiques + Objet :} Si V₁ et V₂ sont monosyllabiques et partagent l'objet, l'objet se place souvent à la fin : \zh{他一边唱一边跳歌。} (Peu courant). Préférer : \zh{他一边唱歌，一边跳舞。}
\end{itemize}
\end{propriete}

\begin{center}
\begin{tabularx}{\linewidth}{|X|X|X|}
\hline
\rowcolor{popblueL} Phrase correcte & Phrase incorrecte & Explication \\
\hline
\zh{我一边吃饭，一边看电视。} & \zh{我吃饭，一边看电视。} & Manque le premier \zh{一边}. \\
\hline
\zh{他一边跑步，一边听歌。} & \zh{他一边跑步，听歌。} & Manque le second \zh{一边}. \\
\hline
\zh{我们一边聊天，一边喝茶。} & \zh{我们一边聊天，我们一边喝茶。} & Sujet répété (interdit). \\
\hline
\zh{姐姐一边洗衣服，一边听广播。} & \zh{姐姐洗衣服，一边听广播。} & Structure incomplète (manque \zh{一边} devant V₁). \\
\hline
\end{tabularx}
\end{center}

\courssub{Comparaison avec le français et pièges des francophones}

\begin{center}
\begin{tabularx}{\linewidth}{|X|X|}
\hline
\rowcolor{popblueL} Français & Chinois (\zh{一边……一边……}) \\
\hline
« Tout en + gérondif, je fais... » / « Je fais... tout en + gérondif » & \zh{Sujet + 一边 + V₁, 一边 + V₂} (Ordre fixe : V₁ puis V₂). \\
\hline
« Je fais A pendant que je fais B » (subordonnée) & Structure \textbf{parataxique} (coordination), pas de subordination. \\
\hline
Sujet répété possible : « \textbf{Je} mange, \textbf{je} lis. » & Sujet \textbf{unique} en tête, jamais répété. \\
\hline
Pas de signe obligatoire entre les deux verbes & \textbf{Virgule \zh{，} obligatoire} entre les deux volets. \\
\hline
Négation : « Je ne mange pas en lisant. » & Négation globale préposée : \zh{不/没 + 一边……一边……} (rare) ou reformulation. \\
\hline
\end{tabularx}
\end{center}

\begin{attention}[Les 3 pièges majeurs pour francophones]
\begin{enumerate}
\item \textbf{Oublier le premier \zh{一边}} : \emph{« Je mange, tout en lisant »} $\rightarrow$ \zh{我吃饭，一边看书。} \textbf{FAUX}. Correct : \zh{我一边吃饭，一边看书。}
\item \textbf{Répéter le sujet} : \zh{我一边吃饭，我一边看书。} \textbf{FAUX}. Le sujet \zh{我} ne s'écrit qu'une fois.
\item \textbf{Oublier la virgule \zh{，}} : \zh{我一边吃饭一边看书。} \textbf{FAUTE DE PONCTUATION} (pénalisée à l'examen). La virgule marque la frontière prosodique et syntaxique entre les deux procès.
\end{enumerate}
\end{attention}

\begin{methode}[Construire une phrase avec \zh{一边……一边……}]
\begin{enumerate}
\item Identifier le \textbf{sujet unique} (qui fait les deux actions).
\item Choisir deux \textbf{verbes d'action durative} (V₁, V₂).
\item Construire : \zh{Sujet + 一边 + V₁ (+ compl.), 一边 + V₂ (+ compl.)。}
\item Vérifier : \cle{Sujet unique} ? \cle{Deux \zh{一边}} ? \cle{Virgule \zh{，}} ? \cle{Verbes duratifs} ?
\end{enumerate}
\emph{Exemple : Sujet = \zh{爸爸} ; V₁ = \zh{开车} ; V₂ = \zh{听新闻} $\rightarrow$ \zh{爸爸一边开车，一边听新闻。}}
\end{methode}

\courssub{La ponctuation en chinois}

\begin{textebox}[Un message sur WeChat]
\zh{妈妈问：“你今晚回家吃饭吗？”我说：“回！我买了鱼、青菜和水果。”妈妈很高兴。她一边做饭，一边听音乐……家里真热闹！}\\
Māma wèn: “Nǐ jīnwǎn huí jiā chīfàn ma?” Wǒ shuō: “Huí! Wǒ mǎi le yú、qīngcài hé shuǐguǒ.” Māma hěn gāoxìng. Tā yìbiān zuòfàn, yìbiān tīng yīnyuè…… Jiā lǐ zhēn rènao!\\
« Maman demande : “Tu rentres dîner à la maison ce soir ?” Je réponds : “Oui ! J'ai acheté du poisson, des légumes verts et des fruits.” Maman est très contente. Tout en cuisinant, elle écoute de la musique… Comme c'est animé à la maison ! »
\end{textebox}

Observons ce message. On y trouve : des deux-points \zh{：} qui annoncent les paroles, des guillemets \zh{“ ”}, un point d'interrogation \zh{？}, un point d'exclamation \zh{！}, une virgule énumérative \zh{、} entre \zh{鱼} et \zh{青菜}, une virgule ordinaire \zh{，} dans la phrase avec \zh{一边……一边……}, des points de suspension \zh{……} et enfin le petit cercle \zh{。} qui clôt les phrases déclaratives. Chacun de ces signes a sa forme et sa fonction propres.

\begin{propriete}[La règle d'or de la ponctuation chinoise : la pleine chasse]
Tout signe de ponctuation chinois occupe la place d'\cle{un caractère entier} : on dit qu'il est à \cle{pleine chasse}. Concrètement :
\begin{itemize}
\item Le signe remplit un carré d'écriture complet, comme un caractère (ex: \zh{我} \zh{是} \zh{学} \zh{生} \zh{。} = 5 cases).
\item On ne laisse \cle{jamais d'espace} ni avant ni après le signe : on écrit \zh{我是学生。} et non « 我是学生 。 ».
\item Un signe de ponctuation ne se place \cle{jamais en début de ligne} : si la phrase se termine en fin de ligne, le signe reste « collé » au dernier caractère, quitte à dépasser légèrement la marge.
\end{itemize}
\end{propriete}

\begin{popfigure}
\begin{tikzpicture}[scale=1.1]
\draw[step=1, popline, thin] (0,0) grid (5,1);
\draw[popblue, dashed] (0,0) -- (1,1);
\draw[popblue, dashed] (0,1) -- (1,0);
\draw[popblue, dashed] (0.5,0) -- (0.5,1);
\draw[popblue, dashed] (0,0.5) -- (1,0.5);
\node[font=\Huge, popdark] at (0.5,0.5) {我};
\node[font=\Huge, popdark] at (1.5,0.5) {是};
\node[font=\Huge, popdark] at (2.5,0.5) {学};
\node[font=\Huge, popdark] at (3.5,0.5) {生};
\node[font=\Huge, poporange] at (4.5,0.5) {。};
\draw[poporange, dashed, thick] (4,0) rectangle (5,1);
\node[poporange, align=center, font=\small] at (4.5,-0.6) {le signe occupe\\une case entière};
\end{tikzpicture}
\legende{Chaque caractère et chaque signe de ponctuation occupe un carré complet : \zh{我是学生。} = 5 cases.}
\end{popfigure}

\begin{attention}[Pas d'espace avant la ponctuation — Erreur n°1 des francophones]
En français, on met une espace fine avant \textbf{? ! : ;}. En chinois, c'est \cle{interdit} : le signe suit immédiatement le dernier caractère.
\begin{itemize}
\item Correct : \zh{你来吗？} (\emph{Nǐ lái ma?}) — \zh{太好了！} (\emph{Tài hǎo le!})
\item Incorrect : « 你来吗 ？ » — « 太好了 ！ »
\end{itemize}
C'est l'erreur de présentation la plus fréquente et la plus pénalisée en examen.
\end{attention}

\courssub{Le point 。 et la virgule ，}

\begin{definition}[Le point 。 (\emph{jùhào})]
Le point chinois \zh{。} est un \cle{petit cercle plein-chasse}, tracé dans le coin inférieur gauche de sa case. Il marque la fin d'une phrase déclarative.
\zh{我是学生。} \emph{Wǒ shì xuésheng.} « Je suis élève. »
\end{definition}

\begin{definition}[La virgule ， (\emph{dòuhào})]
La virgule \zh{，} marque une pause à l'intérieur de la phrase : elle sépare des \cle{propositions} ou des \cle{membres de phrase}. C'est elle que l'on utilise \textbf{obligatoirement} entre les deux volets de la structure \zh{一边……一边……} :
\zh{他一边吃饭，一边看报纸。} \emph{Tā yìbiān chīfàn, yìbiān kàn bàozhǐ.} « Il mange tout en lisant le journal. »
\end{definition}

\begin{exemplebox}[。 et ， en action]
\zh{我喜欢喀麦隆。雅温得很大，杜阿拉很热闹。}\\
\emph{Wǒ xǐhuan Kāmàilóng. Yǎwēndé hěn dà, Dù'ālā hěn rènao.}\\
« J'aime le Cameroun. Yaoundé est grande, Douala est animée. »
Ici, \zh{。} clôt chaque phrase complète ; \zh{，} sépare deux propositions coordonnées au sein d'une même phrase.
\end{exemplebox}

\courssub{La virgule énumérative 、 : un signe sans équivalent français}

\begin{definition}[La virgule énumérative 、 (\emph{dùnhào})]
Le signe \zh{、} est une petite barre oblique pleine chasse, \cle{spécifique au chinois}. Il sert \cle{uniquement} à séparer des mots ou groupes de mots de même nature (noms, verbes, adjectifs) dans une énumération. Le dernier élément est introduit par \zh{和} (\emph{hé}, « et »), sans signe devant.
\zh{我买了书、笔和本子。} \emph{Wǒ mǎi le shū、bǐ hé běnzi.} « J'ai acheté des livres, des stylos et des cahiers. »
\end{definition}

\begin{attention}[Ne jamais confondre ， et 、 — Piège n°2]
Le français utilise la même virgule pour les deux fonctions. Le chinois les distingue radicalement.
\begin{itemize}
\item \zh{、} = énumération de \cle{mots coordonnés} (noms, adjectifs) : \zh{我喜欢看报纸、杂志和新闻。} \emph{Wǒ xǐhuan kàn bàozhǐ、zázhì hé xīnwén.} (Noms énumérés).
\item \zh{，} = pause entre \cle{propositions} (sujet + verbe) : \zh{他一边喝茶，一边看报纸。} \emph{Tā yìbiān hē chá, yìbiān kàn bàozhǐ.} (Deux actions/propositions).
\end{itemize}
\textbf{Astuce :} Si tu peux remplacer le signe par « et » entre deux \textbf{noms}, utilise \zh{、}. Si le signe sépare deux \textbf{propositions verbales}, utilise \zh{，}.
\end{attention}

\courssub{Le point-virgule ； et les deux-points ：}

\begin{definition}[Le point-virgule ； (\emph{fēnhào})]
Le point-virgule \zh{；} sépare des \cle{propositions parallèles}, de construction semblable, au sein d'une même phrase complexe.
\zh{白天，他工作；晚上，他学习。} \emph{Báitiān, tā gōngzuò; wǎnshang, tā xuéxí.} « Le jour, il travaille ; le soir, il étudie. »
\end{definition}

\begin{definition}[Les deux-points ： (\emph{màohào})]
Les deux-points \zh{：} annoncent une \cle{citation}, une \cle{explication}, une \cle{raison} ou une \cle{liste}. Ils précèdent presque toujours les paroles rapportées au discours direct.
\zh{老师说：“明天有考试。”} \emph{Lǎoshī shuō: “Míngtiān yǒu kǎoshì.”} « Le professeur dit : “Demain, il y a un contrôle.” »
\end{definition}

\courssub{Le point d'interrogation ？ et le point d'exclamation ！}

Ces deux signes existent en chinois sous une forme \cle{pleine chasse} : \zh{？} et \zh{！}. Leur emploi est proche du français, avec la différence majeure : \cle{aucune espace} ne les précède.

\begin{exemplebox}[Interrogation et exclamation]
\zh{你明天来吗？} \emph{Nǐ míngtiān lái ma?} « Tu viens demain ? »\\
\zh{太好了！} \emph{Tài hǎo le!} « C'est formidable ! »\\
\zh{杜阿拉真大！} \emph{Dù'ālā zhēn dà!} « Douala est vraiment grande ! »
\end{exemplebox}

\begin{propriete}[Point d'interrogation et particule \zh{吗}]
Même avec le point d'interrogation \zh{？}, la particule interrogative \zh{吗} reste \textbf{indispensable} dans une question fermée (oui/non) ; le signe \zh{？} ne la remplace pas.
\zh{你是学生吗？} (Correct) \quad vs \quad \zh{你是学生？} (Incomplet / Oral seulement).
\end{propriete}

\begin{exemplebox}[Les trois signes réunis : le discours direct]
\zh{他问：“你明天来吗？”}\\
\emph{Tā wèn: “Nǐ míngtiān lái ma?”}\\
« Il demande : “Tu viens demain ?” »
On observe la combinaison typique : \cle{deux-points} \zh{：} + \cle{guillemets} \zh{“ ”} + \cle{point d'interrogation} \zh{？} placé \emph{à l'intérieur} des guillemets.
\end{exemplebox}

\courssub{Les guillemets “ ” et les marques de titre 《》}

\begin{definition}[Les guillemets “ ” (\emph{yǐnhào})]
Les guillemets chinois \zh{“ ”} encadrent les citations, les paroles, les dialogues ou les mots mis en relief. Ils sont pleins chasse.
\zh{他说：“我是喀麦隆人。”} \emph{Tā shuō: “Wǒ shì Kāmàilóng rén.”}
\end{definition}

\begin{definition}[Les marques de titre 《》 (\emph{shūmínghào})]
Les signes \zh{《》} entourent les \cle{titres d'œuvres} : livres, journaux, films, chansons, émissions de télévision, pièces de théâtre, lois, documents. Là où le français utilise l'italique ou les guillemets, le chinois écrit le titre entre \zh{《} et \zh{》}.
\begin{itemize}
\item \zh{《人民日报》} \emph{Rénmín Rìbào} « Le Quotidien du Peuple » (journal).
\item \zh{《红楼梦》} \emph{Hónglóu Mèng} « Le Rêve dans le pavillon rouge » (roman).
\item \zh{我在看《新闻联播》。} \emph{Wǒ zài kàn 《Xīnwén Liánbō》.} « Je regarde le Journal télévisé national. »
\end{itemize}
\end{definition}

\begin{attention}[《》 n'est pas un guillemet ordinaire — Piège n°3]
Ne traduis jamais \zh{《》} par « » dans tes copies en pensant que ce sont des guillemets de citation.
\begin{itemize}
\item \zh{《人民日报》} = Titre du journal (italique en français : \emph{Le Quotidien du Peuple}).
\item \zh{“人民日报”} = Citation des trois mots « Quotidien du Peuple » (guillemets en français).
\end{itemize}
Écrire \zh{“人民日报”} pour désigner le nom du journal est une \textbf{faute de ponctuation}.
\end{attention}

\courssub{Les points de suspension …… et le tiret ——}

\begin{definition}[Suspension et tiret]
\begin{itemize}
\item \textbf{Points de suspension \zh{……} (\emph{shěnglüèhào}) :} S'écrivent avec \cle{six points} (deux groupes de trois), occupant \cle{deux cases} pleines chasse. Ils marquent une énumération inachevée, une hésitation ou une pensée en suspens.
\zh{市场里有很多东西：水果、蔬菜、鱼……} \emph{Shìchǎng li yǒu hěn duō dōngxi: shuǐguǒ、shūcài、yú……} « Il y a beaucoup de choses au marché : fruits, légumes, poisson… »
\item \textbf{Tiret \zh{——} (\emph{pòzhéhào}) :} Est \cle{double} (deux tirets pleins chasse), occupant \cle{deux cases}. Il introduit une explication, une conclusion ou marque une rupture brutale.
\zh{他只有一个愿望——去中国学习。} \emph{Tā zhǐ yǒu yí ge yuànwàng —— qù Zhōngguó xuéxí.} « Il n'a qu'un souhait : aller étudier en Chine. »
\end{itemize}
\end{definition}

\courssub{Tableau récapitulatif des signes de ponctuation}

\begin{center}
\begin{tabularx}{\linewidth}{|X|X|X|X|}
\hline
\rowcolor{popblueL} Signe & Nom (pinyin) & Équivalent français & Exemple \\
\hline
\zh{。} & 句号 \emph{jùhào} & point (.) & \zh{我是学生。} \\
\hline
\zh{，} & 逗号 \emph{dòuhào} & virgule (,) & \zh{他一边吃，一边说。} \\
\hline
\zh{、} & 顿号 \emph{dùnhào} & aucun (virgule d'énumération) & \zh{书、笔和本子} \\
\hline
\zh{；} & 分号 \emph{fēnhào} & point-virgule (;) & \zh{白天工作；晚上学习。} \\
\hline
\zh{：} & 冒号 \emph{màohào} & deux-points (:) & \zh{他说：“你好。”} \\
\hline
\zh{？} & 问号 \emph{wènhào} & point d'interrogation (?) & \zh{你来吗？} \\
\hline
\zh{！} & 叹号 \emph{tànhào} & point d'exclamation (!) & \zh{太好了！} \\
\hline
\zh{“ ”} & 引号 \emph{yǐnhào} & guillemets (« ») & \zh{“你明天来吗？”} \\
\hline
\zh{《》} & 书名号 \emph{shūmínghào} & italique / guillemets de titre & \zh{《人民日报》} \\
\hline
\zh{……} & 省略号 \emph{shěnglüèhào} & points de suspension (…) & \zh{水果、蔬菜……} \\
\hline
\zh{——} & 破折号 \emph{pòzhéhào} & tiret (—) & \zh{愿望——去中国} \\
\hline
\end{tabularx}
\end{center}

\begin{popfigure}
\begin{tikzpicture}[font=\small, node distance=0.4cm and 1.2cm]
\node[draw=popblue, fill=popblueL, rounded corners, align=center, font=\bfseries] (centre) {Ponctuation\\chinoise};
\node[draw=popgreen, fill=popgreenL, rounded corners, align=center, above left=0.6cm and 2.8cm of centre] (fin) {Fin de phrase\\\zh{。 ？ ！}};
\node[draw=poporange, fill=poporangeL, rounded corners, align=center, left=2.8cm of centre] (pause) {Pauses internes\\\zh{， 、 ；}};
\node[draw=poppurple, fill=poppurpleL, rounded corners, align=center, below left=0.6cm and 2.8cm of centre] (annonce) {Annonce et titres\\\zh{： “ ” 《》}};
\node[draw=popgold, fill=popgoldL, rounded corners, align=center, below=0.9cm of centre] (autres) {Rupture / Suite\\\zh{…… ——}};
\draw[-{Stealth}, thick, popblue] (centre) -- (fin);
\draw[-{Stealth}, thick, popblue] (centre) -- (pause);
\draw[-{Stealth}, thick, popblue] (centre) -- (annonce);
\draw[-{Stealth}, thick, popblue] (centre) -- (autres);
\end{tikzpicture}
\legende{Carte mentale : les quatre grandes fonctions de la ponctuation chinoise.}
\end{popfigure}

\courssub{Comparaison systématique avec le français}

\begin{center}
\begin{tabularx}{\linewidth}{|X|X|X|}
\hline
\rowcolor{popblueL} Point de comparaison & Français & Chinois \\
\hline
Espace avant ? ! : ; & Obligatoire (espace fine) & Interdite : \zh{你来吗？} \\
\hline
Fin de phrase déclarative & Point rond . & Petit cercle \zh{。} \\
\hline
Énumération de noms & Virgule , & Signe spécial \zh{、} \\
\hline
Titres d'œuvres & Italique ou « » & Marques \zh{《》} \\
\hline
Points de suspension & Trois points … & Six points \zh{……} (2 cases) \\
\hline
Tiret & Simple — & Double \zh{——} (2 cases) \\
\hline
Taille du signe & Demi-chasse & Pleine chasse (1 case) \\
\hline
Début de ligne & Ponctuation autorisée & Ponctuation \textbf{interdite} en début de ligne \\
\hline
\end{tabularx}
\end{center}

\begin{exoresolu}[Ponctuer correctement un texte intégrant \zh{一边……一边……}]
Énoncé : Recopie ce texte en ajoutant la ponctuation chinoise correcte (caractères, pinyin et traduction) :
\zh{老师问 你们昨天做了什么 小王说 我买了苹果 香蕉和橙子 然后 我一边吃水果 一边看电视 真舒服}
\tcblower
Solution détaillée, étape par étape :
\begin{enumerate}
\item \zh{老师问} annonce une citation $\rightarrow$ deux-points + guillemets ouvrants : \zh{老师问：“}
\item Citation = question fermée avec \zh{什么} $\rightarrow$ point d'interrogation \textbf{à l'intérieur} des guillemets : \zh{你们昨天做了什么？”}
\item \zh{小王说} annonce une nouvelle citation $\rightarrow$ \zh{小王说：“}
\item \zh{苹果 香蕉} = noms énumérés $\rightarrow$ virgule énumérative \zh{、}, \zh{和} avant dernier : \zh{我买了苹果、香蕉和橙子。} (Point final à l'intérieur des guillemets).
\item Fermeture guillemets. \zh{然后} (ensuite) introduit phrase suivante $\rightarrow$ virgule \zh{，} après \zh{然后} : \zh{”然后，}
\item Structure \zh{一边……一边……} $\rightarrow$ virgule ordinaire \zh{，} entre les deux actions : \zh{我一边吃水果，一边看电视。}
\item Exclamation finale $\rightarrow$ \zh{真舒服！}
\end{enumerate}
\textbf{Texte ponctué :}
\zh{老师问：“你们昨天做了什么？”小王说：“我买了苹果、香蕉和橙子。”然后，我一边吃水果，一边看电视。真舒服！}\\
\emph{Lǎoshī wèn: “Nǐmen zuótiān zuò le shénme?” Xiǎo Wáng shuō: “Wǒ mǎi le píngguǒ、xiāngjiāo hé chéngzi.” Ránhòu, wǒ yìbiān chī shuǐguǒ, yìbiān kàn diànshì. Zhēn shūfu!}\\
« Le professeur demande : “Qu'avez-vous fait hier ?” Xiao Wang dit : “J'ai acheté des pommes, des bananes et des oranges.” Ensuite, tout en mangeant des fruits, j'ai regardé la télévision. Quel confort ! »
\end{exoresolu}

\courssub{Entraînement intégré : structure + ponctuation}

\begin{exercice}[Exercice 1 : Construire et ponctuer]
Complétez les phrases avec \zh{一边……一边……} et ajoutez la ponctuation correcte (virgules, point final).
\begin{enumerate}
\item \zh{妹妹 \_\_\_\_ 唱歌 \_\_\_\_ \_\_\_\_ 跳舞}
\item \zh{爸爸 \_\_\_\_ 看报纸 \_\_\_\_ \_\_\_\_ 喝茶}
\item \zh{我们 \_\_\_\_ 走路 \_\_\_\_ \_\_\_\_ 聊天}
\item \zh{学生 \_\_\_\_ 听课 \_\_\_\_ \_\_\_\_ 记笔记}
\end{enumerate}
\end{exercice}

\begin{exercice}[Exercice 2 : Traduction thème (Français $\rightarrow$ Chinois)]
Traduisez en chinois (caractères + pinyin) en respectant la structure \zh{一边……一边……} et la ponctuation.
\begin{enumerate}
\item Ma sœur mange tout en regardant la télévision.
\item Le professeur parle tout en écrivant au tableau.
\item Nous marchons tout en chantant.
\item Xiao Li conduit tout en écoutant la radio. (Conduire = \zh{开车} \emph{kāichē})
\end{enumerate}
\end{exercice}

\begin{exercice}[Exercice 3 : Ponctuation complète]
Ponctuez les phrases suivantes avec les signes chinois corrects (\zh{。 ， 、 ； ： “ ” 《》 ？ ！ …… ——}).
\begin{enumerate}
\item \zh{妈妈问 你晚上回家吃饭吗}
\item \zh{超市里有牛奶 面包 鸡蛋和水果}
\item \zh{他一边听音乐 一边做作业}
\item \zh{我在读 人民日报 这是一份很有名的报纸}
\item \zh{中国的城市真多 北京 上海 广州}
\item \zh{他说 我最喜欢的书是 红楼梦}
\item \zh{请坐 我们马上开始}
\end{enumerate}
\end{exercice}

\begin{corrige}[Exercices 1, 2 et 3]

\textbf{Exercice 1 :}
\begin{enumerate}
\item \zh{妹妹一边唱歌，一边跳舞。} (\emph{Mèimei yìbiān chànggē, yìbiān tiàowǔ.})
\item \zh{爸爸一边看报纸，一边喝茶。} (\emph{Bàba yìbiān kàn bàozhǐ, yìbiān hē chá.})
\item \zh{我们一边走路，一边聊天。} (\emph{Wǒmen yìbiān zǒulù, yìbiān liáotiān.})
\item \zh{学生一边听课，一边记笔记。} (\emph{Xuésheng yìbiān tīng kè, yìbiān jì bǐjì.})
\end{enumerate}

\textbf{Exercice 2 :}
\begin{enumerate}
\item \zh{我妹妹一边吃饭，一边看电视。} (\emph{Wǒ mèimei yìbiān chīfàn, yìbiān kàn diànshì.})
\item \zh{老师一边说话，一边写黑板。} (\emph{Lǎoshī yìbiān shuō huà, yìbiān xiě hēibǎn.})
\item \zh{我们一边走路，一边唱歌。} (\emph{Wǒmen yìbiān zǒulù, yìbiān chànggē.})
\item \zh{小李一边开车，一边听收音机。} (\emph{Xiǎo Lǐ yìbiān kāichē, yìbiān tīng shōuyīnjī.})
\end{enumerate}

\textbf{Exercice 3 :}
\begin{enumerate}
\item \zh{妈妈问：“你晚上回家吃饭吗？”} (Citation interrogative : \zh{：“……？”})
\item \zh{超市里有牛奶、面包、鸡蛋和水果。} (Énumération de noms : \zh{、})
\item \zh{他一边听音乐，一边做作业。} (Deux propositions \zh{一边……} : \zh{，})
\item \zh{我在读《人民日报》，这是一份很有名的报纸。} (Titre de journal : \zh{《》} ; deux propositions : \zh{，})
\item \zh{中国的城市真多：北京、上海、广州……} (Annonce liste : \zh{：} ; énumération noms : \zh{、} ; inachevée : \zh{……})
\item \zh{他说：“我最喜欢的书是《红楼梦》。”} (Citation + titre d'œuvre : \zh{：《》} )
\item \zh{请坐，我们马上开始。} (Impératif + proposition : \zh{，} ; fin : \zh{。})
\end{enumerate}
\end{corrige}

\begin{attention}[Bilan : Les trois fautes fatales en examen]
\begin{enumerate}
\item \textbf{Point anglais} : Écrire « . » au lieu du cercle \zh{。} = \textbf{Faute de présentation}.
\item \textbf{Espace parasite} : Mettre une espace avant \zh{？}, \zh{！}, \zh{，}, \zh{。}, \zh{：}, \zh{；} = \textbf{Faute de mise en page}.
\item \textbf{Confusion ， / 、} : Utiliser \zh{，} dans une énumération de noms (\zh{苹果，香蕉}) au lieu de \zh{、} (\zh{苹果、香蕉}) = \textbf{Faute de grammaire ponctuelle}.
\end{enumerate}
\emph{Relisez systématiquement votre copie pour traquer ces trois erreurs.}
\end{attention}

\begin{savaistu}[Le petit cercle qui coûte des points]
À l'écrit manuscrit chinois, le point \zh{。} est un petit cercle tracé d'un seul mouvement, dans le coin inférieur gauche de sa case (comme le point final d'un trait \zh{丶} \emph{zhǔ}). Utiliser un point anglais « . » dans une copie d'examen — au Cameroun comme en Chine — est considéré comme une \cle{faute de présentation} et peut coûter des points. De même, les correcteurs vérifient qu'aucun signe ne commence une ligne : c'est l'une des premières règles de mise en page enseignées aux écoliers chinois (dès le CP), au même titre que l'ordre des traits des caractères. Respecter la pleine chasse, c'est respecter la géométrie de l'écriture chinoise.
\end{savaistu}

\begin{aretenir}[L'essentiel de la leçon 40]
\begin{itemize}
\item \textbf{Structure \zh{一边……一边……} :} \zh{Sujet + 一边 + V₁, 一边 + V₂。} Exprime la simultanéité de deux actions duratives du même sujet. \cle{Virgule \zh{，} obligatoire} entre les deux volets. Sujet non répété.
\item \textbf{Ponctuation :} Tous les signes sont à \cle{pleine chasse} (1 case, sauf \zh{……} et \zh{——} = 2 cases). \cle{Jamais d'espace} avant/après. \cle{Jamais en début de ligne}.
\item \textbf{Signes clés à distinguer :}
\begin{itemize}
\item \zh{。} (fin phrase) vs \zh{？} \zh{！}
\item \zh{，} (pause propositions) vs \zh{、} (énumération mots) \textbf{\cle{Différence majeure vs français}}.
\item \zh{：} (annonce citation/liste) + \zh{“ ”} (citation) vs \zh{《》} (titre œuvre).
\item \zh{……} (6 points, suite) vs \zh{——} (double tiret, explication/rupture).
\end{itemize}
\item \textbf{Application :} Dans \zh{一边……，一边……}, la virgule est syntaxique (sépare deux propositions). Dans \zh{苹果、香蕉和梨}, la virgule énumérative est lexicale (sépare noms).
\end{itemize}
\end{aretenir}

\begin{experience}[Activité orale : « Journaliste d'un soir »]
En binôme. L'élève A est journaliste, l'élève B est une personnalité (footballeur \zh{埃托奥}, chanteur, élève modèle).
\begin{enumerate}
\item A pose 3 questions sur la routine de B (utiliser \zh{你平时一边……一边……吗？}).
\item B répond en utilisant \zh{我一边……，一边……} (vrai ou inventé).
\item A rédige un micro-titre pour son article : \zh{《XXX一边……，一边……》} (marques de titre !).
\item Restitution orale devant la classe. Attention à la prononciation des tons (3\textsuperscript{e} ton \zh{一} $\rightarrow$ \zh{yì} devant 4\textsuperscript{e} ton \zh{边}).
\end{enumerate}
\end{experience}

\coursec{Entraînement intégré : structure + ponctuation}

Dans les sections précédentes, nous avons appris à placer correctement les signes de ponctuation chinois : le point 。, la virgule ，, le signe d'énumération 、, le point d'interrogation ？, les deux-points ：, les guillemets “” et les marques de titres 《》. Nous savons aussi construire la structure \zh{一边……一边……} pour exprimer deux actions simultanées. Il est temps de réunir ces deux compétences dans un entraînement complet, du plus guidé au plus libre : c'est exactement le chemin que vous suivrez le jour de l'évaluation.

\begin{popfigure}
\begin{tikzpicture}[node distance=0.35cm, every node/.style={font=\small}]
\node[draw=popblue, fill=popblueL, rounded corners, minimum width=2.6cm, minimum height=0.9cm, align=center] (e1) {1. Complétion};
\node[draw=popgreen, fill=popgreenL, rounded corners, minimum width=2.6cm, minimum height=0.9cm, align=center, right=of e1] (e2) {2. Transformation};
\node[draw=poporange, fill=poporangeL, rounded corners, minimum width=2.6cm, minimum height=0.9cm, align=center, right=of e2] (e3) {3. Correction};
\node[draw=poppurple, fill=poppurpleL, rounded corners, minimum width=2.6cm, minimum height=0.9cm, align=center, right=of e3] (e4) {4. Ponctuation};
\node[draw=poppink, fill=poppinkL, rounded corners, minimum width=2.6cm, minimum height=0.9cm, align=center, below=0.5cm of e3] (e5) {5. Traduction};
\node[draw=popgold, fill=popgoldL, rounded corners, minimum width=2.6cm, minimum height=0.9cm, align=center, right=of e5] (e6) {6. Production libre};
\draw[-{Stealth}, thick, popblue] (e1) -- (e2);
\draw[-{Stealth}, thick, popgreen] (e2) -- (e3);
\draw[-{Stealth}, thick, poporange] (e3) -- (e4);
\draw[-{Stealth}, thick, poppurple] (e4.south) |- ++(0,-0.25) -| (e5.north);
\draw[-{Stealth}, thick, poppink] (e5) -- (e6);
\end{tikzpicture}
\legende{Progression de l'entraînement : de la phrase à trous au texte personnel.}
\end{popfigure}

\courssub{Rappel éclair avant de commencer}

\begin{aretenir}[Les deux outils du jour]
\begin{itemize}
\item Structure de simultanéité : \cle{Sujet + 一边 + V1, 一边 + V2}. Le sujet est placé \textbf{avant} le premier \zh{一边} et n'est pas répété. Les deux actions doivent être \textbf{compatibles} (on peut réellement les faire en même temps).
\item Ponctuation : chaque signe chinois occupe \textbf{une case entière} (un carré d'écriture). On distingue bien \zh{、} (énumération de mots ou de groupes de mots) et \zh{，} (pause entre propositions) ; les titres de journaux, de livres et d'émissions prennent les marques de titres \zh{《》}.
\end{itemize}
\end{aretenir}

\begin{exemplebox}[Modèle de référence]
\zh{姐姐一边听广播，一边做运动。}\\
\emph{Jiějie yìbiān tīng guǎngbō, yìbiān zuò yùndòng.}\\
« Ma grande sœur écoute la radio tout en faisant du sport. »\\[4pt]
Observez : un seul sujet (\zh{姐姐}), deux \zh{一边}, une virgule \zh{，} entre les deux propositions, un point \zh{。} final.
\end{exemplebox}

\begin{attention}[La prononciation de 一 dans 一边]
Le caractère \zh{一} seul se prononce \emph{yī} (premier ton), mais devant un premier, deuxième ou troisième ton, il passe au quatrième ton : on dit donc \emph{yìbiān}. C'est une règle de changement de ton (\emph{biàndiào}) : l'écriture ne change pas, seule la prononciation change.
\end{attention}

\courssub{Exercice 1 — Complétion}

\begin{methode}[Comment compléter une phrase à trous avec 一边……一边……]
\begin{enumerate}
\item Repérez le \textbf{sujet} de la phrase : il doit être unique et placé avant le premier \zh{一边}.
\item Identifiez les \textbf{deux actions} : vérifiez qu'elles sont simultanées possibles.
\item Insérez \zh{一边} devant chaque verbe, sans répéter le sujet.
\item Relisez la phrase entière : ordre des mots, virgule \zh{，} au milieu, point \zh{。} à la fin.
\end{enumerate}
\end{methode}

\begin{exercice}[Complétez avec 一边……一边……]
\begin{enumerate}
\item 弟弟\underline{\hspace{1.2cm}}看电视，\underline{\hspace{1.2cm}}吃水果。(dìdi / kàn diànshì / chī shuǐguǒ)
\item 我们\underline{\hspace{1.2cm}}听新闻，\underline{\hspace{1.2cm}}做作业。(wǒmen / tīng xīnwén / zuò zuòyè)
\item 妈妈\underline{\hspace{1.2cm}}打电话，\underline{\hspace{1.2cm}}做饭。(māma / dǎ diànhuà / zuò fàn)
\item 学生们\underline{\hspace{1.2cm}}看报纸，\underline{\hspace{1.2cm}}说汉语。(xuéshengmen / kàn bàozhǐ / shuō Hànyǔ)
\end{enumerate}
\end{exercice}

\begin{corrige}[Exercice 1]
\begin{enumerate}
\item 弟弟一边看电视，一边吃水果。\emph{Dìdi yìbiān kàn diànshì, yìbiān chī shuǐguǒ.} « Le petit frère regarde la télévision tout en mangeant des fruits. »
\item 我们一边听新闻，一边做作业。\emph{Wǒmen yìbiān tīng xīnwén, yìbiān zuò zuòyè.} « Nous écoutons les informations tout en faisant nos devoirs. »
\item 妈妈一边打电话，一边做饭。\emph{Māma yìbiān dǎ diànhuà, yìbiān zuò fàn.} « Maman téléphone tout en préparant le repas. »
\item 学生们一边看报纸，一边说汉语。\emph{Xuéshengmen yìbiān kàn bàozhǐ, yìbiān shuō Hànyǔ.} « Les élèves lisent le journal tout en parlant chinois. »
\end{enumerate}
\end{corrige}

\courssub{Exercice 2 — Transformation : fusionner deux phrases}

\begin{methode}[Transformer deux phrases simples en une phrase à actions simultanées]
\begin{enumerate}
\item Vérifiez que le \textbf{sujet est identique} dans les deux phrases. Si les sujets sont différents, la fusion est impossible à ce niveau.
\item Vérifiez que les actions sont \textbf{simultanées possibles} (regarder + boire : oui ; dormir + courir : non).
\item Appliquez le schéma : \cle{Sujet + 一边 + V1, 一边 + V2 + 。} — on supprime le second sujet et le premier point devient une virgule.
\end{enumerate}
\end{methode}

\begin{exoresolu}[Fusion de deux phrases]
Énoncé : fusionnez en une seule phrase avec \zh{一边……一边……} : 他看新闻。他喝茶。\emph{Tā kàn xīnwén. Tā hē chá.}
\tcblower
Solution détaillée :
\begin{enumerate}
\item Sujet identique ? Oui : \zh{他} dans les deux phrases.
\item Actions compatibles ? Oui : on peut regarder les informations et boire du thé en même temps.
\item Application du schéma : on garde le sujet une seule fois, on place \zh{一边} devant chaque verbe, on remplace le premier point par une virgule \zh{，}.
\end{enumerate}
Résultat : \zh{他一边看新闻，一边喝茶。}\emph{Tā yìbiān kàn xīnwén, yìbiān hē chá.} « Il regarde les informations tout en buvant du thé. »
\end{exoresolu}

\begin{exercice}[Fusionnez les phrases suivantes]
\begin{enumerate}
\item 我听广播。我吃饭。(wǒ / tīng guǎngbō / chī fàn)
\item 她上网。她唱歌。(tā / shàngwǎng / chàng gē)
\item 爸爸看《新闻联播》。爸爸喝茶。(bàba / kàn / hē chá)
\item 同学们走路。同学们说话。(tóngxuémen / zǒu lù / shuō huà)
\end{enumerate}
\end{exercice}

\begin{corrige}[Exercice 2]
\begin{enumerate}
\item 我一边听广播，一边吃饭。\emph{Wǒ yìbiān tīng guǎngbō, yìbiān chī fàn.} « J'écoute la radio tout en mangeant. »
\item 她一边上网，一边唱歌。\emph{Tā yìbiān shàngwǎng, yìbiān chàng gē.} « Elle surfe sur Internet tout en chantant. »
\item 爸爸一边看《新闻联播》，一边喝茶。\emph{Bàba yìbiān kàn 《Xīnwén Liánbō》, yìbiān hē chá.} « Papa regarde le Journal télévisé tout en buvant du thé. » — notez les marques de titre \zh{《》} autour du nom de l'émission.
\item 同学们一边走路，一边说话。\emph{Tóngxuémen yìbiān zǒu lù, yìbiān shuō huà.} « Les camarades marchent tout en parlant. »
\end{enumerate}
\end{corrige}

\courssub{Exercice 3 — Correction de fautes}

\begin{attention}[Les trois pièges à traquer]
\begin{itemize}
\item \textbf{Second 一边 oublié} : *\zh{他一边看新闻，喝茶。} est fautif ; il faut deux \zh{一边}.
\item \textbf{Sujet mal placé ou répété} : *\zh{一边他看新闻，一边他喝茶。} — le sujet se met une seule fois, avant le premier \zh{一边}.
\item \textbf{Actions incompatibles} : *\zh{他一边睡觉，一边跑步。} — on ne peut pas dormir et courir en même temps ; la phrase est grammaticale mais absurde.
\end{itemize}
\end{attention}

\begin{exercice}[Corrigez les phrases fautives]
\begin{enumerate}
\item 妹妹一边听音乐，写作业。(mèimei / tīng yīnyuè / xiě zuòyè)
\item 一边老师看报纸，一边老师喝茶。(lǎoshī / kàn bàozhǐ / hē chá)
\item 我一边睡觉，一边看电视。(wǒ / shuìjiào / kàn diànshì)
\item 他们一边，一边看足球比赛。(tāmen / kàn zúqiú bǐsài)
\end{enumerate}
\end{exercice}

\begin{corrige}[Exercice 3]
\begin{enumerate}
\item Faute : second \zh{一边} manquant. Correction : \zh{妹妹一边听音乐，一边写作业。}\emph{Mèimei yìbiān tīng yīnyuè, yìbiān xiě zuòyè.} « La petite sœur écoute de la musique tout en écrivant ses devoirs. »
\item Faute : sujet répété et mal placé. Correction : \zh{老师一边看报纸，一边喝茶。}\emph{Lǎoshī yìbiān kàn bàozhǐ, yìbiān hē chá.} « Le professeur lit le journal tout en buvant du thé. »
\item Faute : actions incompatibles (dormir et regarder la télévision). Correction possible : \zh{我一边吃饭，一边看电视。}\emph{Wǒ yìbiān chī fàn, yìbiān kàn diànshì.} « Je regarde la télévision tout en mangeant. »
\item Faute : premier \zh{一边} sans verbe. Correction : \zh{他们一边吃饭，一边看足球比赛。}\emph{Tāmen yìbiān chī fàn, yìbiān kàn zúqiú bǐsài.} « Ils regardent le match de football tout en mangeant. »
\end{enumerate}
\end{corrige}

\courssub{Exercice 4 — Ponctuation d'un mini-texte}

\begin{attention}[Règles d'or de la ponctuation]
Dans cet exercice, chaque signe doit occuper \textbf{une case} (un carré d'écriture). Vérifiez en particulier : \zh{、} pour les énumérations de mots, \zh{《》} pour les titres de journaux ou d'émissions, \zh{，} entre les deux propositions de \zh{一边……一边……}, et \zh{。} à la fin de chaque phrase complète.
\end{attention}

\begin{textebox}[Texte à ponctuer]
晚上 我们一家人一边看电视 一边吃饭 爸爸喜欢看 新闻联播 妈妈喜欢看电视剧 我呢 我喜欢一边上网 一边听音乐\\[4pt]
\emph{Wǎnshang wǒmen yì jiā rén yìbiān kàn diànshì yìbiān chī fàn bàba xǐhuan kàn Xīnwén Liánbō māma xǐhuan kàn diànshìjù wǒ ne wǒ xǐhuan yìbiān shàngwǎng yìbiān tīng yīnyuè}
\end{textebox}

\begin{exoresolu}[Ponctuer le texte]
Énoncé : ponctuez le texte ci-dessus avec les signes \zh{。，、《》}.
\tcblower
Solution détaillée :
\begin{enumerate}
\item Première phrase : deux propositions liées par \zh{一边……一边……} → virgule \zh{，} au milieu, point \zh{。} à la fin.
\item \zh{新闻联播} est le titre d'une émission → marques de titre \zh{《》}.
\item \zh{我呢} introduit une petite pause → virgule \zh{，}.
\end{enumerate}
Texte ponctué :\\
\zh{晚上，我们一家人一边看电视，一边吃饭。爸爸喜欢看《新闻联播》，妈妈喜欢看电视剧。我呢，我喜欢一边上网，一边听音乐。}\\
\emph{Wǎnshang, wǒmen yì jiā rén yìbiān kàn diànshì, yìbiān chī fàn. Bàba xǐhuan kàn 《Xīnwén Liánbō》, māma xǐhuan kàn diànshìjù. Wǒ ne, wǒ xǐhuan yìbiān shàngwǎng, yìbiān tīng yīnyuè.}\\
« Le soir, toute notre famille regarde la télévision tout en dînant. Papa aime regarder le Journal télévisé, maman aime regarder les séries. Moi, j'aime surfer sur Internet tout en écoutant de la musique. »
\end{exoresolu}

\courssub{Exercice 5 — Traduction thème (français → chinois)}

\begin{exercice}[Traduisez en chinois]
\begin{enumerate}
\item Le soir, mon père écoute la radio tout en lisant le journal.
\item Les élèves de la classe de Première A aiment regarder des vidéos tout en apprenant le chinois.
\item Pendant le match de football, mon grand frère mange des arachides tout en regardant la télévision.
\end{enumerate}
\end{exercice}

\begin{corrige}[Exercice 5]
\begin{enumerate}
\item 晚上，我爸爸一边听广播，一边看报纸。\emph{Wǎnshang, wǒ bàba yìbiān tīng guǎngbō, yìbiān kàn bàozhǐ.}
\item 高二A班的学生们一边看视频，一边学汉语。\emph{Gāo'èr A bān de xuéshengmen yìbiān kàn shìpín, yìbiān xué Hànyǔ.} — en Chine, la classe de Première correspond à la deuxième année de lycée (\zh{高二}).
\item 看足球比赛的时候，我哥哥一边吃花生，一边看电视。\emph{Kàn zúqiú bǐsài de shíhou, wǒ gēge yìbiān chī huāshēng, yìbiān kàn diànshì.}
\end{enumerate}
Justification : dans les trois cas, le sujet est unique, les actions sont compatibles, et le schéma \cle{Sujet + 一边 + V1, 一边 + V2} s'applique directement. Le français « tout en + gérondif » se rend toujours par \zh{一边……一边……}, jamais par une seule occurrence de \zh{一边}.
\end{corrige}

\courssub{Exercice 6 — Production guidée : « Une soirée médias dans ma famille »}

\begin{methode}[Rédiger un court paragraphe avec 一边……一边……]
\begin{enumerate}
\item Annoncez le cadre : quand (\zh{晚上}), qui (\zh{我们一家人}), où (\zh{在客厅里}).
\item Décrivez les actions de chacun : au moins \textbf{deux} phrases avec \zh{一边……一边……}.
\item Soignez la ponctuation : \zh{，} entre les propositions, \zh{。} en fin de phrase, \zh{《》} pour les titres d'émissions, \zh{！} pour l'exclamation finale.
\item Relisez à voix haute : chaque signe = une case, chaque phrase = un sens complet.
\end{enumerate}
\end{methode}

\begin{exemplebox}[Modèle de production — quatre phrases]
\zh{晚上，我们一家人在客厅里。爸爸一边看《新闻联播》，一边喝茶。妈妈一边打电话，一边做饭。我呢，我一边听音乐，一边做作业。我们家真热闹！}\\
\emph{Wǎnshang, wǒmen yì jiā rén zài kètīng li. Bàba yìbiān kàn 《Xīnwén Liánbō》, yìbiān hē chá. Māma yìbiān dǎ diànhuà, yìbiān zuò fàn. Wǒ ne, wǒ yìbiān tīng yīnyuè, yìbiān zuò zuòyè. Wǒmen jiā zhēn rènao!}\\
« Le soir, toute notre famille est au salon. Papa regarde le Journal télévisé tout en buvant du thé. Maman téléphone tout en préparant le repas. Moi, j'écoute de la musique tout en faisant mes devoirs. Quelle animation dans notre famille ! »
\end{exemplebox}

\begin{experience}[Production en classe puis correction collective au tableau]
\begin{enumerate}
\item Chaque élève rédige \textbf{quatre phrases} sur le thème « Une soirée médias dans ma famille », avec au moins deux structures \zh{一边……一边……} et une ponctuation irréprochable (10 minutes).
\item Deux ou trois copies sont reproduites au tableau.
\item La classe corrige collectivement en \textbf{justifiant} chaque choix : place du sujet, présence des deux \zh{一边}, compatibilité des actions, choix entre \zh{，} et \zh{、}, usage de \zh{《》}.
\item Le professeur fait verbaliser la règle à chaque correction : « Pourquoi une virgule ici ? Pourquoi des marques de titre là ? »
\end{enumerate}
\end{experience}

\begin{center}
\begin{tabularx}{\linewidth}{|X|X|X|}
\hline
\rowcolor{popblueL} Point vérifié & Critère de réussite & Exemple correct \\
\hline
Sujet unique & placé avant le premier 一边, non répété & 他一边看新闻，一边喝茶。 \\
\hline
Deux 一边 & jamais un seul & 妹妹一边听音乐，一边写作业。 \\
\hline
Actions compatibles & simultanéité réellement possible & 我一边吃饭，一边看电视。 \\
\hline
Virgule ， & entre les deux propositions & 她一边上网，一边听音乐。 \\
\hline
Titres 《》 & journaux, livres et émissions & 《新闻联播》 \\
\hline
Point 。 & fin de chaque phrase complète & 我们一家人在客厅里。 \\
\hline
\end{tabularx}
\end{center}

\begin{savaistu}[Pour aller plus loin : deux sujets différents]
À un niveau plus avancé, \zh{一边……一边……} peut aussi relier les actions de \textbf{deux personnes différentes} ; chaque sujet se place alors après son \zh{一边} : \zh{一边爸爸看报纸，一边妈妈做饭。}\emph{Yìbiān bàba kàn bàozhǐ, yìbiān māma zuò fàn.} « Pendant que papa lit le journal, maman prépare le repas. » En Première, on vous demande seulement de maîtriser le cas du sujet unique — mais reconnaître cette variante à la lecture est un plus.
\end{savaistu}

\begin{aretenir}[Bilan de l'entraînement intégré]
Vous savez désormais : compléter, transformer, corriger, ponctuer, traduire et produire avec \zh{一边……一边……}. La clé de la réussite est toujours la même démarche en trois temps : \textbf{sujet unique} → \textbf{actions compatibles} → \textbf{schéma Sujet + 一边 + V1, 一边 + V2}, avec une ponctuation soignée où chaque signe occupe sa case. Cette rigueur, acquise sur le thème des médias, vous servira dans tous les modules suivants.
\end{aretenir}

\newpage

\coursec{Activité d'intégration}

\coursec{Leçon 40 — Grammaire : 一边……一边…… ; La ponctuation en chinois}

\courssub{Objectifs de la leçon}

À l'issue de cette leçon, tu seras capable de :
\begin{itemize}
\item comprendre et utiliser la structure \cle{一边……一边……} pour exprimer la simultanéité de deux actions effectuées par un même sujet ;
\item identifier et employer correctement les principaux signes de ponctuation chinois : le point \zh{。}, la virgule \zh{，}, la virgule d'énumération \zh{、}, les guillemets \zh{“ ”}, les marques de titre \zh{《》}, le point d'interrogation \zh{？} et le point d'exclamation \zh{！} ;
\item construire des phrases correctes en respectant l'ordre des mots chinois (sujet en tête, double \zh{一边} obligatoire) et la compatibilité sémantique des actions ;
\item éviter les erreurs fréquentes des francophones (traduction littérale de « en même temps », omission du second \zh{一边}, confusion entre \zh{、} et \zh{，}, point français vs point chinois) ;
\item rédiger et présenter un court texte authentique (script de journal télévisé) mobilisant la structure et la ponctuation.
\end{itemize}

\courssub{Observation et découverte de la structure}

\courssub{Situation de communication}

Le club de chinois de ton lycée à Yaoundé prépare une soirée « Médias et langues ». Chaque groupe doit présenter un mini journal télévisé en chinois. Le sujet imposé est le suivant : \emph{le soir dans une famille camerounaise, chacun utilise les médias à sa façon, souvent en faisant deux choses en même temps}. Ton groupe incarne la rédaction de la chaîne et doit décrire cette scène familiale : le père regarde les informations, la mère écoute la radio, la grande sœur utilise son téléphone, le petit frère lit le journal.

Pour vous inspirer, voici le début du script rédigé par la rédaction d'une autre classe :

\begin{textebox}[Extrait du journal télévisé — Classe de Terminale]
晚上好！欢迎收看《家庭新闻》。\\
Wǎnshang hǎo! Huānyíng shōukàn 《Jiātíng xīnwén》.\\
« Bonsoir ! Bienvenue au journal “Nouvelles de la famille”. »\\[4pt]
今天晚上，我们看看一个喀麦隆家庭。\\
Jīntiān wǎnshang, wǒmen kànkan yí ge Kāmàilóng jiātíng.\\
« Ce soir, nous observons une famille camerounaise. »\\[4pt]
爸爸一边看电视，一边喝茶。\\
Bàba yìbiān kàn diànshì, yìbiān hē chá.\\
« Le papa regarde la télévision tout en buvant du thé. »
\end{textebox}

\courssub{Analyse guidée}

Observe l'extrait ci-dessus et réponds aux questions :
\begin{enumerate}
\item Repère la structure qui exprime « faire deux choses en même temps ». Combien de fois le mot \zh{一边} apparaît-il dans la phrase 3 ?
\item Identifie le sujet de la phrase 3. Est-il répété devant le second \zh{一边} ?
\item Observe la ponctuation : quel signe sépare les deux actions ? Quel signe termine la phrase ? Quel signe entoure le titre de l'émission ?
\item Dans la phrase 1, quel signe est utilisé pour l'exclamation ? Et pour la salutation ?
\end{enumerate}

\begin{aretenir}[Découverte : la structure \cle{一边……一边……}]
\begin{itemize}
\item \textbf{Forme :} \zh{一边} (yībiān) signifie « d'un côté / pendant ce temps ». La structure répète \zh{一边} deux fois.
\item \textbf{Sens :} Elle exprime que \textbf{deux actions se déroulent simultanément}, effectuées par \textbf{un même sujet}.
\item \textbf{Exemple observé :} \zh{爸爸一边看电视，一边喝茶。} (Le papa regarde la télé \textbf{et en même temps} boit du thé.)
\end{itemize}
\end{aretenir}

\courssub{Schéma et construction de la structure}

\begin{propriete}[Schéma grammatical de \zh{一边……一边……}]
\textbf{Structure de base :}\\
\zh{Sujet + 一边 + Verbe 1 (+ Complément 1)，一边 + Verbe 2 (+ Complément 2) 。}

\begin{center}
\begin{tabularx}{\linewidth}{|X|X|X|}
\hline
\rowcolor{popblueL} \textbf{Composant} & \textbf{Exemple} & \textbf{Rôle} \\
\hline
Sujet & \zh{爸爸} (Bàba) & Celui qui fait les deux actions \\
\hline
\zh{一边} (1er) & \zh{一边} & Marque le début de la 1ère action \\
\hline
Verbe 1 + Compl. & \zh{看电视} (kàn diànshì) & Action principale / plus longue \\
\hline
\zh{，} & \zh{，} & Pause obligatoire entre les deux actions \\
\hline
\zh{一边} (2e) & \zh{一边} & Marque le début de la 2e action (obligatoire) \\
\hline
Verbe 2 + Compl. & \zh{喝茶} (hē chá) & Action secondaire / accompagnement \\
\hline
\zh{。} & \zh{。} & Point final chinois (pleine chasse) \\
\hline
\end{tabularx}
\end{center}

\textbf{Règles de construction :}
\begin{enumerate}
\item \textbf{Sujet unique :} Il ne se répète \textbf{pas} devant le second \zh{一边}. \zh{爸爸一边看电视，爸爸一边喝茶} $\rightarrow$ \textbf{Incorrect}.
\item \textbf{Double \zh{一边} obligatoire :} On ne peut pas dire \zh{*爸爸看电视，一边喝茶} ni \zh{*爸爸一边看电视，喝茶}.
\item \textbf{Ordre des actions :} L'action principale (souvent plus longue) vient généralement en premier, l'action d'accompagnement en second.
\item \textbf{Compatibilité :} Les deux actions doivent pouvoir se faire physiquement en même temps (regarder + manger \checkmark ; dormir + courir \cross).
\end{enumerate}
\end{propriete}

\begin{exemplebox}[Exemples variés avec pinyin et traduction]
\begin{itemize}
\item \zh{我一边吃早饭，一边听新闻。} \\
      \emph{Wǒ yìbiān chī zǎofàn, yìbiān tīng xīnwén.} \\
      « Je prends mon petit-déjeuner tout en écoutant les informations. »
\item \zh{妹妹一边做作业，一边听音乐。} \\
      \emph{Mèimei yìbiān zuò zuòyè, yìbiān tīng yīnyuè.} \\
      « Ma petite sœur fait ses devoirs tout en écoutant de la musique. »
\item \zh{我们一边走路，一边聊天。} \\
      \emph{Wǒmen yìbiān zǒulù, yìbiān liáotiān.} \\
      « Nous marchons tout en discutant. »
\item \zh{老师一边写字，一边说话。} \\
      \emph{Lǎoshī yìbiān xiě zì, yìbiān shuōhuà.} \\
      « Le professeur écrit au tableau tout en parlant. »
\end{itemize}
\end{exemplebox}

\courssub{Emplois, limites et exceptions}

\begin{propriete}[Emplois principaux et limites]
\textbf{A. Actions habituelles / descriptives (le plus courant en HSK 1-3) :}\\
Décrit une scène, une habitude ou un fait présent/général.
\zh{他常常一边抽烟，一边开车。} (Dangereux ! Mais structure correcte.)

\textbf{B. Actions en cours (avec 正在 / 在) :}\\
On peut ajouter \zh{在} ou \zh{正在} \textbf{avant le sujet} ou \textbf{devant le premier 一边} pour insister sur le processus en cours.
\zh{看！弟弟正在一边看动画片，一边吃零食。} \\
\emph{Kàn! Dìdi zhèngzài yìbiān kàn dònghuàpiàn, yìbiān chī língshí.} \\
« Regarde ! Le petit frère est en train de regarder des dessins animés tout en mangeant des grignotines. »

\textbf{C. Limite : Pas d'action ponctuelle / instantanée.}\\
Les verbes d'état momentané (arriver, mourir, casser, trouver) ne s'utilisent pas avec \zh{一边}.
\zh{*他一边到北京，一边给我打电话。} $\rightarrow$ \textbf{Incorrect} (\zh{到} est ponctuel).
\emph{Correction :} \zh{他到了北京，给我打电话。} (Arrivé à Pékin, il m'a appelé.)

\textbf{D. Limite : Sujets différents = impossible.}\\
\zh{*爸爸一边看电视，妈妈一边喝茶。} $\rightarrow$ \textbf{Incorrect} pour cette structure.
\emph{Alternative :} \zh{爸爸看电视，妈妈喝茶。} (Deux phrases coordonnées) ou \zh{爸爸看电视的时候，妈妈在喝茶。}
\end{propriete}

\courssub{Comparaison avec le français et pièges des francophones}

\begin{attention}[Erreurs fréquentes des apprenants francophones — \zh{一边……一边……}]
\begin{itemize}
\item \textbf{Piège 1 : Traduction littérale de « en même temps » / « tout en ».}\\
      \emph{Français :} « Il mange en regardant la télé. » / « Il mange en même temps qu'il regarde la télé. »\\
      \emph{Chinois incorrect :} \zh{*他吃饭的时候看电视。} / \zh{*他吃饭，同时看电视。}\\
      \emph{Chinois correct :} \zh{他一边吃饭，一边看电视。} $\leftarrow$ \textbf{Structure figée, double \zh{一边} obligatoire.}
\item \textbf{Piège 2 : Omission du second \zh{一边}.}\\
      \emph{Incorrect :} \zh{*妈妈一边做饭，洗菜。} (Manque le 2e \zh{一边}).\\
      \emph{Correct :} \zh{妈妈一边做饭，一边洗菜。} ou \zh{妈妈一边洗菜，一边做饭。} (Ordre logique : laver puis cuisiner).
\item \textbf{Piège 3 : Répétition du sujet.}\\
      \emph{Incorrect :} \zh{*哥哥一边打球，哥哥一边喝水。}\\
      \emph{Correct :} \zh{哥哥一边打球，一边喝水。}
\item \textbf{Piège 4 : Confusion avec \zh{一面……一面……}.}\\
      \zh{一面……一面……} (yīmiàn... yīmiàn...) est plus formel, souvent écrit, et peut exprimer une contradiction (« d'un côté... de l'autre... »). \zh{一边……一边……} est neutre, oral et écrit, pour actions physiques simultanées. \textbf{Au programme LV2 Première : maîtriser \zh{一边……一边……}.}
\item \textbf{Piège 5 : Prononciation de \zh{一} (yī) dans \zh{一边}.}\\
      Règle de changement de ton (sandhi) : \zh{一} (1er ton) $\rightarrow$ \zh{二} (2e ton / yí) devant un 4e ton (\zh{biān} est 1er ton, mais \zh{边} seul est 1er ton ; cependant \zh{一边} se dit souvent \emph{yíbiān} par assimilation ou \emph{yībiān} en citation). \textbf{En pratique orale : prononce \emph{yíbiān} [yǐbiān] naturellement.}
\end{itemize}
\end{attention}

\courssub{La ponctuation en chinois}

\begin{propriete}[Tableau récapitulatif des signes de ponctuation chinois (HSK 1-4)]
\begin{center}
\begin{tabularx}{\linewidth}{|X|X|X|X|}
\hline
\rowcolor{popblueL} \textbf{Signe} \& \textbf{Nom chinois} \& \textbf{Usage principal} \& \textbf{Exemple} \\
\hline
\zh{。} \& 句号 (jùhào) \& Point final (déclaratif). Petit cercle centré. \& \zh{我爱喀麦隆。} \\
\hline
\zh{，} \& 逗号 (dòuhào) \& Virgule : pause à l'intérieur de la phrase (sujet, compléments, clauses). \& \zh{今天晚上，我们看电视。} \\
\hline
\zh{、} \& 顿号 (dùnhào) \& Virgule d'énumération : \textbf{uniquement} entre mots/groupes de \textbf{même nature} (noms, verbes, adjectifs). Pas de \zh{、} avant le dernier si \zh{和} / \zh{或} présent. \& \zh{电视、广播、报纸和手机} \\
\hline
\zh{“ ”} \& 引号 (yǐnhào) \& Guillemets : citations directes, titres d'articles/poèmes courts, emphase. \& \zh{他说：“我很忙。”} \\
\hline
\zh{《》} \& 书名号 (shūmínghào) \& Marques de titre : livres, journaux, films, émissions, journaux, lois, navires. \& \zh{《人民日报》} \zh{《家庭新闻》} \\
\hline
\zh{？} \& 问号 (wènhào) \& Point d'interrogation (phrase interrogative). \& \zh{你好吗？} \\
\hline
\zh{！} \& 感叹号 (gǎntànhào) \& Point d'exclamation (exclamation, ordre, salutations). \& \zh{晚上好！} \\
\hline
\zh{：} \& 冒号 (màohào) \& Deux-points : annonce une citation, une liste, une explication. \& \zh{爸爸说：“回家吃饭。”} \\
\hline
\zh{；} \& 分号 (fēnhào) \& Point-virgule : sépare des clauses longues et parallèles dans une phrase complexe. \& \zh{天很热；大家都去游泳。} \\
\hline
\end{tabularx}
\end{center}

\end{propriete}

\begin{aretenir}[Règles d'or de la ponctuation chinoise]
\begin{enumerate}
\item \textbf{Pleine chasse :} Tous les signes occupent \textbf{un carré plein} (largeur d'un caractère). Pas d'espace avant/après dans le flux de texte.
\item \textbf{Point final :} C'est un \textbf{cercle} \zh{。}, jamais un point bas \zh{.}.
\item \textbf{Virgule \zh{，} vs \zh{、} :} \zh{，} structure la phrase (sujet, temps, lieu, propositions). \zh{、} énumère des \textbf{éléments de même niveau} (liste de noms, de verbes). \textbf{Jamais} \zh{、} entre deux phrases complètes.
\item \textbf{Citation :} \zh{说：} (dire :) + \zh{“ } + contenu + \zh{ ”}. Le point final de la citation \zh{。/！/？} se place \textbf{à l'intérieur} des guillemets.
\item \textbf{Titres :} \textbf{Jamais} de guillemets français « » ni guillemets droits "" pour les titres d'œuvres longues. Toujours \zh{《》}.
\end{enumerate}
\end{aretenir}

\begin{exemplebox}[Illustration de la ponctuation dans un paragraphe]
\zh{今天早上，爸爸买了报纸、杂志和牛奶。他一边喝牛奶，一边看《喀麦隆论坛报》。突然，他问道：“这篇文章是谁写的？”妈妈回答：“是张伟写的！”} \\
\emph{Jīntiān zǎoshang, bàba mǎile bàozhǐ、zázhì hé niúnǎi. Tā yìbiān hē niúnǎi, yìbiān kàn 《Kāmàilóng Lùntán Bào》。Tūrán, tā wèn dào:“Zhè piān wénzhāng shì shéi xiě de?” Māma huídá:“Shì Zhāng Wěi xiě de!”} \\
« Ce matin, papa a acheté le journal, le magazine et le lait. Il buvait son lait tout en lisant le "Forum du Cameroun". Soudain, il a demandé : "Qui a écrit cet article ?" Maman a répondu : "C'est Zhang Wei qui l'a écrit !" »
\end{exemplebox}

\courssub{Entraînement intégré : Structure + Ponctuation}

\courssub{Mise en situation : « Le Journal de la Famille »}

\begin{methode}[Déroulement de l'activité de groupe (50 min)]
\begin{enumerate}
\item \textbf{Brainstorming lexical (5 min).} En groupe de 4, listez au brouillon au moins 10 mots utiles : membres de la famille (\zh{爸爸、妈妈、哥哥、姐姐、弟弟、妹妹}), médias (\zh{电视、广播、报纸、手机、电脑、网络}), actions (\zh{看、听、读、写、玩、吃、喝、做、打、聊}). Vérifiez pinyin et tons.
\item \textbf{Rédaction du script (20 min).} Rédigez un script de \textbf{6 à 8 phrases} respectant \textbf{toutes} les contraintes :
      \begin{itemize}
      \item Au moins \textbf{4 phrases} avec \zh{一边……一边……} (couvrir les 4 membres de la famille).
      \item Au moins \textbf{une énumération} de 3+ médias avec \zh{、}.
      \item Un \textbf{titre d'émission} entre \zh{《》}.
      \item Au moins \textbf{une citation directe} (parole d'un membre) entre \zh{“ ”} introduite par \zh{说：}.
      \item Ponctuation irréprochable : \zh{。} (fin), \zh{，} (pauses), \zh{、} (liste), \zh{？/！} si needed.
      \item Contexte camerounais bienvenu (ndolé, football, CRTV, Mt. Fébé, etc.).
      \end{itemize}
\item \textbf{Relecture croisée (5 min).} Échangez avec le groupe voisin. Check-list : \ding{51} Sujet unique ? \ding{51} Double \zh{一边} ? \ding{51} Virgule \zh{，} entre les deux \zh{一边} ? \ding{51} \zh{、} seulement pour liste ? \ding{51} \zh{《》} pour titre ? \ding{51} \zh{“ ”} pour citation ? \ding{51} Point \zh{。} final ? Corrigez au stylo vert.
\item \textbf{Répétition orale (5 min).} Répartissez les rôles (1 présentateur, 3-4 membres famille). Lisez à voix haute : soignez les tons, le \emph{yíbiān}, la fluidité.
\item \textbf{Présentation \& Écoute active (15 min).} Chaque groupe présente. Auditeurs : complétez la grille d'écoute (voir \texttt{experience} ci-dessous).
\item \textbf{Vote.} 3 critères : Grammaire (structure + ordre), Ponctuation (visuelle), Prononciation (tons + fluidité).
\end{enumerate}
\end{methode}

\courssub{Ressources lexicales mobilisables}

\begin{center}
\begin{tabularx}{\linewidth}{|X|X|X|X|}
\hline
\rowcolor{popblueL} \textbf{Caractères} \& \textbf{Pinyin} \& \textbf{Nature} \& \textbf{Traduction} \\
\hline
看电视 \& kàn diànshì \& V + Compl. \& regarder la télévision \\
\hline
听广播 \& tīng guǎngbō \& V + Compl. \& écouter la radio \\
\hline
看报纸 \& kàn bàozhǐ \& V + Compl. \& lire le journal \\
\hline
玩手机 \& wán shǒujī \& V + Compl. \& utiliser/jouer au téléphone \\
\hline
上网 \& shàng wǎng \& V + Compl. \& surfer sur Internet \\
\hline
吃晚饭 \& chī wǎnfàn \& V + Compl. \& dîner \\
\hline
喝茶 / 喝咖啡 \& hē chá / hē kāfēi \& V + Compl. \& boire du thé / du café \\
\hline
做饭 / 做菜 \& zuò fàn / zuò cài \& V + Compl. \& cuisiner / préparer le repas \\
\hline
聊天 \& liáo tiān \& V (séparable) \& discuter, bavarder \\
\hline
新闻 \& xīnwén \& Nom \& informations, nouvelles \\
\hline
节目 \& jiémù \& Nom \& émission, programme \\
\hline
有意思 \& yǒu yìsi \& Adj. \& intéressant, amusant \\
\hline
喀麦隆 \& Kāmàilóng \& Nom propre \& Cameroun \\
\hline
雅温得 \& Yǎwēndé \& Nom propre \& Yaoundé \\
\hline
\end{tabularx}
\end{center}

\begin{attention}[Rappel des pièges pour la rédaction]
\begin{itemize}
\item \zh{看报纸} (lire le journal) — \textbf{pas} \zh{读报纸} (moins courant pour « lire le journal » au sens de le consulter).
\item \zh{玩手机} (utiliser le téléphone pour loisir/réseaux) vs \zh{打电话} (passer un appel).
\item \zh{做饭} (cuisiner le repas principal) vs \zh{做菜} (cuisiner des plats).
\item Énumération : \zh{电视、广播、报纸、手机} (4 éléments, 3 \zh{、}). Si \zh{和} avant le dernier : \zh{电视、广播、报纸和手机} (3 \zh{、} + \zh{和}). Les deux sont corrects.
\item Citation : \zh{弟弟说：“我喜欢看足球比赛！”} (Deux-points \zh{：} après \zh{说}, point \zh{！} \textbf{dans} les guillemets).
\end{itemize}
\end{attention}

\courssub{Proposition de production et corrigé commenté}

\begin{exoresolu}[Script modèle — « Journal de la famille Mvondo »]
Rédigez votre script, puis comparez-le à ce modèle.
\tcblower
\textbf{Script modèle (7 phrases) :}\\[4pt]
\zh{晚上好！欢迎收看《喀麦隆家庭观察》。} \\
\emph{Wǎnshang hǎo! Huānyíng shōukàn 《Kāmàilóng Jiātíng Guānchá》.} \\
« Bonsoir ! Bienvenue à l'émission “Observation des familles camerounaises”. »\\[4pt]
\zh{今天晚上，这个家庭用电视、广播、手机和报纸。} \\
\emph{Jīntiān wǎnshang, zhège jiātíng yòng diànshì、guǎngbō、shǒujī hé bàozhǐ.} \\
« Ce soir, cette famille utilise la télé, la radio, le téléphone et le journal. »\\[4pt]
\zh{爸爸一边看足球新闻，一边喝茶。} \\
\emph{Bàba yìbiān kàn zúqiú xīnwén, yìbiān hē chá.} \\
« Papa regarde les infos foot tout en buvant du thé. »\\[4pt]
\zh{妈妈一边听广播，一边做恩多莱。} \\
\emph{Māma yìbiān tīng guǎngbō, yìbiān zuò ēnduólái.} \\
« Maman écoute la radio tout en préparant le ndolé. »\\[4pt]
\zh{哥哥一边玩手机，一边吃芒果。} \\
\emph{Gēge yìbiān wán shǒujī, yìbiān chī mángguǒ.} \\
« Grand frère joue sur son téléphone tout en mangeant une mangue. »\\[4pt]
\zh{姐姐一边上网，一边写作业。} \\
\emph{Jiějie yìbiān shàng wǎng, yìbiān xiě zuòyè.} \\
« Grande sœur surfe sur le net tout en faisant ses devoirs. »\\[4pt]
\zh{弟弟说：“这个节目真有意思！”} \\
\emph{Dìdi shuō:“Zhège jiémù zhēn yǒu yìsi!”} \\
« Petit frère dit : “Cette émission est vraiment intéressante !” »\\[6pt]

\textbf{Commentaire détaillé des choix linguistiques :}
\begin{itemize}
\item \textbf{Structure \zh{一边……一边……} (Phrases 3-6) :} 4 occurrences. Sujets distincts (\zh{爸爸、妈妈、哥哥、姐姐}) mais structure identique. Actions compatibles : regarder/boire, écouter/cuisiner, téléphone/manger, internet/écrire. Ordre logique : action média (longue) + action quotidienne (courte).
\item \textbf{Énumération \zh{、} (Phrase 2) :} 4 médias \zh{电视、广播、手机和报纸}. \zh{、} sépare les 3 premiers, \zh{和} lie le dernier. Nature identique (noms d'objets/médias).
\item \textbf{Titre \zh{《》} (Phrase 1) :} \zh{《喀麦隆家庭观察》} titre d'émission TV. Jamais de guillemets.
\item \textbf{Citation \zh{“ ”} (Phrase 7) :} \zh{弟弟说：} (verbe + deux-points) + \zh{“内容!”}. Ponctuation finale \zh{！} à l'intérieur.
\item \textbf{Ponctuation générale :} \zh{！} salutation. \zh{，} après \zh{今天晚上} (complément de temps en tête). \zh{。} fin de chaque phrase déclarative. \zh{？} absent ici mais prêt à l'emploi.
\item \textbf{Contexte camerounais :} \zh{恩多莱} (ndolé, plat national), \zh{足球新闻} (football roi), \zh{芒果} (fruit local), \zh{喀麦隆} dans le titre.
\end{itemize}
\end{exoresolu}

\begin{experience}[Grille d'écoute active — Présentations des groupes]
Pendant les présentations, complète ce tableau sur ton cahier :

\begin{center}
\begin{tabularx}{\linewidth}{|X|X|X|X|}
\hline
\rowcolor{popblueL} \textbf{Groupe n°} \& \textbf{Phrases \zh{一边……一边……} entendues (Sujet + V1 + V2)} \& \textbf{Ponctuation repérée sur l'affiche (\zh{《》}, \zh{“ ”}, \zh{、})} \& \textbf{Erreur relevée + Correction proposée} \\
\hline
1 \& \& \& \\
\hline
2 \& \& \& \\
\hline
3 \& \& \& \\
\hline
4 \& \& \& \\
\hline
\end{tabularx}
\end{center}

\textbf{Consigne :} Sois bienveillant mais rigoureux. Note exactement ce que tu vois/entends. À la fin, la classe vote : 1 pt Grammaire, 1 pt Ponctuation, 1 pt Prononciation. Total max 3 pts. Le groupe gagnant reçoit le titre de « Meilleure Rédaction Médias » du trimestre.
\end{experience}

\courssub{Bilan de la leçon}

Cette leçon t'a permis de maîtriser deux piliers de l'expression écrite et orale en chinois :
\begin{enumerate}
\item \textbf{La simultanéité :} La structure \cle{一边……一边……} est l'outil standard (HSK 2-3) pour dire « faire A tout en faisant B ». Rappelle-toi : \textbf{Sujet unique}, \textbf{double \zh{一边} obligatoire}, \textbf{virgule \zh{，} au milieu}, \textbf{actions compatibles}. Ne tombe pas dans le piège du français « en + gérondif » ou « en même temps que ».
\item \textbf{La ponctuation :} Elle n'est pas décorative, elle structure la pensée. Le point \zh{。} (cercle), la virgule \zh{，} (structure), la virgule d'énumération \zh{、} (liste homogène), les guillemets \zh{“ ”} (citation), les marques \zh{《》} (titres d'œuvres) sont non négociables à l'écrit. Un texte sans ponctuation chinoise correcte est considéré comme fautif en examen.
\end{enumerate}

\cle{Conseil pour le portfolio :} Conserve ton script corrigé (version verte + version finale). Il constitue une pièce de choix pour ton dossier d'évaluation du Module 5 « Mass médias et télécommunication ». Tu pourras t'en servir comme modèle pour la rédaction d'un article de blog, d'une lettre ou d'une saynète plus longue en Terminale.

\begin{savaistu}[Le saviez-vous ? — Ponctuation et histoire]
La ponctuation chinoise moderne (标点符号 \zh{biāodiǎn fúhào}) a été normalisée au début du XXe siècle (1919-1930) sous l'influence de la ponctuation occidentale (notamment via le japonais). Avant cela, les textes classiques chinois s'écrivaient \textbf{sans ponctuation} (文言文), la respiration et la structure se devinaient au rythme des caractères (souvent 4 ou 7 par vers). Le point \zh{。} et la virgule \zh{，} sont les premiers adoptés. Les marques de titre \zh{《》} viennent des guillemets français « » retournés et agrandis pour occuper un carré plein. Aujourd'hui, Taïwan et Hong Kong utilisent parfois des guillemets « » pour les citations et 《》 pour les titres, tandis que la Chine continentale standardise \zh{“ ”} et \zh{《》}.
\end{savaistu}

\begin{exercice}[Entraînement individuel — Devoirs / Renforcement]
\begin{enumerate}
\item \textbf{Transformez} les phrases françaises en chinois avec \zh{一边……一边……} :
      \begin{enumerate}
      \item Maman cuisine en écoutant la radio.
      \item Le grand frère regarde un match de foot en mangeant des arachides.
      \item Nous (les élèves) lisons le texte en écoutant l'enregistrement.
      \item Le prof écrit au tableau en expliquant la grammaire.
      \end{enumerate}
\item \textbf{Ponctuez} correctement le paragraphe suivant (ajoutez \zh{。，、：“”《》？！}) :
      \zh{今天 我 买了 苹果 香蕉 和 西瓜 妈妈 说 这些 水果 很 新鲜 我 一边 吃 苹果 一边 看 《西游记》 真 有意思}
\item \textbf{Corrigez} les erreurs dans les phrases suivantes :
      \begin{enumerate}
      \item \zh{*妹妹一边跳舞，唱歌。}
      \item \zh{*爸爸看电视，一边喝酒。}
      \item \zh{*他说：“我饿了。}
      \item \zh{*我们用笔、纸、书包和。}
      \item \zh{*欢迎收看“新闻联播”。}
      \end{enumerate}
\item \textbf{Rédigez} (50-80 caractères) : « Décris ton soir typique à la maison : qui fait quoi, avec quels médias, en faisant quoi d'autre ? » Utilise \zh{一边……一边……} 2 fois, une énumération \zh{、}, un titre \zh{《》} (inventé) et une citation \zh{“ ”}.
\end{enumerate}
\end{exercice}


\begin{corrige}[Exercices de renforcement]
\textbf{1. Transformation :}
\begin{enumerate}
\item \zh{妈妈一边做饭，一边听广播。} (\emph{Māma yìbiān zuò fàn, yìbiān tīng guǎngbō.})
\item \zh{哥哥一边看足球比赛，一边吃花生。} (\emph{Gēge yìbiān kàn zúqiú bǐsài, yìbiān chī huāshēng.})
\item \zh{我们一边读课文，一边听录音。} (\emph{Wǒmen yìbiān dú kèwén, yìbiān tīng lùyīn.})
\item \zh{老师一边写黑板，一边讲语法。} (\emph{Lǎoshī yìbiān xiě hēibǎn, yìbiān jiǎng yǔfǎ.})
\end{enumerate}

\textbf{2. Ponctuation :}
\zh{今天，我买了苹果、香蕉和西瓜。妈妈说：“这些水果很新鲜！”我一边吃苹果，一边看《西游记》，真有意思！}

\textbf{3. Correction :}
\begin{enumerate}
\item \zh{妹妹一边跳舞，\textbf{一边}唱歌。} (Manquait 2e \zh{一边})
\item \zh{\textbf{爸爸一边}看电视，一边喝酒。} (Manquait 1er \zh{一边} + sujet)
\item \zh{他说：“我饿了\textbf{。”}} (Point final \zh{。} \textbf{à l'intérieur} des guillemets + guillemets fermants)
\item \zh{我们用笔、纸、书包\textbf{和文具盒}。} (Énumération incomplète : \zh{和} nécessite un élément après lui. Ou retirer \zh{和} : \zh{我们用笔、纸、书包。})
\item \zh{欢迎收看\textbf{《}新闻联播\textbf{》}。} (Titre d'émission $\rightarrow$ \zh{《》}, pas guillemets).
\end{enumerate}

\textbf{4. Production libre (exemple) :}
\zh{晚上好！欢迎收看《我家的晚上》。我家用电视、手机和电脑。爸爸一边看新闻，一边喝啤酒。妈妈一边刷视频，一边聊天。我说：“今晚真开心！”}
\end{corrige}

\newpage

\coursec{Méthodes et exercices résolus}

\coursec{Leçon 40 — Grammaire : 一边……一边…… ; La ponctuation en chinois}

\begin{textebox}[Situation de communication : Une soirée en famille à Yaoundé]
\zh{晚上，我们一家人在客厅。} Wǎnshang, wǒmen yì jiā rén zài kètīng. \\
\zh{爸爸一边看新闻，一边喝茶。} Bàba yíbiān kàn xīnwén, yíbiān hē chá. \\
\zh{妈妈一边做饭，一边唱歌。} Māma yíbiān zuòfàn, yíbiān chànggē. \\
\zh{我一边做作业，一边听音乐。} Wǒ yíbiān zuò zuòyè, yíbiān tīng yīnyuè. \\
\zh{弟弟一边吃水果，一边看动画片。} Dìdi yíbiān chī shuǐguǒ, yíbiān kàn dònghuàpiàn. \\
\zh{我们家真热闹！} Wǒmen jiā zhēn rènao! \\
« Ce soir, toute la famille est au salon. Papa regarde les informations en buvant du thé. Maman cuisine en chantant. Je fais mes devoirs en écoutant de la musique. Mon petit frère mange des fruits en regardant un dessin animé. Comme c'est animé chez nous ! »
\end{textebox}

\begin{propriete}[Schéma de la structure 一边……一边……]
La structure \zh{一边……一边……} (yíbiān……yíbiān……) exprime la \textbf{simultanéité de deux actions} effectuées par \textbf{un même sujet}.

\cle{Sujet + 一边 + Verbe$_1$ (+ Complément$_1$) ，一边 + Verbe$_2$ (+ Complément$_2$) 。}

\begin{center}
\begin{tabularx}{\linewidth}{|X|X|X|}
\hline
\rowcolor{popblueL} Composant & Exemple & Rôle \\
\hline
Sujet & \zh{他} (tā) & Unique, placé en tête \\
\hline
Premier membre & \zh{一边吃饭} (yíbiān chīfàn) & Action 1 + son complément \\
\hline
Séparateur & \zh{，} & Virgule chinoise (pleine chasse) \\
\hline
Deuxième membre & \zh{一边看电视} (yíbiān kàn diànshì) & Action 2 + son complément \\
\hline
Marque de fin & \zh{。} & Point chinois \\
\hline
\end{tabularx}
\end{center}

\textbf{Attention au tone sandhi :} \zh{一} (yī) $\rightarrow$ \zh{yí} (2e ton) devant \zh{边} (biān, 1er ton). On écrit \textbf{yíbiān}, jamais *yìbiān.
\end{propriete}

\begin{definition}[Vocabulaire clé de la leçon (HSK 2--3)]
\begin{center}
\begin{tabularx}{\linewidth}{|X|X|X|X|}
\hline
\rowcolor{popblueL} Caractères & Pinyin & Nature & Traduction \\
\hline
一边……一边…… & yíbiān……yíbiān…… & structure & en même temps que, tout en... \\
\hline
看电视 & kàn diànshì & VO & regarder la télévision \\
\hline
听音乐 & tīng yīnyuè & VO & écouter de la musique \\
\hline
做作业 & zuò zuòyè & VO & faire ses devoirs \\
\hline
做饭 & zuòfàn & V / VO & cuisiner, faire la cuisine \\
\hline
唱歌 & chànggē & V / VO & chanter \\
\hline
走路 & zǒulù & V / VO & marcher \\
\hline
说话 & shuōhuà & V / VO & parler \\
\hline
开车 & kāichē & V / VO & conduire \\
\hline
打电话 & dǎ diànhuà & VO & téléphoner \\
\hline
跑步 & pǎobù & V & courir, faire du jogging \\
\hline
睡觉 & shuìjiào & VO & dormir \\
\hline
看比赛 & kàn bǐsài & VO & regarder un match \\
\hline
吃花生 & chī huāshēng & VO & manger des cacahuètes \\
\hline
喊 / 加油 & hǎn / jiāyóu & V / V & crier / Allez ! (encourager) \\
\hline
进球 & jìnqiú & V / N & marquer un but / but \\
\hline
热闹 & rènao & Adj & animé, vivant \\
\hline
雅温得 & Yǎwēndé & Npr & Yaoundé (capitale du Cameroun) \\
\hline
喀麦隆 & Kāmàilóng & Npr & Cameroun \\
\hline
\end{tabularx}
\end{center}
\end{definition}

\begin{popfigure}
\begin{tikzpicture}[
  node distance=1.2cm and 1.5cm,
  main/.style={draw=popblue, thick, rounded corners, fill=popblueL, text width=3.2cm, align=center, font=\sffamily},
  arrow/.style={-Stealth, thick, popdark},
  labelnode/.style={font=\small\sffamily, popmuted}
]
\node[main, align=center] (sujet){Sujet unique\\ \zh{他 / 我们 / 妈妈}};
\node[main, right=of sujet, align=center] (bian1){\zh{一边} + Verbe$_1$ (+ Compl.)\\ \zh{一边吃饭}};
\node[main, right=of bian1] (comma) {\zh{，}};
\node[main, right=of comma, align=center] (bian2){\zh{一边} + Verbe$_2$ (+ Compl.)\\ \zh{一边看电视}};
\node[main, right=of bian2] (final) {\zh{。}};

\draw[arrow] (sujet) -- (bian1) node[midway, above, labelnode] {fait action 1};
\draw[arrow] (bian1) -- (comma);
\draw[arrow] (comma) -- (bian2) node[midway, above, labelnode] {pendant ce temps};
\draw[arrow] (bian2) -- (final);

\node[labelnode, below=0.8cm of bian1] (v1) {Action 1 : manger};
\node[labelnode, below=0.8cm of bian2] (v2) {Action 2 : regarder TV};
\draw[popline, dashed] (bian1.south) -- (v1.north);
\draw[popline, dashed] (bian2.south) -- (v2.north);
\end{tikzpicture}
\legende{Architecture de la phrase avec 一边……一边…… : un sujet, deux actions simultanées, ponctuation pleine chasse.}
\end{popfigure}

\courssub{Fiche méthode 1 : Construire une phrase avec 一边……一边……}

\begin{methode}[Construire 一边……一边…… en 5 étapes]
La structure \zh{一边……一边……} (yíbiān……yíbiān……) exprime deux actions réalisées \textbf{en même temps} par le même sujet. Voici la démarche à suivre.

\begin{enumerate}
\item \textbf{Identifier le sujet unique.} Les deux actions doivent être faites par la même personne (ou le même groupe). Si les sujets sont différents, la structure est impossible : il faut deux phrases séparées ou une structure \zh{……的时候……}.
\item \textbf{Placer le sujet une seule fois}, avant le premier \zh{一边} :\\
\cle{Sujet + 一边 + Verbe$_1$ (+ Complément$_1$) ，一边 + Verbe$_2$ (+ Complément$_2$) 。}
\item \textbf{Choisir deux actions compatibles} (qu'on peut physiquement faire ensemble : écouter + marcher, chanter + danser ; pas dormir + manger, pas courir + dormir).
\item \textbf{Mettre la virgule chinoise} \zh{，} entre les deux membres, et le point chinois \zh{。} à la fin. Pas d'espace avant/après les signes.
\item \textbf{Vérifier l'ordre des mots} dans chaque membre : le complément suit son verbe (\zh{一边听音乐，一边做作业}, non *\zh{一边听，一边做音乐作业}).
\end{enumerate}

\begin{center}
\begin{tabularx}{\linewidth}{|X|X|X|}
\hline
\rowcolor{popblueL} Structure & Exemple (caractères + pinyin) & Traduction \\
\hline
Sujet + 一边 + V$_1$ ，一边 + V$_2$ 。 & \zh{他一边吃饭，一边看电视。} Tā yíbiān chīfàn, yíbiān kàn diànshì. & Il mange en regardant la télévision. \\
\hline
Sujet + 一边 + V$_1$ ，一边 + V$_2$ 。 & \zh{我们一边走路，一边说话。} Wǒmen yíbiān zǒulù, yíbiān shuōhuà. & Nous marchons en parlant. \\
\hline
Lieu + Sujet + 一边 + V$_1$ ，一边 + V$_2$ 。 & \zh{在雅温得，年轻人一边走路，一边听音乐。} Zài Yǎwēndé, niánqīngrén yíbiān zǒulù, yíbiān tīng yīnyuè. & À Yaoundé, les jeunes marchent en écoutant de la musique. \\
\hline
\end{tabularx}
\end{center}

\textbf{Réflexe examen :} devant une traduction du type « en + gérondif » français (\emph{en chantant, en écoutant}), « tout en + infinitif » ou « pendant que + même sujet », pensez immédiatement à \zh{一边……一边……}.
\end{methode}

\begin{exoresolu}[Construction de phrases — Type Bac]
\textbf{Énoncé.} Traduisez en chinois, avec la structure \zh{一边……一边……}, les phrases suivantes :\\
1. Marie écoute de la musique en faisant ses devoirs.\\
2. Les élèves chantent en dansant.\\
3. Mon père conduit en téléphonant. (Attention : est-ce une bonne idée ?)
\tcblower
\textbf{Solution détaillée.}

\textbf{Phrase 1.} Sujet unique : Marie (\zh{玛丽} Mǎlì). Les deux actions : écouter de la musique \zh{听音乐} et faire ses devoirs \zh{做作业}. On place le sujet une fois, puis les deux membres séparés par \zh{，} :\\
\zh{玛丽一边听音乐，一边做作业。}\\
\emph{Mǎlì yíbiān tīng yīnyuè, yíbiān zuò zuòyè.}\\
« Marie écoute de la musique en faisant ses devoirs. »\\
Justification : le sujet \zh{玛丽} n'est pas répété ; chaque verbe garde son complément (\zh{音乐}, \zh{作业}) juste après lui.

\textbf{Phrase 2.} Sujet : les élèves \zh{学生们} (xuéshengmen). Actions : chanter \zh{唱歌} et danser \zh{跳舞}.\\
\zh{学生们一边唱歌，一边跳舞。}\\
\emph{Xuéshengmen yíbiān chànggē, yíbiān tiàowǔ.}\\
« Les élèves chantent en dansant. »

\textbf{Phrase 3.} Sujet : mon père \zh{我爸爸} (wǒ bàba). Actions : conduire \zh{开车} et téléphoner \zh{打电话}. La structure est grammaticalement correcte :\\
\zh{我爸爸一边开车，一边打电话。}\\
\emph{Wǒ bàba yíbiān kāichē, yíbiān dǎ diànhuà.}\\
« Mon père conduit en téléphonant. »\\
Remarque culturelle : la phrase est correcte, mais l'action est dangereuse et interdite ; en Chine comme au Cameroun, téléphoner au volant est sanctionné. On peut ajouter un commentaire : \zh{这很危险！} \emph{Zhè hěn wēixiǎn!} « C'est très dangereux ! »

\textbf{Conclusion.} Dans les trois cas : un seul sujet, deux verbes d'action simultanée, virgule chinoise entre les membres, point chinois final.
\end{exoresolu}

\courssub{Fiche méthode 2 : Emplois, limites et transformation de la structure}

\begin{methode}[Bien employer 一边……一边…… : les 4 réflexes]
\begin{enumerate}
\item \textbf{Un seul sujet.} \zh{我一边看书，弟弟一边玩} est \textbf{faux} : deux sujets différents. Il faut dire : \zh{我看书，弟弟玩。} ou \zh{我看书的时候，弟弟在玩。} (Quand je lis, mon petit frère joue).
\item \textbf{Deux actions actives et volontaires.} On évite les verbes d'état, psychiques ou involontaires : on ne dit pas *\zh{一边知道，一边明白} ni *\zh{一边喜欢，一边爱}. La structure demande des verbes d'action dynamique (\zh{听、说、唱、走、吃、做、看、跑、喝}).
\item \textbf{Variante soutenue :} \zh{一面……一面……} (yímiàn……yímiàn……) est un synonyme plus écrit, plus formel : \zh{他一面工作，一面学习。} \emph{Tā yímiàn gōngzuò, yímiàn xuéxí.} « Il travaille tout en étudiant. » (À l'oral, on préfère \zh{一边}).
\item \textbf{Négation et question.}
\begin{itemize}
\item Négation partielle : \zh{不} devant le verbe concerné : \zh{他一边吃饭，一边不说话。} (Il mange sans parler).
\item Négation totale : \zh{他不是一边吃饭，一边看电视。} (Ce n'est pas le cas qu'il mange en regardant la télé).
\item Question : \zh{你一边吃饭，一边看电视吗？} ou \zh{你是不是一边吃饭，一边看电视？}
\end{itemize}
\end{enumerate}
\textbf{Transformation type examen :} deux phrases simples \zh{他吃饭。他看电视。} peuvent être fusionnées en une phrase à \zh{一边……一边……} si le sujet est identique et les actions simultanées.
\end{methode}

\begin{exoresolu}[Transformation et correction — Type Bac]
\textbf{Énoncé.}\\
A. Fusionnez les deux phrases en une seule avec \zh{一边……一边……} :\\
1. 妈妈做饭。妈妈听音乐。\\
2. 我走路。我说话。\\
B. Corrigez la phrase fautive et expliquez l'erreur :\\
3. 我一边睡觉，一边跑步。
\tcblower
\textbf{Solution détaillée.}

\textbf{A.1.} Les deux phrases ont le même sujet \zh{妈妈} et décrivent des actions simultanées possibles. On supprime la répétition du sujet et on applique le schéma :\\
\zh{妈妈一边做饭，一边听音乐。}\\
\emph{Māma yíbiān zuòfàn, yíbiān tīng yīnyuè.}\\
« Maman cuisine en écoutant de la musique. »

\textbf{A.2.} Même raisonnement :\\
\zh{我一边走路，一边说话。}\\
\emph{Wǒ yíbiān zǒulù, yíbiān shuōhuà.}\\
« Je marche en parlant. »

\textbf{B.3.} La phrase \zh{我一边睡觉，一边跑步。} est \textbf{grammaticalement construite} mais \textbf{logiquement impossible} : on ne peut pas dormir et courir en même temps. La structure \zh{一边……一边……} exige deux actions \emph{compatibles}. Correction possible :\\
\zh{我一边跑步，一边听音乐。}\\
\emph{Wǒ yíbiān pǎobù, yíbiān tīng yīnyuè.}\\
« Je cours en écoutant de la musique. »

\textbf{Conclusion.} La transformation exige trois vérifications : même sujet, actions simultanées, actions compatibles. Au Bac, une phrase grammaticalement correcte mais absurde est pénalisée.
\end{exoresolu}

\courssub{Fiche méthode 3 : Comparaison avec le français et pièges des francophones}

\begin{methode}[Du français au chinois sans calque]
\begin{enumerate}
\item \textbf{Repérer l'équivalent.} Le français « en + gérondif » (\emph{en mangeant}), « tout en + infinitif » et « pendant qu'il/elle + verbe » (même sujet) se traduisent par \zh{一边……一边……}.
\item \textbf{Ne pas traduire mot à mot.} Le français « en » n'a pas d'équivalent isolé en chinois : il est rendu par la \emph{structure entière}. N'écrivez jamais un mot-outil supplémentaire pour traduire « en » (pas de \zh{在}, \zh{当}, \zh{时} ici).
\item \textbf{Ne pas répéter le sujet.} Le français répète souvent le sujet ou le pronom (« il mange et il regarde ») ; le chinois ne le dit qu'une fois : \zh{他一边吃饭，一边看电视。}
\item \textbf{Attention à 和 (hé).} \zh{和} relie des \emph{noms}, pas des actions simultanées : *\zh{他吃饭和看电视} est incorrect pour exprimer la simultanéité. Pour relier deux actions simultanées, seule la structure \zh{一边……一边……} convient.
\item \textbf{Pas de 了 (le) dans les deux membres} pour une habitude présente ou une scène en cours : *\zh{他一边吃了饭，一边看了电视} est fautif ; \zh{了} marque l'accomplissement/le changement d'état, pas la simultanéité. On dit : \zh{他一边吃饭，一边看电视。}
\end{enumerate}
\end{methode}

\begin{exoresolu}[Thème : du français au chinois — Type Bac]
\textbf{Énoncé.} Traduisez en chinois :\\
1. À Yaoundé, les jeunes marchent en écoutant de la musique.\\
2. Elle apprend le chinois tout en regardant des films chinois.\\
3. Il mange et regarde la télévision en même temps.
\tcblower
\textbf{Solution détaillée.}

\textbf{Phrase 1.} « En écoutant » = gérondif français $\rightarrow$ \zh{一边……一边……}. Sujet : les jeunes \zh{年轻人} ; lieu \zh{在雅温得} placé après le sujet (le complément de lieu précède le verbe en chinois).\\
\zh{在雅温得，年轻人一边走路，一边听音乐。}\\
\emph{Zài Yǎwēndé, niánqīngrén yíbiān zǒulù, yíbiān tīng yīnyuè.}\\
« À Yaoundé, les jeunes marchent en écoutant de la musique. »

\textbf{Phrase 2.} « Tout en regardant » $\rightarrow$ même structure. Apprendre le chinois : \zh{学汉语} ; regarder des films chinois : \zh{看中国电影}.\\
\zh{她一边看中国电影，一边学汉语。}\\
\emph{Tā yíbiān kàn Zhōngguó diànyǐng, yíbiān xué Hànyǔ.}\\
« Elle apprend le chinois tout en regardant des films chinois. »

\textbf{Phrase 3.} Piège : « et » + « en même temps » ne se traduit pas par \zh{和}. Il faut la structure de simultanéité :\\
\zh{他一边吃饭，一边看电视。}\\
\emph{Tā yíbiān chīfàn, yíbiān kàn diànshì.}\\
« Il mange et regarde la télévision en même temps. »

\textbf{Conclusion.} Trois formulations françaises différentes (gérondif, « tout en », « et... en même temps ») aboutissent à une seule structure chinoise : \zh{一边……一边……}. C'est le réflexe à acquérir.
\end{exoresolu}

\courssub{Fiche méthode 4 : La ponctuation en chinois}

\begin{methode}[Utiliser les signes de ponctuation chinois]
Les signes chinois sont \textbf{pleine chasse} (ils occupent la largeur d'un caractère) et se collent au caractère précédent, sans espace. Ne jamais utiliser la ponctuation française (.,;:?!) dans un texte chinois.

\begin{center}
\begin{tabularx}{\linewidth}{|X|X|X|}
\hline
\rowcolor{popblueL} Signe & Nom et fonction & Exemple (caractères + pinyin) \\
\hline
\zh{。} & Point (fin de phrase déclarative) & \zh{我喜欢汉语。} Wǒ xǐhuan Hànyǔ. \\
\hline
\zh{，} & Virgule (pause entre propositions) & \zh{我一边吃饭，一边看电视。} Wǒ yíbiān chīfàn, yíbiān kàn diànshì. \\
\hline
\zh{、} & Virgule énumérative (entre mots d'une liste) & \zh{我买了苹果、香蕉和橙子。} Wǒ mǎi le píngguǒ, xiāngjiāo hé chéngzi. \\
\hline
\zh{？} & Point d'interrogation & \zh{你好吗？} Nǐ hǎo ma? \\
\hline
\zh{！} & Point d'exclamation & \zh{太好了！} Tài hǎo le! \\
\hline
\zh{：} & Deux-points (annonce, citation) & \zh{他说：« 你好。 »} Tā shuō: « Nǐ hǎo. » \\
\hline
\zh{；} & Point-virgule (deux propositions liées) & \zh{他学汉语；我学法语。} Tā xué Hànyǔ; wǒ xué Fǎyǔ. \\
\hline
\zh{《 》} & Guillemets de titres (livres, films, journaux) & \zh{我看《人民日报》。} Wǒ kàn « Rénmín Rìbào ». \\
\hline
\zh{« »} & Guillemets de citation (paroles, dialogues) & \zh{老师说：« 请坐。 »} Lǎoshī shuō: « Qǐng zuò. » \\
\hline
\zh{……} & Points de suspension (hésitation, liste ouverte) & \zh{我想……我不知道。} Wǒ xiǎng…… wǒ bù zhīdào. \\
\hline
\end{tabularx}
\end{center}

\textbf{Réflexes indispensables :}
\begin{enumerate}
\item Jamais de point français « . » ni de virgule française « , » dans une phrase chinoise.
\item \zh{、} ne s'emploie \textbf{qu'entre des mots énumérés} (noms, verbes, adjectifs), jamais entre deux propositions complètes.
\item Avec \zh{吗}, le point d'interrogation \zh{？} reste obligatoire : \zh{你是学生吗？}
\item Le dernier élément d'une énumération est souvent introduit par \zh{和} (ou \zh{跟}), sans \zh{、} devant : \zh{苹果、香蕉和橙子}.
\item Les guillemets \zh{« »} encadrent la citation ; le point final se place \textbf{à l'intérieur} des guillemets si la citation est une phrase complète.
\end{enumerate}
\end{methode}

\begin{exoresolu}[Ponctuation — Type Bac]
\textbf{Énoncé.} Ponctuez correctement les textes suivants (recopiez-les avec les signes chinois appropriés) :\\
1. 你叫什么名字 我叫李明\\
2. 我喜欢苹果 香蕉 芒果\\
3. 他一边唱歌 一边跳舞 真棒\\
4. 老师说 同学们好
\tcblower
\textbf{Solution détaillée.}

\textbf{Phrase 1.} La première proposition est une question $\rightarrow$ \zh{？} ; la réponse est déclarative $\rightarrow$ \zh{。}\\
\zh{你叫什么名字？我叫李明。}\\
\emph{Nǐ jiào shénme míngzi? Wǒ jiào Lǐ Míng.}\\
« Comment t'appelles-tu ? Je m'appelle Li Ming. »

\textbf{Phrase 2.} Énumération de noms $\rightarrow$ virgule énumérative \zh{、} (et non \zh{，}) ; point final \zh{。}\\
\zh{我喜欢苹果、香蕉、芒果。}\\
\emph{Wǒ xǐhuan píngguǒ, xiāngjiāo, mángguǒ.}\\
« J'aime les pommes, les bananes et les mangues. »

\textbf{Phrase 3.} Structure \zh{一边……一边……} $\rightarrow$ virgule \zh{，} entre les deux membres ; « 真棒 » est une exclamation $\rightarrow$ \zh{！}\\
\zh{他一边唱歌，一边跳舞，真棒！}\\
\emph{Tā yíbiān chànggē, yíbiān tiàowǔ, zhēn bàng!}\\
« Il chante en dansant, c'est formidable ! »

\textbf{Phrase 4.} Parole rapportée $\rightarrow$ deux-points \zh{：} puis guillemets de citation \zh{« »} ; point à l'intérieur des guillemets :\\
\zh{老师说：« 同学们好。 »}\\
\emph{Lǎoshī shuō: « Tóngxuémen hǎo. »}\\
« Le professeur dit : "Bonjour, les élèves." »

\textbf{Conclusion.} Chaque signe a une fonction précise : \zh{、} pour les listes de mots, \zh{，} pour les pauses entre propositions, \zh{：« »} pour introduire une parole. Confondre \zh{、} et \zh{，} est l'erreur la plus fréquente.
\end{exoresolu}

\courssub{Fiche méthode 5 : Entraînement intégré — Rédiger un court texte}

\begin{methode}[Rédiger un court texte avec 一边……一边…… et une ponctuation correcte]
\begin{enumerate}
\item \textbf{Choisir une scène de la vie quotidienne} (à la maison, au marché de Douala, en classe, devant la télévision, dans un taxi).
\item \textbf{Préparer 2 ou 3 paires d'actions simultanées} : \zh{听音乐 + 做作业}, \zh{做饭 + 唱歌}, \zh{看比赛 + 吃花生}, \zh{开车 + 听广播}.
\item \textbf{Rédiger des phrases courtes} : sujet unique, structure complète, un complément après chaque verbe.
\item \textbf{Ponctuer au fur et à mesure} : \zh{，} entre les membres de \zh{一边……一边……}, \zh{。} en fin de phrase déclarative, \zh{！} pour une exclamation finale, \zh{：« »} pour le discours direct.
\item \textbf{Relire} : vérifier qu'il n'y a ni point français, ni sujet répété, ni \zh{和} entre deux verbes d'action simultanée, ni \zh{、} entre propositions.
\end{enumerate}
\textbf{Modèle de rédaction :}\\
\zh{晚上，我们一家人在客厅。爸爸一边看新闻，一边喝茶。妈妈一边做饭，一边唱歌。我一边做作业，一边听音乐。我们家真热闹！}\\
\emph{Wǎnshang, wǒmen yì jiā rén zài kètīng. Bàba yíbiān kàn xīnwén, yíbiān hē chá. Māma yíbiān zuòfàn, yíbiān chànggē. Wǒ yíbiān zuò zuòyè, yíbiān tīng yīnyuè. Wǒmen jiā zhēn rènao!}\\
« Le soir, toute notre famille est au salon. Papa regarde les informations en buvant du thé. Maman cuisine en chantant. Je fais mes devoirs en écoutant de la musique. Comme c'est animé chez nous ! »
\end{methode}

\begin{exoresolu}[Composition guidée — Type Bac]
\textbf{Énoncé.} Rédigez un court texte de 4 à 5 phrases en chinois sur le thème « Un soir de match de football à la maison », en utilisant au moins deux fois la structure \zh{一边……一边……} et une ponctuation correcte.
\tcblower
\textbf{Solution détaillée.}

\textbf{Étape 1 — Plan.} Scène : la famille regarde un match des Lions Indomptables. Personnages : papa, moi, mon frère. Actions simultanées : regarder le match + manger ; crier + sauter (joie) ; commenter + boire.

\textbf{Étape 2 — Rédaction.}\\
\zh{今天晚上，我们一家人看足球比赛。爸爸一边看比赛，一边吃花生。我一边看，一边喊：« 加油！ » 喀麦隆队进球了，我们高兴极了！}\\
\emph{Jīntiān wǎnshang, wǒmen yì jiā rén kàn zúqiú bǐsài. Bàba yíbiān kàn bǐsài, yíbiān chī huāshēng. Wǒ yíbiān kàn, yíbiān hǎn: « Jiāyóu! » Kāmàilóng duì jìnqiú le, wǒmen gāoxìng jí le!}\\
« Ce soir, toute notre famille regarde un match de football. Papa regarde le match en mangeant des cacahuètes. Je regarde en criant : "Allez !" L'équipe du Cameroun a marqué, nous sommes très contents ! »

\textbf{Étape 3 — Vérification grammaticale.}
\begin{itemize}
\item Deux emplois de \zh{一边……一边……} : sujet unique à chaque fois (\zh{爸爸}, \zh{我}), actions compatibles (regarder + manger ; regarder + crier).
\item Ponctuation : \zh{，} pour les pauses, \zh{。} pour les phrases déclaratives, \zh{：« »} pour la parole criée, \zh{！} pour l'exclamation.
\item \zh{进球了} : \zh{了} marque ici l'accomplissement du but (changement de situation), ce qui est correct.
\item Aucun calque du français : pas de \zh{和} entre les verbes, pas de sujet répété.
\end{itemize}

\textbf{Conclusion.} Un texte court, vivant et correct : deux structures de simultanéité bien formées et une ponctuation irréprochable suffisent pour obtenir la totalité des points de production écrite.
\end{exoresolu}

\begin{exercice}[Entraînement personnel — Devoirs / Classe]
\textbf{A. Complétez avec la structure 一边……一边…… (choisissez les verbes adaptés) :}
1. 妹妹 \_\_\_\_\_\_\_\_ ，\_\_\_\_\_\_\_\_ 。(regarder la télé / manger des fruits)
2. 我们 \_\_\_\_\_\_\_\_ ，\_\_\_\_\_\_\_\_ 。(marcher / parler)
3. 哥哥 \_\_\_\_\_\_\_\_ ，\_\_\_\_\_\_\_\_ 。(écouter de la musique / faire les devoirs)

\textbf{B. Corrigez les erreurs (grammaire, logique ou ponctuation) :}
1. 我一边吃饭和看电视。
2. 他一边开车，他一边打电话。
3. 妈妈一边做饭、一边听音乐。
4. 学生们一边睡觉，一边上课。
5. 她说：我累了。

\textbf{C. Traduisez en chinois (soignez la ponctuation) :}
1. Ma sœur lit un livre en buvant du café.
2. Au marché de Douala, les gens marchent et parlent.
3. Ne conduis pas en téléphonant ! (Impératif négatif : 别...)

\textbf{D. Production écrite :} Décrivez en 4 phrases ce que fait votre famille un dimanche après-midi (utilisez \zh{一边……一边……} deux fois minimum).
\end{exercice}

\begin{corrige}[Exercice Leçon 40]
\textbf{A.}
1. \zh{妹妹一边看电视，一边吃水果。} (Mèimei yíbiān kàn diànshì, yíbiān chī shuǐguǒ.)
2. \zh{我们一边走路，一边说话。} (Wǒmen yíbiān zǒulù, yíbiān shuōhuà.)
3. \zh{哥哥一边听音乐，一边做作业。} (Gēge yíbiān tīng yīnyuè, yíbiān zuò zuòyè.)

\textbf{B.}
1. \zh{我一边吃饭，一边看电视。} (Erreur : \zh{和} ne relie pas deux actions simultanées).
2. \zh{他一边开车，一边打电话。} (Erreur : répétition du sujet \zh{他}).
3. \zh{妈妈一边做饭，一边听音乐。} (Erreur : \zh{、} s'utilise pour énumérer des mots, pas des propositions).
4. \zh{学生们一边上课，一边听讲。} ou \zh{学生们上课。} (Erreur logique : on ne dort pas en cours. Correction : actions compatibles).
5. \zh{她说：« 我累了。 »} (Erreur : manque deux-points, guillemets chinois, point final).

\textbf{C.}
1. \zh{我妹妹一边看书，一边喝咖啡。} (Wǒ mèimei yíbiān kàn shū, yíbiān hē kāfēi.)
2. \zh{在杜阿拉市场，人们一边走路，一边说话。} (Zài Dùālā shìchǎng, rénmen yíbiān zǒulù, yíbiān shuōhuà.)
3. \zh{别一边开车，一边打电话！} (Bié yíbiān kāichē, yíbiān dǎ diànhuà!)

\textbf{D. (Exemple de production)} \\
\zh{星期天下午，我们全家在家里。爸爸一边看报纸，一边喝茶。妈妈一边打扫卫生，一边听广播。我和弟弟一边踢足球，一边聊天。大家都很开心！} \\
\emph{Xīngqítiān xiàwǔ, wǒmen quánjiā zài jiālǐ. Bàba yíbiān kàn bàozhǐ, yíbiān hē chá. Māma yíbiān dǎsǎo wèishēng, yíbiān tīng guǎngbō. Wǒ hé dìdi yíbiān tī zúqiú, yíbiān liáotiān. Dàjiā dōu hěn kāixīn!}
\end{corrige}

\begin{experience}[Activité orale : « Journaliste d'un soir »]
\textbf{Consigne :} Par binômes. L'un est journaliste (à la télévision chinoise CCTV ou camerounaise CRTV), l'autre est un témoin dans la rue (à Yaoundé ou Pékin).
\textbf{Scénario :} « Comment les gens passent-ils leur temps libre ? »
\textbf{Modèle :}
\begin{itemize}
\item Journaliste : \zh{您好，请问您平时喜欢怎么放松？} (Nín hǎo, qǐngwèn nín píngshí xǐhuan zěnme fàngsōng?)
\item Témoin : \zh{我喜欢一边听音乐，一边跑步。} (Wǒ xǐhuan yíbiān tīng yīnyuè, yíbiān pǎobù.)
\item Journaliste : \zh{还有呢？} (Hái yǒu ne?)
\item Témoin : \zh{有时候一边喝茶，一边看书。} (Yǒushíhou yíbiān hē chá, yíbiān kàn shū.)
\end{itemize}
\textbf{Contrainte :} Utiliser au moins 3 fois \zh{一边……一边……}. Le journaliste note les réponses en caractères + pinyin. Échange des rôles après 3 minutes.
\end{experience}

\begin{savaistu}[Le saviez-vous ? Ponctuation et histoire]
Le chinois classique s'écrivait sans ponctuation (文言文 sans espaces ni signes). Les lettrés lisaient à voix haute pour « segmenter » le texte (句读, jùdòu). La ponctuation moderne (标点符号, biāodiǎn fúhào) a été officialisée en \textbf{1919} (école de Pékin) puis \textbf{1930} (ministère de l'Éducation), inspirée des signes occidentaux mais adaptée à la grille carrée (pleine chasse).
\textbf{Anecdote :} Le signe \zh{、} (dùn hào) est une invention purement chinoise pour l'énumération de mots courts, là où l'occidental met une virgule. En Chine, un texte officiel sans ponctuation correcte est considéré comme « non abouti » (不规范).
\end{savaistu}

\begin{attention}[Les 5 erreurs qui coûtent des points au Bac]
\begin{enumerate}
\item \textbf{Répéter le sujet} dans le deuxième membre : *\zh{他一边吃饭，他一边看电视}. Le sujet ne se dit qu'une fois, avant le premier \zh{一边}.
\item \textbf{Traduire « et » par 和 entre deux actions simultanées} : *\zh{他吃饭和看电视}. \zh{和} relie des noms ; pour deux actions en même temps, seule la structure \zh{一边……一边……} convient.
\item \textbf{Confondre 、 et ，} : *\zh{他一边唱歌、一边跳舞}. \zh{、} sépare des \emph{mots} dans une liste ; entre deux propositions, il faut \zh{，}. Utiliser le point français « . » ou la virgule française « , » dans un texte chinois est également pénalisé.
\item \textbf{Associer des actions incompatibles} : *\zh{一边睡觉，一边跑步}. Le correcteur sanctionne une phrase grammaticalement juste mais logiquement impossible (sémantique).
\item \textbf{Mal placer les compléments / Tons et caractères} : *\zh{一边听，一边做音乐作业} (ordre incorrect). Attention aux tons : \zh{一} $\rightarrow$ \textbf{yí} (2e ton) devant \zh{边} (1er ton). Écriture : \zh{边} (biān, clé \zh{辶} chuò, 5 traits) $\neq$ \zh{过} (guò) $\neq$ \zh{这} (zhè). Un trait manquant sur \zh{辶} ou ordre des traits faux (horizontal avant vertical) coûte le point d'écriture.
\end{enumerate}
\end{attention}

\newpage

\coursec{Exercices}

\courssub{Je vérifie mes connaissances}

\begin{exercice}[QCM : la bonne phrase]
Pour chaque question, une seule phrase est correcte. Entoure la lettre correspondante.

\textbf{1.} « Il écoute de la musique en faisant ses devoirs. »
\begin{itemize}
\item[A.] 他一边做作业，一边听音乐。(Tā yìbiān zuò zuòyè, yìbiān tīng yīnyuè.)
\item[B.] 一边他做作业，一边听音乐。
\item[C.] 他一边做作业，听音乐。
\item[D.] 他听音乐一边做作业一边。
\end{itemize}

\textbf{2.} « Maman cuisine tout en parlant au téléphone. »
\begin{itemize}
\item[A.] 妈妈做饭一边，一边打电话。
\item[B.] 妈妈一边打电话，一边做饭。(Māma yìbiān dǎ diànhuà, yìbiān zuò fàn.)
\item[C.] 妈妈一边打电话，做饭。
\item[D.] 一边妈妈打电话，一边做饭。
\end{itemize}

\textbf{3.} Quelle phrase est impossible avec \zh{一边……一边……} ?
\begin{itemize}
\item[A.] 他一边吃饭，一边看报纸。
\item[B.] 我们一边走路，一边聊天。
\item[C.] 他一边睡觉，一边跑步。(Tā yìbiān shuìjiào, yìbiān pǎobù.)
\item[D.] 她一边唱歌，一边跳舞。
\end{itemize}

\textbf{4.} Quel signe utilise-t-on pour une énumération en chinois ?
\begin{itemize}
\item[A.] ， \hspace{1cm} B. 、 \hspace{1cm} C. 。 \hspace{1cm} D. ？
\end{itemize}
\end{exercice}

\begin{exercice}[Vrai ou faux ? Justifie ta réponse]
Dis si les affirmations suivantes sont vraies ou fausses, puis justifie en une phrase en français.
\begin{enumerate}
\item Dans la structure \zh{一边……一边……}, le sujet se place avant le premier \zh{一边}.
\item On peut écrire \zh{他一边吃饭，看电视。} sans répéter \zh{一边}.
\item Les deux actions reliées par \zh{一边……一边……} doivent pouvoir se faire en même temps.
\item En chinois, le point final s'écrit \zh{。} et non pas « . ».
\item Les guillemets chinois s'écrivent \zh{“”} et le titre d'une émission ou d'un livre se met entre \zh{《》}.
\end{enumerate}
\end{exercice}

\begin{exercice}[Vocabulaire : les médias]
Complète le tableau en écrivant la traduction française de chaque mot.

\begin{center}
\begin{tabularx}{\linewidth}{|X|X|X|X|}
\hline
\rowcolor{popblueL} Caractères & Pinyin & Nature & Traduction \\
\hline
报纸 & bàozhǐ & nom & \ldots\ldots\ldots\ldots \\
\hline
杂志 & zázhì & nom & \ldots\ldots\ldots\ldots \\
\hline
电视 & diànshì & nom & \ldots\ldots\ldots\ldots \\
\hline
广播 & guǎngbō & nom & \ldots\ldots\ldots\ldots \\
\hline
节目 & jiémù & nom & \ldots\ldots\ldots\ldots \\
\hline
手机 & shǒujī & nom & \ldots\ldots\ldots\ldots \\
\hline
\end{tabularx}
\end{center}
\end{exercice}

\begin{exercice}[Associe le pinyin aux caractères]
Relie chaque mot en caractères à son pinyin.

\begin{center}
\begin{tabularx}{\linewidth}{|X|X|}
\hline
\rowcolor{popblueL} Caractères & Pinyin à associer \\
\hline
1. 听 & a. kàn \\
\hline
2. 看 & b. xiě \\
\hline
3. 写 & c. tīng \\
\hline
4. 说 & d. chàng \\
\hline
5. 唱 & e. shuō \\
\hline
6. 一边 & f. yìbiān \\
\hline
\end{tabularx}
\end{center}
\end{exercice}

\courssub{Je m'entraîne}

\begin{exercice}[Traduis en français]
Traduis les phrases suivantes en français correct.
\begin{enumerate}
\item 爸爸一边看报纸，一边喝茶。(Bàba yìbiān kàn bàozhǐ, yìbiān hē chá.)
\item 我们一边走路，一边说汉语。(Wǒmen yìbiān zǒulù, yìbiān shuō Hànyǔ.)
\item 她一边做饭，一边听广播。(Tā yìbiān zuò fàn, yìbiān tīng guǎngbō.)
\item 学生们一边听老师讲课，一边记笔记。(Xuéshengmen yìbiān tīng lǎoshī jiǎngkè, yìbiān jì bǐjì.)
\end{enumerate}
\end{exercice}

\begin{exercice}[Transformation : fusionne les phrases]
Fusionne chaque paire de phrases en une seule phrase avec \zh{一边……一边……}. Écris la phrase en caractères, avec le pinyin et la traduction française.

\begin{enumerate}
\item 他开车。他打电话。(Tā kāichē. Tā dǎ diànhuà.)
\item 学生们听老师讲课。学生们记笔记。(Xuéshengmen tīng lǎoshī jiǎngkè. Xuéshengmen jì bǐjì.)
\item 我看电视。我吃苹果。(Wǒ kàn diànshì. Wǒ chī píngguǒ.)
\item 妈妈做饭。妈妈唱歌。(Māma zuò fàn. Māma chànggē.)
\end{enumerate}
\end{exercice}

\begin{exercice}[Texte à trous]
Complète le texte suivant avec les mots de la liste : \zh{一边} (deux fois), \zh{、}, \zh{。}, \zh{《》}.

\begin{textebox}[Le soir chez moi]
晚上，我爸爸\ldots\ldots 看报纸，\ldots\ldots 听广播\ldots\ldots 妈妈做饭，我看电视节目\ldots\ldots 新闻联播\ldots\ldots 。桌子上有很多东西：报纸\ldots\ldots 杂志和书。

(Wǎnshang, wǒ bàba \ldots\ldots kàn bàozhǐ, \ldots\ldots tīng guǎngbō \ldots\ldots Māma zuò fàn, wǒ kàn diànshì jiémù \ldots\ldots Xīnwén Liánbō \ldots\ldots . Zhuōzi shang yǒu hěn duō dōngxi : bàozhǐ \ldots\ldots zázhì hé shū.)

« Le soir, mon père lit le journal tout en écoutant la radio. Maman cuisine, je regarde l'émission \emph{Journal télévisé}. Sur la table, il y a beaucoup de choses : des journaux, des magazines et des livres. »
\end{textebox}
\end{exercice}

\begin{exercice}[Remets les mots dans l'ordre]
Remets les mots suivants dans le bon ordre pour former une phrase correcte. Écris la phrase en caractères avec le pinyin et la traduction.

\begin{enumerate}
\item 一边 / 她 / 唱歌 / 跳舞 / 一边
\item 一边 / 我 / 一边 / 吃饭 / 看手机
\item 打电话 / 他 / 开车 / 一边 / 一边
\item 听 / 我们 / 音乐 / 做作业 / 一边 / 一边
\end{enumerate}
\end{exercice}

\courssub{Je me prépare au Bac}

\begin{exercice}[Type Bac — Compréhension de texte et questions de langue]
Lis le texte suivant, puis réponds aux questions.

\begin{textebox}[晚上的一家人]
晚上七点，我们一家人都在家。爸爸一边看报纸，一边喝茶。妈妈一边做饭，一边听广播。我姐姐一边看电视，一边用手机给朋友发消息。我呢？我一边做作业，一边听音乐。我们都很喜欢看电视节目《新闻联播》。桌子上有很多东西：报纸、杂志、书和手机。

(Wǎnshang qī diǎn, wǒmen yì jiā rén dōu zài jiā. Bàba yìbiān kàn bàozhǐ, yìbiān hē chá. Māma yìbiān zuò fàn, yìbiān tīng guǎngbō. Wǒ jiějie yìbiān kàn diànshì, yìbiān yòng shǒujī gěi péngyou fā xiāoxi. Wǒ ne ? Wǒ yìbiān zuò zuòyè, yìbiān tīng yīnyuè. Wǒmen dōu hěn xǐhuan kàn diànshì jiémù 《Xīnwén Liánbō》. Zhuōzi shang yǒu hěn duō dōngxi : bàozhǐ, zázhì, shū hé shǒujī.)

« À sept heures du soir, toute notre famille est à la maison. Papa lit le journal tout en buvant du thé. Maman cuisine tout en écoutant la radio. Ma grande sœur regarde la télévision tout en envoyant des messages à ses amis avec son téléphone portable. Et moi ? Je fais mes devoirs tout en écoutant de la musique. Nous aimons tous beaucoup regarder l'émission \emph{Journal télévisé}. Sur la table, il y a beaucoup de choses : des journaux, des magazines, des livres et des téléphones portables. »
\end{textebox}

\textbf{A. Questions de compréhension} (réponds en français)
\begin{enumerate}
\item Que font le père et la mère le soir ?
\item Que fait la grande sœur avec son téléphone portable ?
\item Quelle émission de télévision la famille aime-t-elle regarder ?
\item Cite les quatre choses qui se trouvent sur la table.
\end{enumerate}

\textbf{B. Questions de langue}
\begin{enumerate}
\item Releve dans le texte deux phrases construites avec \zh{一边……一边……} et recopie-les avec leur pinyin.
\item Pourquoi le signe \zh{、} est-il utilisé dans la dernière phrase ? Explique la règle.
\item Pourquoi le titre \zh{《新闻联播》} est-il écrit entre \zh{《》} ?
\end{enumerate}
\end{exercice}

\begin{exercice}[Type Bac — Thème et version]
\textbf{A. Thème.} Traduis en chinois (caractères et pinyin) :

« Le soir, mon père lit le journal tout en écoutant la radio, et ma sœur envoie des messages tout en regardant la télévision. »

\textbf{B. Version.} Traduis en français :

\zh{星期天早上，我一边吃早饭，一边看报纸。我弟弟一边听音乐，一边做作业。}

(Xīngqītiān zǎoshang, wǒ yìbiān chī zǎofàn, yìbiān kàn bàozhǐ. Wǒ dìdi yìbiān tīng yīnyuè, yìbiān zuò zuòyè.)
\end{exercice}

\begin{exercice}[Type Bac — Correction, ponctuation et expression écrite]
\textbf{A. Correction.} Chaque phrase contient une erreur. Corrige-la et justifie ta correction en français.
\begin{enumerate}
\item *他一边吃饭，看电视。
\item *一边他唱歌，一边跳舞。
\item *他一边睡觉，一边跑步。
\end{enumerate}

\textbf{B. Ponctuation.} Réécris le texte suivant en ajoutant les signes de ponctuation corrects (\zh{。 ， 、 ？ ： “” 《》}) :

\zh{妈妈问我 你做完作业了吗 我说 我一边做作业 一边听了 新闻联播 我还买了报纸 杂志和书}

(Māma wèn wǒ \ldots nǐ zuò wán zuòyè le ma \ldots wǒ shuō \ldots wǒ yìbiān zuò zuòyè \ldots yìbiān tīng le \ldots Xīnwén Liánbō \ldots wǒ hái mǎi le bàozhǐ \ldots zázhì hé shū \ldots)

\textbf{C. Expression écrite.} Rédige un court paragraphe en chinois (40 à 60 caractères) intitulé \zh{我家的晚上} (La soirée chez moi). Tu dois utiliser :
\begin{itemize}
\item deux fois la structure \zh{一边……一边……} ;
\item une énumération avec le signe \zh{、} ;
\item le titre d'une émission entre \zh{《》}.
\end{itemize}
\end{exercice}

\newpage

\coursec{Corrigés des exercices}

\begin{corrige}[Exercice 1 — QCM : la bonne phrase]
\textbf{1. Réponse A.} \zh{他一边做作业，一边听音乐。} (Tā yìbiān zuò zuòyè, yìbiān tīng yīnyuè.) « Il écoute de la musique en faisant ses devoirs. » Le sujet \zh{他} se place avant le premier \zh{一边}, et \zh{一边} est répété devant chaque action. B est faux car le sujet ne peut pas se placer après \zh{一边} ; C est incomplet (il manque le deuxième \zh{一边}) ; D a un ordre des mots impossible.

\textbf{2. Réponse B.} \zh{妈妈一边打电话，一边做饭。} (Māma yìbiān dǎ diànhuà, yìbiān zuò fàn.) « Maman cuisine tout en parlant au téléphone. » Même justification : Sujet + \zh{一边} + verbe 1, \zh{一边} + verbe 2.

\textbf{3. Réponse C.} \zh{他一边睡觉，一边跑步。} est impossible : on ne peut pas dormir et courir en même temps. La structure \zh{一边……一边……} exige deux actions \emph{compatibles}, réalisables simultanément par le même sujet.

\textbf{4. Réponse B.} Le signe \zh{、} (dùnhào) sert à séparer les éléments d'une énumération. \zh{，} sépare les propositions, \zh{。} marque la fin d'une phrase, \zh{？} marque la question.
\end{corrige}

\begin{corrige}[Exercice 2 — Vrai ou faux ?]
\begin{enumerate}
\item \textbf{Vrai.} Le sujet se place une seule fois, avant le premier \zh{一边} : \zh{他一边吃饭，一边看电视。} On ne dit jamais \zh{一边他吃饭……}.
\item \textbf{Faux.} \zh{一边} doit être répété devant chaque action : \zh{他一边吃饭，一边看电视。} Contrairement au français (\emph{« en »} n'apparaît qu'une fois), le chinois exige la répétition.
\item \textbf{Vrai.} Les deux actions doivent pouvoir se dérouler simultanément : \zh{一边睡觉，一边跑步} est absurde.
\item \textbf{Vrai.} Le point final chinois est un petit cercle pleine chasse \zh{。}, et non le point occidental « . ».
\item \textbf{Vrai.} Les guillemets chinois sont \zh{“”} pour les paroles et citations ; les titres d'œuvres, d'émissions, de journaux et de livres se mettent entre \zh{《》} (shūmínghào), par exemple \zh{《新闻联播》}.
\end{enumerate}
\end{corrige}

\begin{corrige}[Exercice 3 — Vocabulaire : les médias]
\begin{center}
\begin{tabularx}{\linewidth}{|X|X|X|X|}
\hline
\rowcolor{popblueL} Caractères & Pinyin & Nature & Traduction \\
\hline
报纸 & bàozhǐ & nom & le journal (presse écrite) \\
\hline
杂志 & zázhì & nom & le magazine, la revue \\
\hline
电视 & diànshì & nom & la télévision \\
\hline
广播 & guǎngbō & nom & la radio (radiodiffusion) \\
\hline
节目 & jiémù & nom & l'émission, le programme \\
\hline
手机 & shǒujī & nom & le téléphone portable \\
\hline
\end{tabularx}
\end{center}
\textbf{Points de vigilance :} ne pas confondre \zh{电视} (la télévision, l'appareil ou le média) et \zh{节目} (l'émission) ; \zh{手机} se compose de \zh{手} (main) et \zh{机} (machine) : littéralement « machine de main ».
\end{corrige}

\begin{corrige}[Exercice 4 — Associe le pinyin aux caractères]
\begin{itemize}
\item 1. 听 — c. tīng (écouter)
\item 2. 看 — a. kàn (regarder, lire)
\item 3. 写 — b. xiě (écrire)
\item 4. 说 — e. shuō (parler, dire)
\item 5. 唱 — d. chàng (chanter)
\item 6. 一边 — f. yìbiān (d'un côté ; en même temps)
\end{itemize}
\textbf{Attention :} \zh{听} (tīng, 1er ton) et \zh{唱} (chàng, 4e ton) ont tous deux la clé de la bouche \zh{口}, mais des prononciations très différentes. \zh{看} (kàn) et \zh{说} (shuō) sont des verbes d'action fréquents dans la structure \zh{一边……一边……}.
\end{corrige}

\begin{corrige}[Exercice 5 — Traduis en français]
\begin{enumerate}
\item \zh{爸爸一边看报纸，一边喝茶。} — « Papa lit le journal tout en buvant du thé. » (\zh{看报纸} = lire le journal ; \zh{喝茶} = boire du thé.)
\item \zh{我们一边走路，一边说汉语。} — « Nous parlons chinois tout en marchant. » (\zh{走路} = marcher ; \zh{说汉语} = parler chinois.)
\item \zh{她一边做饭，一边听广播。} — « Elle cuisine tout en écoutant la radio. » (\zh{做饭} = faire la cuisine ; \zh{听广播} = écouter la radio.)
\item \zh{学生们一边听老师讲课，一边记笔记。} — « Les élèves écoutent le cours du professeur tout en prenant des notes. » (\zh{讲课} = faire cours ; \zh{记笔记} = prendre des notes.)
\end{enumerate}
\textbf{Méthode :} en français, on rend \zh{一边……一边……} par « tout en + gérondif » ou « en + gérondif » ; l'action introduite par « en » peut correspondre au premier ou au deuxième membre, selon ce qui est le plus naturel.
\end{corrige}

\begin{corrige}[Exercice 6 — Transformation : fusionne les phrases]
\textbf{Méthode :} on garde le sujet une seule fois au début, puis on place \zh{一边} devant chaque verbe, en supprimant la répétition du sujet.

\begin{enumerate}
\item \zh{他一边开车，一边打电话。} (Tā yìbiān kāichē, yìbiān dǎ diànhuà.) « Il téléphone tout en conduisant. » (Phrase grammaticalement correcte, mais dangereuse dans la réalité : au Cameroun comme en Chine, téléphoner au volant est interdit !)
\item \zh{学生们一边听老师讲课，一边记笔记。} (Xuéshengmen yìbiān tīng lǎoshī jiǎngkè, yìbiān jì bǐjì.) « Les élèves écoutent le cours du professeur tout en prenant des notes. »
\item \zh{我一边看电视，一边吃苹果。} (Wǒ yìbiān kàn diànshì, yìbiān chī píngguǒ.) « Je mange une pomme tout en regardant la télévision. »
\item \zh{妈妈一边做饭，一边唱歌。} (Māma yìbiān zuò fàn, yìbiān chànggē.) « Maman chante tout en faisant la cuisine. »
\end{enumerate}
\textbf{Point de vigilance :} ne jamais répéter le sujet après la virgule : *\zh{他一边开车，他一边打电话} est faux.
\end{corrige}

\begin{corrige}[Exercice 7 — Texte à trous]
Texte complété :

\zh{晚上，我爸爸一边看报纸，一边听广播。妈妈做饭，我看电视节目《新闻联播》。桌子上有很多东西：报纸、杂志和书。}

(Wǎnshàng, wǒ bàba yìbiān kàn bàozhǐ, yìbiān tīng guǎngbō. Māma zuò fàn, wǒ kàn diànshì jiémù 《Xīnwén Liánbō》. Zhuōzi shàng yǒu hěn duō dōngxi : bàozhǐ, zázhì hé shū.)

« Le soir, mon père lit le journal tout en écoutant la radio. Maman cuisine, je regarde l'émission \emph{Journal télévisé}. Sur la table, il y a beaucoup de choses : des journaux, des magazines et des livres. »

\textbf{Justifications :}
\begin{itemize}
\item Les deux premiers blancs prennent \zh{一边} : deux actions simultanées du même sujet (\zh{爸爸}).
\item Le troisième blanc prend \zh{。} : fin de la première phrase.
\item \zh{新闻联播} est le titre d'une émission : il se met entre \zh{《》}.
\item Le dernier blanc prend \zh{、} : énumération de noms (\zh{报纸、杂志和书}) ; devant \zh{和}, on ne met pas de signe.
\end{itemize}
\end{corrige}

\begin{corrige}[Exercice 8 — Remets les mots dans l'ordre]
\begin{enumerate}
\item \zh{她一边唱歌，一边跳舞。} (Tā yìbiān chànggē, yìbiān tiàowǔ.) « Elle chante tout en dansant. »
\item \zh{我一边吃饭，一边看手机。} (Wǒ yìbiān chīfàn, yìbiān kàn shǒujī.) « Je regarde mon téléphone tout en mangeant. »
\item \zh{他一边开车，一边打电话。} (Tā yìbiān kāichē, yìbiān dǎ diànhuà.) « Il téléphone tout en conduisant. »
\item \zh{我们一边做作业，一边听音乐。} (Wǒmen yìbiān zuò zuòyè, yìbiān tīng yīnyuè.) « Nous écoutons de la musique tout en faisant nos devoirs. »
\end{enumerate}
\textbf{Méthode :} repérer d'abord le sujet (un pronom ou un nom), le placer en tête, puis alterner \zh{一边} + groupe verbal, \zh{一边} + groupe verbal, et terminer par \zh{。}.
\end{corrige}

\begin{corrige}[Exercice 9 — Type Bac : compréhension de texte et questions de langue]
\textbf{A. Questions de compréhension} (4 points, 1 point par réponse)
\begin{enumerate}
\item Le père lit le journal tout en buvant du thé ; la mère cuisine tout en écoutant la radio.
\item Elle envoie des messages à ses amis avec son téléphone portable (tout en regardant la télévision).
\item La famille aime regarder l'émission \zh{《新闻联播》} (le Journal télévisé).
\item Sur la table, il y a des journaux, des magazines et des livres.
\end{enumerate}

\textbf{B. Questions de langue} (6 points)
\begin{enumerate}
\item (2 points) Deux phrases possibles parmi : \zh{爸爸一边看报纸，一边喝茶。} (Bàba yìbiān kàn bàozhǐ, yìbiān hē chá.) ; \zh{妈妈一边做饭，一边听广播。} (Māma yìbiān zuò fàn, yìbiān tīng guǎngbō.) ; \zh{我一边做作业，一边听音乐。} (Wǒ yìbiān zuò zuòyè, yìbiān tīng yīnyuè.)
\item (2 points) Le signe \zh{、} est utilisé parce que la dernière phrase contient une énumération de noms coordonnés (\zh{报纸、杂志、书和手机}). Règle : en chinois, les éléments d'une énumération sont séparés par \zh{、} et non par la virgule \zh{，} ; devant \zh{和} (et), on ne met aucun signe.
\item (2 points) \zh{新闻联播} est le titre d'une émission de télévision ; en chinois, les titres d'œuvres, de journaux, de livres et d'émissions s'écrivent entre les marques \zh{《》} (shūmínghào), qui remplacent l'italique ou les guillemets français.
\end{enumerate}

\textbf{Points de vigilance :} recopier les phrases chinoises sans faute de caractères ; ne pas traduire \zh{《新闻联播》} par des guillemets français dans la réponse en chinois.
\end{corrige}

\begin{corrige}[Exercice 10 — Type Bac : thème et version]
\textbf{A. Thème} (5 points)

\zh{晚上，我爸爸一边看报纸，一边听广播，我姐姐一边看电视，一边发消息。}

(Wǎnshàng, wǒ bàba yìbiān kàn bàozhǐ, yìbiān tīng guǎngbō, wǒ jiějie yìbiān kàn diànshì, yìbiān fā xiāoxi.)

Barème indicatif : sujet correctement placé avant \zh{一边} (1 pt) ; répétition de \zh{一边} devant chaque action (2 pts) ; vocabulaire exact (\zh{看报纸}, \zh{听广播}, \zh{发消息}) (1 pt) ; ponctuation pleine chasse \zh{，} et \zh{。} (1 pt).

\textbf{B. Version} (5 points)

« Dimanche matin, je lis le journal tout en prenant mon petit-déjeuner. Mon petit frère fait ses devoirs tout en écoutant de la musique. »

Barème indicatif : sens global respecté (2 pts) ; structure « tout en + gérondif » correctement rendue (2 pts) ; vocabulaire précis (\zh{星期天} = dimanche, \zh{早饭} = petit-déjeuner, \zh{弟弟} = petit frère) (1 pt).

\textbf{Points de vigilance :} \zh{弟弟} signifie « petit frère » et non « frère » en général ; \zh{吃早饭} se traduit par « prendre le petit-déjeuner », pas « manger le riz du matin ».
\end{corrige}

\begin{corrige}[Exercice 11 — Type Bac : correction, ponctuation et expression écrite]
\textbf{A. Correction} (6 points, 2 points par phrase)
\begin{enumerate}
\item *\zh{他一边吃饭，看电视。} → \zh{他一边吃饭，一边看电视。} (Tā yìbiān chīfàn, yìbiān kàn diànshì.) Justification : \zh{一边} doit être répété devant la deuxième action ; le chinois ne permet pas de l'omettre comme le français omet « en ».
\item *\zh{一边他唱歌，一边跳舞。} → \zh{他一边唱歌，一边跳舞。} (Tā yìbiān chànggē, yìbiān tiàowǔ.) Justification : le sujet \zh{他} se place obligatoirement avant le premier \zh{一边}.
\item *\zh{他一边睡觉，一边跑步。} → phrase sémantiquement impossible : on ne peut pas dormir et courir en même temps. Correction possible : \zh{他一边跑步，一边听音乐。} (Tā yìbiān pǎobù, yìbiān tīng yīnyuè.) « Il écoute de la musique tout en courant. » Justification : les deux actions reliées par \zh{一边……一边……} doivent être compatibles.
\end{enumerate}

\textbf{B. Ponctuation} (6 points)

\zh{妈妈问我：“你做完作业了吗？”我说：“我一边做作业，一边听了《新闻联播》，我还买了报纸、杂志和书。”}

(Māma wèn wǒ : “Nǐ zuò wán zuòyè le ma ?” Wǒ shuō : “Wǒ yìbiān zuò zuòyè, yìbiān tīng le 《Xīnwén Liánbō》, wǒ hái mǎi le bàozhǐ, zázhì hé shū.”)

« Maman m'a demandé : “As-tu fini tes devoirs ?” J'ai répondu : “J'ai fait mes devoirs tout en écoutant le \emph{Journal télévisé}, et j'ai aussi acheté des journaux, des magazines et des livres.” »

Justifications : les paroles directes sont introduites par \zh{：} et placées entre \zh{“”} ; la question se termine par \zh{？} ; le titre de l'émission prend \zh{《》} ; l'énumération \zh{报纸、杂志} prend \zh{、} (rien devant \zh{和}) ; la phrase déclarative se termine par \zh{。} placé avant le guillemet fermant.

\textbf{C. Expression écrite — proposition de rédaction modèle} (8 points)

\zh{我家的晚上}

\zh{晚上，我们一家人都在家。爸爸一边看报纸，一边喝茶。妈妈一边做饭，一边听广播。我一边做作业，一边听音乐。我们都喜欢看《新闻联播》。桌子上有很多东西：报纸、杂志和书。}

(Wǎnshàng, wǒmen yì jiā rén dōu zài jiā. Bàba yìbiān kàn bàozhǐ, yìbiān hē chá. Māma yìbiān zuò fàn, yìbiān tīng guǎngbō. Wǒ yìbiān zuò zuòyè, yìbiān tīng yīnyuè. Wǒmen dōu xǐhuan kàn 《Xīnwén Liánbō》. Zhuōzi shàng yǒu hěn duō dōngxi : bàozhǐ, zázhì hé shū.)

« Le soir, toute notre famille est à la maison. Papa lit le journal tout en buvant du thé. Maman cuisine tout en écoutant la radio. Je fais mes devoirs tout en écoutant de la musique. Nous aimons tous regarder le \emph{Journal télévisé}. Sur la table, il y a beaucoup de choses : des journaux, des magazines et des livres. »

Barème indicatif : deux structures \zh{一边……一边……} correctes (3 pts) ; énumération avec \zh{、} (1 pt) ; titre d'émission entre \zh{《》} (1 pt) ; respect de la longueur 40–60 caractères (1 pt) ; exactitude des caractères et de la ponctuation (2 pts).

\textbf{Points de vigilance :} compter les caractères (la ponctuation ne compte pas) ; ne pas oublier le titre \zh{我家的晚上} ; vérifier que chaque \zh{一边} est bien suivi d'un groupe verbal complet.
\end{corrige}

\newpage

\coursec{Fiche bilan}

\begin{popfigure}
\begin{tikzpicture}[mindmap, grow cyclic, every node/.style={concept, font=\small}, text width=2.6cm, align=center,
root concept/.append style={concept color=popblue, font=\bfseries},
level 1/.append style={level distance=3.4cm, sibling angle=60},
level 1 concept/.append style={font=\small}]
\node[root concept, align=center]{Leçon 40\\ \zh{一边……一边……}\\ Ponctuation}
  child[concept color=poporange]{ node[align=center]{Schéma\\ \zh{一边} + V1\\ \zh{一边} + V2} }
  child[concept color=popgreen]{ node[align=center]{Sens\\ deux actions\\ simultanées\\ (même sujet)} }
  child[concept color=poppurple]{ node[align=center]{Variante\\ \zh{边……边……}\\ (verbes courts)} }
  child[concept color=poppink]{ node[align=center]{Pièges\\ pas deux sujets\\ différents ;\\ pas avec état\\ long (\zh{在})} }
  child[concept color=popgold]{ node[align=center]{Ponctuation\\ \zh{。，、？！：；} « 》} }
  child[concept color=popblue!60]{ node[align=center]{Méthode\\ repérer 2 actions\\ simultanées\\ puis traduire\\ « en + gérondif »} };
\end{tikzpicture}
\legende{Carte mentale de la leçon 40 : la simultanéité avec \zh{一边……一边……} et la ponctuation chinoise.}
\end{popfigure}

\begin{bilan}
\textbf{1. Lexique clé (à connaître par cœur)}
\begin{center}
\begin{tabularx}{\linewidth}{|X|X|X|X|}
\hline
\rowcolor{popblueL} Caractères & Pinyin & Nature & Traduction \\
\hline
\zh{一边……一边……} & yìbiān... yìbiān... & structure & en même temps, tout en... \\
\hline
\zh{边……边……} & biān... biān... & structure & variante courte (verbes monosyllabiques) \\
\hline
\zh{同时} & tóngshí & adverbe & en même temps \\
\hline
\zh{听} & tīng & verbe & écouter \\
\hline
\zh{唱} & chàng & verbe & chanter \\
\hline
\zh{广播} & guǎngbō & nom & radio, radiodiffusion \\
\hline
\zh{新闻} & xīnwén & nom & informations, nouvelles \\
\hline
\zh{打电话} & dǎ diànhuà & verbe + obj. & téléphoner \\
\hline
\end{tabularx}
\end{center}

\textbf{2. Règles de grammaire à connaître par cœur}
\begin{center}
\begin{tabularx}{\linewidth}{|X|X|X|}
\hline
\rowcolor{popblueL} Règle & Exemple (caractères + pinyin) & Traduction \\
\hline
Sujet + \zh{一边} + V1 + \zh{一边} + V2 : deux actions faites en même temps par le même sujet. & \zh{他一边吃饭，一边看电视。} Tā yìbiān chīfàn, yìbiān kàn diànshì. & Il mange en regardant la télévision. \\
\hline
Les deux verbes sont des \cle{actions actives} et duratives ; pas de verbes d'état (\zh{是}, \zh{喜欢} seuls). & \zh{我们一边走，一边聊天。} Wǒmen yìbiān zǒu, yìbiān liáotiān. & Nous marchons en bavardant. \\
\hline
Variante courte \zh{边……边……} avec verbes monosyllabiques. & \zh{她边唱边跳。} Tā biān chàng biān tiào. & Elle chante en dansant. \\
\hline
Deux sujets différents : utiliser \zh{……的时候} ou \zh{同时}, jamais \zh{一边……一边……}. & \zh{我学习的时候，弟弟在玩儿。} Wǒ xuéxí de shíhou, dìdi zài wánr. & Pendant que j'étudie, mon petit frère joue. \\
\hline
\zh{了} se place après chaque verbe si les deux actions sont accomplies (rare) ; le plus souvent sans \zh{了}. & \zh{妈妈一边做饭，一边听广播。} Māma yìbiān zuòfàn, yìbiān tīng guǎngbō. & Maman cuisine en écoutant la radio. \\
\hline
\end{tabularx}
\end{center}

\textbf{3. La ponctuation chinoise (essentiel)}
\begin{itemize}
  \item \zh{。} point final (un petit cercle, pas un point) ; \zh{，} virgule ; \zh{、} virgule d'énumération entre mots coordonnés : \zh{我喜欢唱歌、跳舞和画画。}
  \item \zh{？} interrogation (même avec \zh{吗} : \zh{你好吗？}) ; \zh{！} exclamation ; \zh{：} deux-points avant une citation ou une explication ; \zh{；} point-virgule entre deux propositions liées.
  \item Guillemets : « 》 pour les citations ; \zh{《》} pour les titres d'œuvres (livres, films, journaux) : \zh{《人民日报》}.
  \item Chaque signe occupe \cle{une case entière} dans l'écriture ; jamais de point anglais « . » ni de virgule « , » dans un texte chinois.
\end{itemize}

\textbf{4. Méthode express (traduction)}
\begin{enumerate}
  \item Repérer les deux actions simultanées et vérifier qu'elles ont le même sujet.
  \item Choisir \zh{一边……一边……} (ou la forme courte \zh{边……边……}).
  \item Traduire en français par « en + gérondif » ou « tout en + infinitif ».
  \item Relire : ordre Sujet + \zh{一边} + V1, \zh{一边} + V2 ; ponctuation pleine chasse.
\end{enumerate}
\end{bilan}

\begin{aretenir}[Checklist avant le Bac]
\begin{itemize}
  \item $\square$ Je connais le schéma : Sujet + \zh{一边} + V1 + \zh{一边} + V2.
  \item $\square$ Je sais que les deux actions doivent avoir le \cle{même sujet}.
  \item $\square$ Je sais employer la variante courte \zh{边……边……} avec des verbes monosyllabiques.
  \item $\square$ Je n'utilise pas \zh{一边……一边……} avec deux sujets différents (je prends \zh{……的时候}).
  \item $\square$ Je traduis correctement par « en + gérondif » en français.
  \item $\square$ Je place la virgule chinoise \zh{，} entre les deux propositions.
  \item $\square$ Je connais la différence entre \zh{，} (virgule) et \zh{、} (énumération).
  \item $\square$ J'écris le point chinois \zh{。} (cercle) et non le point français « . ».
  \item $\square$ Je sais que \zh{《》} sert pour les titres d'œuvres et « 》 pour les citations.
  \item $\square$ Je mets \zh{？} même dans une question avec \zh{吗}.
  \item $\square$ Je sais produire trois phrases personnelles avec \zh{一边……一边……} (radio, téléphone, devoirs).
  \item $\square$ Je relis mes textes : un signe de ponctuation = une case, ponctuation pleine chasse.
\end{itemize}
\end{aretenir}

\findecours

\end{document}
